% TODO: use \bar\psi for sample-average estimand (matching retargeted-mean paper)
% === FORMAT SELECTION ===
% For arXiv: uncomment \arxivtrue
% For Biometrika main: leave as is
% For Biometrika supplement: use supplement.tex instead
\newif\ifarxiv
\newif\ifsupp
\arxivtrue    % arXiv: article class, biblatex, includes appendix
%\arxivfalse   % Biometrika main: biometrika class, natbib, no appendix
%\supptrue    % Biometrika supplement (use via make supp)
\suppfalse

% Optional: blind author names
\newif\ifblind
\blindfalse  % set \blindtrue for anonymous submission

% === END FORMAT SELECTION ===

\ifarxiv
  \documentclass{article}
  \usepackage[margin=1.0in, top=1.0in]{geometry}
  \usepackage[natbib=true, style=apa]{biblatex}
  \setlength\bibitemsep{.7\baselineskip}
  \setlength{\bibhang}{10pt}
\else
  \ifsupp
    \documentclass[supplementary]{biometrika}
  \else
    \documentclass[lineno]{biometrika}
  \fi
  \usepackage{natbib}
  \bibliographystyle{biometrika}
  % biometrika.cls doesn't define these
  \newcommand{\subsubsection}[1]{\par\noindent\textit{#1}\enspace}
  \newcommand{\paragraph}[1]{\par\noindent\textit{#1}\enspace}
\fi

\synctex=1

\usepackage{amsmath}
\usepackage{graphicx}
\usepackage{enumerate}
\usepackage{url}

% Derive appendix setting from format choice
% arXiv and supplement both include appendix
\newif\ifappendix
\ifarxiv\appendixtrue\else\ifsupp\appendixtrue\else\appendixfalse\fi\fi

% Main body is included except in supplement mode
\newif\ifmainbody
\ifsupp\mainbodyfalse\else\mainbodytrue\fi

% For appendix propositions with primed numbers matching body propositions
\ifappendix
\newenvironment{propositionprime}[1]{%
  \par\medskip\noindent\textbf{Proposition #1$'$.}\itshape\space
}{%
  \par\medskip
}
\newenvironment{corollaryprime}[1]{%
  \par\medskip\noindent\textbf{Corollary #1$'$.}\itshape\space
}{%
  \par\medskip
}
\fi

%% Our stuff
\usepackage{amssymb}
\usepackage{mathtools}
\usepackage{epstopdf}
\usepackage[multiple]{footmisc}
\usepackage{placeins}
\DeclareGraphicsRule{.tif}{png}{.png}{`convert #1 `dirname #1`/`basename #1 .tif`.png}
\DeclarePairedDelimiter\norm{\lVert}{\rVert}
\DeclarePairedDelimiter\abs{\lvert}{\rvert}
\DeclarePairedDelimiter\set{\{}{\}}
\DeclarePairedDelimiter\cb{\{}{\}}
\DeclarePairedDelimiter\p{(}{)}
\DeclarePairedDelimiter\sqb{[}{]}
\DeclarePairedDelimiter\inner{\langle}{\rangle}
\DeclareMathOperator{\vectorspan}{span}
\DeclareMathOperator{\expo}{expo}
\usepackage{bbm}
\usepackage{amsfonts}
\newcommand{\indicator}[1]{\mathbbm{1}_{ {#1} }}
\usepackage{lscape}
\usepackage{rotating}
\usepackage{setspace}
\usepackage{threeparttable}
\usepackage{booktabs}

\usepackage{lmodern}
\fontfamily{lmtt}\selectfont
%\renewcommand*\familydefault{\sfdefault} % If the base font of the document is to be sans serif
\usepackage[T1]{fontenc}
\DeclareMathOperator*{\argmax}{argmax}
\DeclareMathOperator*{\argmin}{argmin}

% \bibpunct{(}{)}{;}{a}{}{,}
% \usepackage{hyperref}
\RequirePackage[colorlinks,allcolors=blue]{hyperref}
% Cross-references between main and supplement (Biometrika only)
% Build order: main first, then supplement, then main again (for forward refs)
\ifarxiv\else
  \usepackage{xr-hyper}
  \ifsupp
    \externaldocument{main}
  \else
    \externaldocument{supplement}
  \fi
\fi
\ifarxiv
  \usepackage{titling}
  \newcommand{\subtitle}[1]{%
  \posttitle{%
  \par\end{center}
  \begin{center}\large#1\end{center}
  \vskip0.5em}%
  }
  \usepackage{amsthm}
  \theoremstyle{definition}
  \newtheorem{theorem}{Theorem}[section]
  \newtheorem{lemma}[theorem]{Lemma}
  \newtheorem{proposition}[theorem]{Proposition}
  \newtheorem{corollary}[theorem]{Corollary}
  \newtheorem{problem}[theorem]{Problem}
  \newtheorem{algorithm}[theorem]{Algorithm}
  \newtheorem{assumption}{Assumption}
  \newtheorem{remark}{Remark}
  \newtheorem{example}{Example}
  \newtheorem{condition}{Condition}
  \newtheorem{definition}{Definition}
\fi
\def\ci{\perp\!\!\!\perp}
\usepackage{color}
%\usepackage{fontspec}
%\newfontfamily\DejaSans{DejaVu Sans}

\newcommand\indep{\protect\mathpalette{\protect\independenT}{\perp}}
\def\independenT#1#2{\mathrel{\rlap{$#1#2$}\mkern2mu{#1#2}}}
%%% MACROS
\newcommand{\R}{\ensuremath{\mathbb{R}}}
\newcommand{\B}{\ensuremath{\mathbb{B}}}
\newcommand{\bbL}{\ensuremath{\mathbb{L}}}
\newcommand{\bbC}{\ensuremath{\mathbb{C}}}
\newcommand{\N}{\ensuremath{\mathbb{N}}}
\newcommand{\Q}{\ensuremath{\mathbb{Q}}}
\newcommand{\Z}{\ensuremath{\mathbb{Z}}}
\newcommand{\I}{\ensuremath{\mathcal{I}}}
% \newcommand{\1}{\ensuremath{\mathbbm{1}}}
\newcommand{\calF}{\ensuremath{\mathcal{F}}}
\newcommand{\calC}{\ensuremath{\mathcal{C}}}
\newcommand{\calA}{\ensuremath{\mathcal{A}}}
\newcommand{\calB}{\ensuremath{\mathcal{B}}}
\newcommand{\calO}{\ensuremath{\mathcal{O}}}
\DeclareMathOperator{\E}{E}
\newcommand{\calE}{\ensuremath{\mathcal{E}}}
\newcommand{\calG}{\ensuremath{\mathcal{G}}}
\newcommand{\calR}{\ensuremath{\mathcal{R}}}
\newcommand{\calS}{\ensuremath{\mathcal{S}}}
\newcommand{\calJ}{\ensuremath{\mathcal{J}}}
\newcommand{\calT}{\ensuremath{\mathcal{T}}}
\newcommand{\calU}{\ensuremath{\mathcal{U}}}
\newcommand{\bbH}{\ensuremath{\mathbbm{H}}}
\newcommand{\bbP}{\ensuremath{\mathbbm{P}}}
\newcommand{\calN}{\ensuremath{\mathcal{N}}}
\newcommand{\calM}{\ensuremath{\mathcal{M}}}
\newcommand{\calX}{\ensuremath{\mathcal{X}}}
\newcommand{\calY}{\ensuremath{\mathcal{Y}}}
\newcommand{\calD}{\ensuremath{\mathcal{D}}}
\newcommand{\calL}{\ensuremath{\mathcal{L}}}
\newcommand{\calH}{\ensuremath{\mathcal{H}}}
\DeclareMathOperator{\Cov}{cov}
\DeclareMathOperator{\Var}{var}
\DeclareMathOperator{\MSE}{MSE}
\DeclareMathOperator{\logit}{logit}
\DeclareMathOperator{\trace}{tr}
\newcommand{\conv}{\text{conv}}
\newcommand{\sign}{\text{sign}}
\newcommand{\bipw}{} %\text{bIPW}
\newcommand{\mipw}{} %\text{mIPW}
\newcommand{\ipw}{\text{ipw}}
\DeclarePairedDelimiter\round{\lfloor}{\rceil}
\DeclarePairedDelimiter\floor{\lfloor}{\rfloor}
\DeclarePairedDelimiter\ceil{\lceil}{\rceil}

% essential notation
\newcommand{\estimand}{\psi}
\newcommand{\dispersion}{\chi}
\newcommand{\imbalancescale}{\zeta}

\newcommand{\dgamma}{\delta\gamma}
\newcommand{\dm}{\delta \mu}
\newcommand{\bias}{\imbalance}
\DeclareMathOperator*{\minimize}{minimize}
\DeclareMathOperator{\imbalance}{imbalance}
\DeclareMathOperator{\roundingerror}{rounding}
\DeclareMathOperator{\diameter}{diameter}
\newcommand{\diam}{\diameter}
\newcommand{\qfor}{\quad \text{for} \quad}
\newcommand{\qwhere}{\quad \text{where} \quad}
\newcommand{\qif}{\quad \text{if} \quad}
\newcommand{\qand}{\quad \text{and} \quad}
\DeclareMathOperator{\var}{Var}
\DeclareMathOperator{\disc}{disc}
\DeclareMathOperator{\hdisc}{hdisc}
\DeclareMathOperator{\lindisc}{lindisc}
\newcommand{\F}{\ensuremath{\mathcal{F}}}
\let\P\relax
\DeclareMathOperator{\Pn}{P_n}
\DeclareMathOperator{\P}{P}
\newcommand{\hgamma}{\ensuremath{\hat\gamma}}
%\newcommand{\hgammaipw}{\ensuremath{\hat\gamma^{ipw}}}
\newcommand{\hgammaipw}{\ensuremath{\hat\gamma}}
\newcommand{\gammaipw}{\ensuremath{\gamma^{\ipw}}}
\newcommand{\thetaipw}{\ensuremath{\theta^{ipw}}}
\newcommand{\htheta}{\hat\theta}
\newcommand{\model}{\mathcal{M}}

\newcommand{\aug}{\text{aug}}
\newcommand{\bal}{\text{bal}}

\newcommand{\obs}{\text{obs}}

\def\super{\textsuperscript}
\def\sub{\textsubscript}

\newcommand{\<}{\langle}
\renewcommand{\>}{\rangle}

\newcommand\pderiv[2]{\frac{\partial #1}{\partial #2}}
\newcommand\deriv[2]{\frac{d #1}{d #2}}

\def\b1{\boldsymbol{1}}
\newcommand{\ind}[1]{\b1_{\left\{#1\right\}}}

%% For editing
\newcommand\cmnt[2]{\;
{\textcolor{red}{[{\em #1 --- #2}] \;}
}}
% For a version without comments
%\renewcommand\cmnt[]{}

\newcommand\avi[1]{\cmnt{#1}{Avi}}
\newcommand\eli[1]{\cmnt{#1}{Eli}}
\newcommand\jose[1]{\cmnt{#1}{Jose}}
\newcommand\dah[1]{\cmnt{#1}{Skip}}
\newcommand\rmk[1]{\;\textcolor{red}{{\em #1}\;}}

%%%%%%%%%%%%%%%%%%%%%%%%%%%

\ifarxiv
  \addbibresource{review.bib}
\fi
\newcommand{\hepsilon}{\hat\epsilon}

\begin{document}

% === FRONT MATTER ===
\ifarxiv
  % --- arXiv format ---
  \ifblind
    \bigskip\bigskip\bigskip
    \begin{center}
      {\LARGE\bf The Balancing Act in Causal Inference}
    \end{center}
    \medskip
  \else
    \title{\bf The Balancing Act in Causal Inference\thanks{For comments and suggestions, we thank Tim Armstrong, Peng Ding, Vitor Hadad, Whitney Newey, and Stefan Wager. This work was partly supported through a grant from the Alfred P. Sloan Foundation (G-2020-13946) and an award from the Patient Centered Outcomes Research Initiative (PCORI, ME-2022C1-25648).} \bigskip}
    \author{David A. Hirshberg\thanks{Emory University}
    \and Eli Ben-Michael\thanks{Carnegie Mellon University}
    \and Avi Feller\thanks{UC Berkeley}
    \and Chen Lu\thanks{MIT}
    \hspace{.75cm} Jos\'e R. Zubizarreta\thanks{Harvard University}}
    \maketitle
  \fi
\else
  % --- Biometrika format ---
  \title{The Balancing Act in Causal Inference}
  \author{David A. Hirshberg}
  \affil{Emory University}
  \author{Eli Ben-Michael}
  \affil{Carnegie Mellon University}
  \author{Avi Feller}
  \affil{UC Berkeley}
  \author{Chen Lu}
  \affil{MIT}
  \author{Jos\'e R. Zubizarreta}
  \affil{Harvard University}
  \maketitle
\fi

\ifmainbody
\bigskip

\begin{abstract}
\ifarxiv\singlespacing\fi
The idea of covariate balance is at the core of causal inference.
In randomized experiments, it is an expected outcome of randomization, whereas in observational studies balance is a goal we strive to approximate through adjustment methods such as matching and weighting.
Inverse propensity weights play a central role because they are the unique set of weights that balance the covariate distributions of different treatment groups.
We discuss two broad approaches to estimating these weights:
the more traditional one,
which fits a propensity score model and then uses the reciprocal of the estimated propensity score to construct weights,
and the balancing approach, which estimates the inverse propensity weights essentially by the method of moments, finding weights that achieve balance in the sample.
We connect ideas from the causal inference, sample surveys, and semiparametric estimation literatures, with particular attention to the role of balance as a sufficient condition for robust inference.
We bridge practical approaches to matching in causal inference with the theoretical framework of semiparametric theory.
We focus on the inverse propensity weighting and augmented inverse propensity weighting estimators for the average treatment effect and consider generalizations for a broader class of problems including policy evaluation and the estimation of individualized treatment effects.
\end{abstract}

\ifarxiv
  \noindent\emph{Keywords: Causal Inference; Matching Methods; Propensity Score; Observational Studies; Weighting Methods}
  \vfill
\else
  \begin{keywords}
  Causal inference; Matching methods; Propensity score; Observational studies; Weighting methods
  \end{keywords}
\fi

\newpage

\section{Introduction}
\label{sec:intro}

%The idea of covariate balance is at the core of causal inference.
%Intuitively, we want to compare ``like with like'' when estimating the effect of a treatment: if two groups are similar in every respect except in assignment to treatment, then we can attribute differences in outcomes to the treatment itself, as opposed to other confounding factors \citep{cochran1965planning}.
%% When estimating the average effect of a treatment, the intuition is to compare like-with-like, such that if the treatment groups are similar in every respect except in the assignment to treatment, then differences in outcomes can be safely attributed to differences in the treatment as opposed to other confounding factors \citep{cochran1965planning}.
%While this intuition is widespread, it is typically used informally and without connection to other concepts used to motivate and study treatment effect estimators. In this paper, we explore
%% explain
%the role covariate balance plays in statistical inference about treatment effects, focusing on these connections.
The idea of covariate balance is at the core of causal inference.
Intuitively, when estimating the effect of a treatment, we want to compare ``like with like'': if two groups are similar in every respect except in assignment to treatment, then we can attribute differences in outcomes to the treatment itself, as opposed to other confounding factors \citep{cochran1965planning}.

This intuition is widespread. However, any two groups will differ in some respect, so when we talk about a like with like comparison,
we rely on a consensus about which
differences matter and which can be ignored.
For example, we might focus on imbalances in the means of covariates, or in their marginal distributions, or in summaries of their joint distribution. In this paper,
we explore the role covariate balance plays in statistical inference about treatment effects; its relationship to common notions of balance; and connections to the causal inference, sample surveys, and semiparametrics literatures.

\emph{Inverse propensity weights} play a central role because they are the unique set of weights that balance the covariate distributions of different treatment groups.
However, outside randomized experiments, these weights are unknown and must be estimated.

The most common approach to estimating inverse propensity weights is to estimate the conditional probability of treatment, or propensity score, and then take its reciprocal.
With some estimation methods, this approach can be unstable, especially when this probability is close to zero, and can lead to poor balance in the sample at hand when the propensity score is not estimated accurately.
Researchers often address these limitations with post-hoc fixes, e.g., truncating the extreme weights that arise from propensity score estimates that are too near zero or checking balance and re-estimating the propensity score if it is poor.

An alternative approach is to find weights that optimize balance in the data at hand.
Because balance is a defining property of the
inverse propensity weights, this is essentially a method of moments estimate of the inverse propensity weights.
Researchers typically estimate such weights by solving a constrained optimization problem, e.g., by finding weights that restrict the maximum imbalance in covariate means between groups.
Balance checks are therefore built in: we choose weights to pass the checks, rather than performing them post-hoc.
Examples of proposals in line with this approach include
\citet{athey2018approximate},
\citet{breidt2017model},
\citet{chan2016globally},
\citet{deville1992calibration},
\citet{graham2012inverse},
\citet{hainmueller2012balancing},
\citet{hirshberg2021augmented},
\citet{imai2014covariate},
\citet{kallus2020generalized},
\citet{li2018balancing},
\citet{tan2020regularized},
\citet{wong2018kernel},
\citet{zhao2019covariate}, and
\citet{zubizarreta2015stable}.

We will organize our discussion around a specific form of balance that is necessary and sufficient for straightforward inference based on the widely-used inverse probability weighting (IPW) and augmented inverse probability weighting (AIPW) estimators. For IPW, we need balance in the mean outcome conditional on treatment and covariates, and for AIPW, balance in the errors we make when estimating that conditional mean.
This happens when we estimate the inverse propensity weights and this conditional mean function sufficiently well, an idea that gives rise to the concept of double-robustness, and also when we balance every function in a set that contains
the conditional mean, an idea that gives rise to balance checks and weights optimized to achieve balance in this sense.
In particular, many widely-used balance checks for weights are based on the \emph{maximal imbalance} that would occur for
a conditional mean in a set chosen to encode structural assumptions about it. %, and weights as controlling this maximal imbalance subject to other priorities or constraints.
This gives us a practical framework to link intuition and assumptions about the data to appropriate balance criteria to check or optimize.

%In Sections \ref{sec:framework}-\ref{sec:asymptotic} below,
%we frame the problem of estimating causal effects using inverse propensity weights, discuss %approaches to estimating these weights, and
%discuss inference for causal effects.
%In Sections \ref{sec:balance_models}-\ref{sec:misfit_toys},
%we discuss the modeling choices involved in the estimation of
%inverse propensity weights, with a focus on approaches arising from the balancing %interpretation. In Section~\ref{sec:other-estimands}, we generalize our discussion from
%the specific problem of estimating average treatment effects to a broader class of
%problems including policy evaluation and the estimation of individualized treatment effects.
%We conclude, in Section~\ref{sec:dual}, by discussing an equivalence between balance criteria
%and loss functions for estimating the propensity score,
%including equivalences or near-equivalences between specific approaches
%described as focused on balance and those described as focused on asymptotic efficiency.

%%%%%%%%%%%%%%%%%%%%%%%%%%%
%%%%%%%%%%%%%%%%%%%%%%%%%%%
%%%%%%%%%%%%%%%%%%%%%%%%%%%
% \vspace{-.5cm}
\section{Framework}
\label{sec:framework}

%%%%%%%%%%%%%%%%%%%%%%%%%%%
%%%%%%%%%%%%%%%%%%%%%%%%%%%
\vspace{-.3cm}
\subsection{Setup}

We posit the existence of independent and identically distributed tuples $(X_i, W_i, Y_i(0), Y_i(1))$ randomly drawn from a population, where $X_i \in \R^p$ are observed covariates, $W_i \in \{0,1\}$ is a binary treatment indicator, and $Y_i(0), Y_i(1) \in \R$ are the potential outcomes under control and treatment, respectively \citep{neyman1923,rubin1974estimating}.
Under the Stable Unit Treatment Value Assumption (SUTVA; \citeauthor{rubin1980}, 1980), we can write the observed outcome $Y_i$ as $Y_i = (1-W_i)Y_i(0) + W_i Y_i(1)$.
Since the tuples are randomly drawn from a population and we are interested in population-level estimands, we generally drop the subscript $i$.

We are primarily interested in the {\it Average Treatment Effect} (ATE), defined as
\begin{equation}
  \label{eq:pate}
  \tau := \E[Y(1) - Y(0)].
\end{equation}
To ease exposition, we separate it into two components, $\psi_1 := \E[Y(1)]$ and $\psi_0 := \E[Y(0)]$, and focus on estimating the treatment-specific mean $\psi_1$; we can estimate $\psi_0$ analogously.


We work with the assumption of strong ignorability \citep{Rosenbaum1983}.
\begin{assumption}[Ignorability]
  \label{a:ignore}
  $W \indep (Y(0), Y(1))\mid X$
\end{assumption}

\begin{assumption}[Overlap]
  \label{a:overlap}
  $e(x) > 0$ where $e(x) := \E[W \mid X = x]$ is the \emph{propensity score}
\end{assumption}

\noindent With  Assumptions~\ref{a:ignore}-\ref{a:overlap}, the treatment-specific mean is the average conditional mean outcome.%\footnote{When estimating the PATE, we use a two-sided version of this assumption, as if $1-e(x) > 0$, $\psi_0 = \E[\mu_0(X)]$.}
\begin{equation}
 \label{eq_identification}
\psi_1 = \E[\mu_1(X)] \quad \text{ where } \quad \mu_w(x) := \E[Y \mid W=w, X = x].
 \end{equation}
The main challenge when estimating $\psi_1$ is that, outside of randomized experiments,
the mix of units that receive one treatment tends to differ from the mix that receive another.
% Looking at the outcomes of those who receive the treatment, we see a combination of treatment and selection effects.
Thus, the average outcome among those who receive the treatment reflects selection into treatment as well as the treatment's effect.
Under Assumption \ref{a:ignore}, the observed covariates explain the salient differences between these groups, % (Assumption \ref{a:ignore}),
so eliminating selection effects reduces to a covariate adjustment problem.
% Specifically, we must adjust for systematic differences between the covariates of those who receive the treatment and those in the entire study.




%%%%%%%%%%%%%%%%%%%%%%%%%%%
%%%%%%%%%%%%%%%%%%%%%%%%%%%
\vspace{-.3cm}
\subsection{Inverse propensity weighting}
\label{sec:true_ipw}

A celebrated result in causal inference is that the inverse propensity weighted average of the outcomes is an unbiased estimator of $\psi_1$ under Assumptions \ref{a:ignore} and \ref{a:overlap} \citep{rosenbaum1987model}:
\begin{equation}
  \label{eq_mipw-identity}
  \psi_1 = \E\left[\mu_1(X_i)\right] =  \E\left[\frac{W_i}{e(X_i)} \mu_{W_i}(X_i)\right] = \E\left[\frac{W_i}{e(X_i)} Y_i \right].
\end{equation}

\noindent In fact, this result holds more generally: the inverse propensity weights are the unique weighting function that, for \emph{every} function $f(x)$,
equates weighted averages of $f$ over the treated units to unweighted averages over the study population.
That is, $\gammaipw(x)=1/e(x)$ is the only function  satisfying
\begin{equation}
    \label{eq_mipw-balance}
    \E\left[W_i \gammaipw(X_i) f(X_i)\right] = \E\left[f(X_i)\right]  \ \text{ for all bounded }\ f(x). \end{equation}
Because the solution is unique, Equation~\ref{eq_mipw-balance} also \emph{defines} the inverse propensity weights $\gammaipw$.
We call Equation~\ref{eq_mipw-balance} the \emph{population balance property}: inverse propensity weights adjust for
 differences (in the averages of all functions $f$ of the covariates) between the subpopulation receiving treatment and the overall population.
%  --- differences in the averages of all functions $f$ of the covariates.
Setting $f=\mu_1$ recovers the unbiasedness result above \eqref{eq_mipw-identity}. The law of large numbers implies a corresponding \emph{sample balance property}:
\begin{equation}
    \label{eq_mipw-sample-balance}
    \frac{1}{n}\sum_{i=1}^n W_i \gammaipw(X_i) f(X_i) \approx \frac{1}{n}\sum_{i=1}^n f(X_i) \   \text{ for any bounded }\ f(x).
\end{equation}
 When using an estimate of the inverse propensity weights $\gammaipw$, it is common to check that sample balance indeed holds for some set of functions. One common approach, in which we check that that each covariate has a similar average in the weighted treatment group and in the whole sample, looks at balance in the set of coordinate functions $f(x)=x_{j}$. Another, in which we compute the marginal histograms of each covariate for the two groups and compare them using the Kolmogorov-Smirnov statistic, looks at balance in a set of additive functions; see Appendix~\ref{sec:variation-models} for further discussion.
 Balance checks like this have a long history in observational causal inference \citep{cochran1965planning}.
 %Skipitty: COMMENT ON INSUFFICIENCY OF THESE PARAMETRIC CHECKS; BALANCING WEIGHTS HAVE MOVED ON, BUT CHECKS HAVE NOT. PERHAPS DUE TO VISUALIZATION DIFFICULTIES?
% PROPOSE SOMETHING AND LOOK AT IT IN SIMS?

%%%%%%%%%%%%%%%%%%%%%%%%%%%
%%%%%%%%%%%%%%%%%%%%%%%%%%%
% \subsection{Weighting for causal effects}
% \label{sec:ipw}
\vspace{-.3cm}
\subsection{Imbalance and estimation error}

Sample balance is not just a signal we use to appraise the accuracy with which we have estimated the inverse propensity weights. It is a key determinant of the error of estimators that use them. As starting point, we consider the error of a general weighting estimator,
\begin{equation}
\label{eq_general_ipw}
\hat{\psi}_1 = \frac{1}{n}\sum_{i=1}^n W_i \hat{\gamma}(X_i) Y_i.
\end{equation}
Throughout, we will consider only  design-based weights $\hat{\gamma}(X_i)$, i.e., weights that depend on the study's \emph{design} --- the treatment $W_i$ and covariates $X_i$ --- but not on the outcomes $Y_i$.\footnotemark\
The error of this estimator is, in terms of $\varepsilon_i :=  Y_i - \mu_{W_i}(X_i)$,
\begin{equation}
    \label{eq:weights_error}
    \hat{\psi}_1 - \psi_1 = \underbrace{\frac{1}{n}\sum_{i=1}^nW_i\hat{\gamma}(X_i)  \mu_1(X_i) - \frac{1}{n}\sum_{i=1}^n \mu_1(X_i)}_{\text{imbalance in $\mu_1$}} + \underbrace{\frac{1}{n}\sum_{i=1}^n W_i\hat{\gamma}(X_i) \varepsilon_i}_{\text{noise}} + \underbrace{\frac{1}{n}\sum_{i=1}^{n} \mu_1(X_i) - \psi_1}_{\text{sampling variation}},
\end{equation}
 The first term in Equation~\ref{eq:weights_error} is the imbalance in the conditional mean $\mu_1(X_i)$. The second and third terms are a weighted average of noise and irreducible sampling variation respectively.
Since these two terms both have mean zero conditional on the design $\mathcal{D}=\{ (W_i, X_i) : i \le n \}$, the estimator's design-conditional bias $\E[\hat\psi_1 - \psi_1 \mid \mathcal{D} ]$ \emph{is} the imbalance in $\mu_1$.

\footnotetext{
%\emph{Design-based} weights are weights that depend on the study design --- the treatment status and covariates of all units,
%$\mathcal{D} := \{(W_1,X_1) \ldots (W_n, X_n)\}$ --- but do not depend on the outcomes. In particular, they are weights satisfying $\hat{\gamma} \indep Y_1, \ldots, Y_n \mid \mathcal{D}$.
Furthermore, as implied by our notation, we will consider only weights that are functions of the covariate $X_i$.
This is almost universal among design-based weights, there being no design-based reason to distinguish units with identical covariate and treatment status by assigning them different weights, but is not always indicated notationally the way we do here: some write $\hgamma_i$ instead of $\hgamma(X_i)$ for the weight of the $i$th unit.}

What makes estimation challenging is that the conditional mean function $\mu_1$ is unknown. One natural approach is to think of $\mu_1$ as belonging to some set of functions $\model$, e.g., the set of linear functions, which serves as a \emph{model}. This is not a model for the propensity score or a parametric likelihood---it is a set of functions that we believe contains the true conditional mean $\mu_1$, and whose geometry determines how well our weights can control bias.
When $\mu_1$ belongs to this set, the absolute design-conditional bias
is bounded by the largest value it could take on for any $\mu_1 \in \model$, which we call the \emph{maximal imbalance}, $\imbalance_\model(\hgamma)$:
\begin{equation}
%   \label{eq:balance-bias-control-inside}
    \label{eq:bias}
    \abs*{\text{design-conditional bias}} \leq \imbalance_{\calM}(\hgamma) := \max_{\mu_1 \in \calM} \abs*{\frac{1}{n}\sum_{i=1}^n \mu_1(X_i) - \frac{1}{n}\sum_{i=1}^n W_i \hgamma(X_i) \mu_1(X_i) }.
\end{equation}
Controlling maximal imbalance is effective at controlling bias if we choose the model well.

% Balancing approaches focus on this, and if they use a large enough model class, also estimate the true inverse propensity weights because they are the unique weights that balance all functions of the covariates \eqref{eq_mipw-balance}.
%\avi{Probably OK to keep; but feels a bit redundant with start of sec 3 below}
%\avi{I tried to clean up this sentence; OK?}
%Overall, balancing approaches to IPW focus on controlling imbalance. And since the true inverse propensity weights are the unique weights that balance all functions of the covariates \eqref{eq_mipw-balance}, the balancing approach also estimates the true IPW (so long as they balance a large enough model class).
% On the other hand,
%Modeling approaches, by contrast, tend to focus more on estimating the propensity score accurately, the intuition being that if we do so well enough, the approach will (approximately) inherit the balance achieved by the true inverse propensity weights  \eqref{eq_mipw-sample-balance}.
%The difference between the two approaches is in what they prioritize when working with limited data; that is, what optimization problem they choose to solve for the data at hand.
%For some, that is balance; for others, it is accurate estimation of the treatment probability. These priorities are reflected in differences in the behavior of the resulting treatment effect estimators.

%%%%%%%%%%%%%%%%%%%%%%%%%%%
\vspace{-.3cm}
\subsection{Augmented Inverse Propensity Weighting}
\label{sec:aipw}

A prominent generalization of the IPW estimator augments it with an adjustment involving an estimate of the conditional mean function $\mu_1$ fit via regression on the treated sample (Robins et al., \citeyear{robins1994estimation}). %athey2018approximate, hirshberg2021augmented
This augmented IPW (AIPW) estimator uses the regression estimator $\hat \mu_1$ to impute and adjust for imbalance in $\mu_1$ that is left over after weighting, or equivalently uses weighting to adjust for imbalance in the regression residuals left over after imputation.
\begin{align}
\hat{\psi}^{\text{aug}}_1
:=& \underbrace{\frac{1}{n}\sum_{i=1}^n W_i\hat{\gamma}(X_i) Y_i}_{\text{weighting estimator}} + \underbrace{\frac{1}{n}\sum_{i=1}^n \left\{\hat \mu_1(X_i) - W_i\hat{\gamma}(X_i)\hat \mu_1(X_i)\right\}}_{\text{bias correction via imputation}} \label{eq:aipw2} \\
&= \underbrace{\frac{1}{n}\sum_{i=1}^n \hat \mu_1(X_i)}_{\text{imputation estimator}} + \underbrace{\frac{1}{n}\sum_{i=1}^n W_i \hat \gamma(X_i) \left\{Y_i - \hat \mu_1(X_i)\right\}}_{\text{bias correction via weighting}} \label{eq:aipw2-imputation}
\end{align}

Though it is traditional to use a modeling approach to estimate the weights used in the AIPW estimator,
many recent proposals combine this bias correction with control of maximal imbalance. Examples include \emph{approximate residual balancing} \citep{athey2018approximate}, \emph{augmented minimax linear estimation} \citep{hirshberg2021augmented}, and \emph{automatic double machine learning} \citep*{chernozhukov2018automatic}.

The augmented estimator has an error decomposition analogous to that of a pure weighting estimator. In it, imbalance in the \emph{regression error function}
$\delta \mu_1 = \mu_1-\hat \mu_1$ replaces imbalance in $\mu_1$ itself.
\begin{equation}
    \label{eq:weights_error_aug}
    \hat{\psi}^{\aug}_1 - \psi_1 = \underbrace{\frac{1}{n}\sum_{i=1}^nW_i\hat{\gamma}(X_i) \delta \mu_1(X_i) - \frac{1}{n}\sum_{i=1}^n \delta \mu_1(X_i)}_{\text{imbalance in $\dm$}} + \underbrace{\frac{1}{n}\sum_{i=1}^n W_i\hat{\gamma}(X_i) \varepsilon_i}_{\text{noise}} + \underbrace{\frac{1}{n}\sum_{i=1}^{n} \mu_1(X_i) - \psi_1}_{\text{sampling variation}}.
\end{equation}
Intuitively, imbalance with augmentation should be smaller than without because the regression error function $\delta m$ should be smaller than the conditional mean function $\mu_1$ itself.
This generalizes our previous error decomposition for $\hat\psi_1$: the IPW estimator is AIPW with $\hat \mu_1=0$, in which case $\delta m=\mu_1$ and the error decomposition reduces to Equation~\ref{eq:weights_error}.
Thus, maximal imbalance can play a similar role in our thinking;
%\avi{I found this next phrase hard to follow?}
however, in this case we will want our model $\model$ to contain the regression error $\dm$, not $\mu_1$ itself.
For brevity, our discussion will be in terms of the more general augmented estimator and the regression error $\dm$,
noting when we require that the regression estimator $\hat \mu_1$ is or is not zero.
% As we discuss in Section TK, augmented estimators often rely on extrapolation \citep{Ben-Michael2019}, unlike balancing weights, which are typically sample bounded. Thus, augmented estimators can be more sensitive to the choice of model.
% \dah{The sample boundedness claim isn't particularly universal --- lots of balancing estimators don't use nonnegative weights summing to one --- and I don't know what this claim about extrapolation and model sensitivity means. I think it's best to leave forward refs like this out and raise the issue where there's space for a clearer discussion.}

\begin{remark}
\label{remark:tmle}
One popular alternative to AIPW is the \emph{targeted maximum likelihood estimator} (TMLE), an imputation estimator  that uses a specialized regression estimator $\hat \mu_1$ so that bias correction is unnecessary. To make this happen, $\hat \mu_1$ is chosen so that the bias correction term in Equation~\ref{eq:aipw2-imputation} is small or even exactly zero. We can do this, for example, using inverse-probability weighted least squares regression in a regression model that includes an unconstrained intercept. The resulting least squares predictor $\hat \mu_1$ minimizes weighted squared error among all functions in the model and therefore among the subset $\{ \hat \mu_1(\cdot) + t : t \in \R \}$ of \emph{translations of $\hat \mu_1$}, and therefore satisfies a zero-derivative condition
$$
0 = \frac{d}{dt}\Big|_{t=0} \ \frac{1}{n}\sum_{i=1}^n W_i \hat\gamma(X_i) \{ Y_i - \hat \mu_1(X_i) - t \}^2 = -2 \times \frac{1}{n}\sum_{i=1}^n W_i \hat\gamma(X_i) \{ Y_i - \hat \mu_1(X_i) \}.
$$
that implies our bias correction term in \eqref{eq:aipw2-imputation} is exactly zero. For the purpose of our discussion below, this \emph{is} an AIPW estimator. The bias correction is there; it just happens to be zero.
\end{remark}




%%%%%%%%%%%%%%%%%%%%%%%%%%%
%%%%%%%%%%%%%%%%%%%%%%%%%%%
%%%%%%%%%%%%%%%%%%%%%%%%%%%
\section{Estimating inverse propensity score weights}
\label{sec:ipw}

In observational studies, the inverse propensity weights $\gammaipw$ are unknown and must be estimated. The two equivalent definitions of
$\gammaipw$, via the explicit form $\gammaipw(x)=1/e(x)$ and via the population balance property \eqref{eq_mipw-balance},
suggest two broad types of estimators for inverse propensity weights.
Traditional modeling approaches estimate $e(x)$ by regressing treatment status on covariates.
Balancing estimators instead directly estimate weights $\gamma(x)$ that satisfy the sample balance property \eqref{eq_mipw-sample-balance}.
We turn to these now.

%%%%%%%%%%%%%%%%%%%%%%%%%%%
%%%%%%%%%%%%%%%%%%%%%%%%%%%
\vspace{-.3cm}
\subsection{Modeling approach to weighting}
\label{sec:mipw}
By far, the most common approach to weighting is to take an estimate of the propensity score, $\hat{e}(x)$, and use its reciprocal to compute inverse propensity weights, $\hgammaipw(X_i)=1/\hat{e}(X_i)$.
Researchers typically estimate $e(x)$ via maximum likelihood, e.g., with a logistic regression of $W$ on $X$, though other learning approaches are used as well.
The sample balance property \eqref{eq_mipw-sample-balance} implies that if the estimated propensity score $\hat{e}(x)$ converges to the true propensity score $e(x)$,
$\hat{\psi}_1$ converges to $\psi_1$. This approach suffers from two major drawbacks.

The first is instability. Observations with small estimated probabilities of treatment tend to dominate the analysis. To some extent, this is inevitable, as we need to give rarely-treated observations large weights to ensure representativeness of the weighted sample. However, small errors in estimating treatment probability can lead to dramatic differences in the weights. If two units are assigned to treatment with probability $.05$, but we estimate that probability to be  $.01$ for one unit and $.09$ for the other, the first unit gets nine times the weight that the second does, though our probability estimate is off by only $.04$ in both cases.

The second is that the estimated weights $1 / \hat{e}(X_i)$ may not satisfy the sample balance property in Equation~\ref{eq_mipw-sample-balance}
with much (or any) accuracy, which can lead to bias. This can occur when the estimated weights are inaccurate, e.g., due to misspecification or a small sample size. In practice, it is often difficult to estimate conditional treatment probabilities, and the instability brought on by inversion compounds this problem.% small errors in estimating the propensity score can result in large errors in estimating its inverse, especially when the propensity score is small.

To deal with the first drawback, investigators sometimes \emph{trim} the inverse propensity weights,
taking $\hgammaipw(X_i)=\min\{ T, \ 1/\hat e(X_i) \}$ for some threshold $T$ \citep{crump2009dealing}. To deal with the second, investigators resort to balance checks like those discussed in Section~\ref{sec:true_ipw}. If imbalance is large, they often fit a new propensity score model, repeating this process until they find one that achieves good balance. This is a good
accompaniment to trimming, which can introduce imbalances that are not present with the untrimmed weights $\hgammaipw(X_i)=1/\hat e(X_i)$.


%%%%%%%%%%%%%%%%%%%%%%%%%%%
%%%%%%%%%%%%%%%%%%%%%%%%%%%
\vspace{-.3cm}
\subsection{Balancing approach to weighting}
\label{sec:bal_weights_intro}

An alternative approach is to find weights $\hat{\gamma}$ by solving for weights that satisfy the sample balance property \eqref{eq_mipw-sample-balance}, sometimes known as \emph{balancing weights}.
While different in its implementation, this is also an inverse propensity weighting estimator:
we can interpret balancing weights as a \emph{method of moments} estimate of the inverse propensity weights.
Specifically, as we discuss in Section \ref{sec:true_ipw}, the inverse propensity score weights $\gammaipw$ are the unique solution to the population balance moment conditions \eqref{eq_mipw-balance}.
% Since the inverse propensity score $\gammaipw$ is the unique solution to the population balance moment conditions \eqref{eq_mipw-balance},
The balancing approach is to find
% We estimate it when we choose
weights $\hgamma$ that satisfy
%enough
the corresponding sample balance moment conditions \eqref{eq_mipw-sample-balance},
i.e., weights for which $\imbalance_{\model}(\hat\gamma)$ is small for some model $\model$.
% In other words, the modeling approach to IPW assesses covariate imbalance as a post-hoc check on the estimated weights; the balancing approach directly uses covariate imbalance to estimate the weights themselves.
%
% \dah{These changes make the implication that solving sample analogs of the moment conditions defining the inverse propensity weights a little hazy. Was that intended? It also loses the idea that you need enough moment conditions --- a large enough model --- to identify it.}
% \avi{just a little confusing to talk about  $\imbalance_{\model}(\hat\gamma) \approx 0$ in this paragraph and then
% $\imbalance_{\model}(\hat\gamma) \leq \lambda$ in the next; trying to smooth that a bit. But otherwise, was just trying to smooth text.}


Typically, this is done by solving an optimization problem to choose a weighting function $\hgamma(\cdot)$ to trade off imbalance  and some measure of the complexity of the weights,
\begin{equation}
\label{eq_balancing-weights}
 \hgamma = \mathop{\argmin}_{\gamma}~~ \Bigg\{ \imbalance_{\model}^2(\gamma) \;\;+\;\; \frac{\sigma^2}{n^2} \sum_{W_i=1}\gamma(X_i)^2  \Bigg\},
\end{equation}
\noindent where $\sigma$ is a tuning parameter,
 or more generally, for increasing convex functions $\imbalancescale$ and $\dispersion$,
\begin{equation}
\label{eq_balancing-weights-general}
 \hgamma = \mathop{\argmin}_{\gamma}~~ \Bigg\{ \imbalancescale\{\imbalance_{\model}(\gamma)\} \;\;+\;\; \frac{1}{n^2} \sum_{W_i=1}\dispersion\{\gamma(X_i)\}  \Bigg\}.
\end{equation}
\noindent In the familiar special case that the weights control the \emph{maximal imbalance} for the linear model with $\ell_1$-bounded coefficients, $\model = \{\beta \cdot x  : \|\beta\|_1 \leq 1\}$,
% We discuss these choices in the next sections.
% An important special case is where the weights control the \emph{maximal imbalance} for the linear model $\model = \{x'\beta: \|\beta\|_1 \leq 1\}$ as in Equation\ref{eq:bias};
% that is, $\imbalance_{\model}(\gamma)$ is the largest imbalance in the individual elements of the covariate vector $x$.
the corresponding imbalance measure is the largest  imbalance in the individual elements of the covariate vector $x$:
\begin{equation} \label{eq:sbw}
\imbalance_{\model}(\gamma) = \max_{j=1,\ldots,p} \left| \frac{1}{n} \sum_{i=1}^n X_{ij} - \frac{1}{n}\sum_{i=1}^n W_i \gamma(X_i) X_{ij} \right| \text{ for }
\model = \{\beta \cdot x : \|\beta\|_1 \leq 1\}.
\end{equation}
%When $\imbalance_{\model}(\gamma) = 0$, the resulting weights achieve \emph{exact balance}; otherwise, the weights only achieve \emph{approximate balance}.

The special case \eqref{eq_balancing-weights} has a natural \emph{minimax} interpretation: the weights minimize the worst-case
design-conditional mean squared error for $\mu_1 \in \model$
and $\Var[Y_i \mid X_i, W_i] \le \sigma^2$.
Furthermore, when the model $\model$ is an absolutely convex set, they have an equivalent dual characterization
as a least-squares estimate of the inverse propensity score. That is, the weights are the function of the covariates $\hat{\gamma}(X_i)$ that solves a penalized least squares problem,
\begin{equation}
  \label{eq:dual_l2_penalized_regression}
  \min_{\gamma} ~~~\underbrace{\frac{1}{n}\sum_{i=1}^n W_i \left\{\gammaipw(X_i) - \gamma(X_i)\right\}^2}_{\text{mean squared error}} \;-\; \underbrace{\left\{\frac{1}{n}\sum_{i=1}^n\gamma(X_i) - \frac{1}{n}\sum_{i=1}^n W_i \gammaipw(X_i)\gamma(X_i) \right\}}_{\text{mean zero ``noise''}} \;+\; \frac{\sigma^2}{n} \|\gamma\|_\calM^2.
\end{equation}
The penalty is proportional to the square of the
\emph{gauge} of the model, $\|\gamma\|_\model = \inf \{\alpha > 0 : \gamma \in \alpha \model\}$ for $\alpha\model = \{ \alpha m : m \in \model \}$, which is the factor $\alpha$ by which we must \emph{inflate} the model to get a set that contains $\gamma$.\footnote{The gauge of an absolutely convex set of functions of $x$ is a norm or seminorm on functions of $x$. For example, for the linear model with $\ell_1$-bounded coefficients $\model = \{\beta \cdot x : \|\beta\|_1 \leq 1\}$ , the gauge is the $\ell_1$-norm of the coefficients
of $\gamma(x) = \beta \cdot x$:  $\|\gamma\|_\model = \|\beta\|_1$. It is infinite if
$\gamma(x)$ is not linear in $x$, i.e., if $\gamma(x)$ is not in the space spanned by the set of functions in the model. If our model is defined as the unit ball of some norm on functions, i.e. if $\model = \{ m : \norm{m} \le 1 \}$, then its gauge is just that norm $\norm{\cdot}$.}
The balancing property of the inverse propensity weights implies the middle term has mean zero, so up to this noise-like term, the minimax weights estimate $\gammaipw$ by minimizing penalized mean squared error.
%Skipitty: COMMENT ON LACK OF NEED TO TRIM AND CHECK BALANCE.
% so up to this noise-like term, we are minimizing the penalized mean squared error of $\gamma$ as an estimator of $\gammaipw$.

%Skipitty: MAKE THIS FIT HERE.
Many existing approaches can be written as examples of Equation~\ref{eq_balancing-weights-general} with various
choices of $\imbalancescale$ and $\dispersion$, including
\emph{entropy balancing} \citep{hainmueller2012balancing}; the \emph{stable balancing weights} approach \citep{zubizarreta2015stable};
the \emph{lasso minimum distance estimator} of the inverse propensity score \citep{chernozhukov2018automatic};
and the \emph{regularized calibrated propensity score estimator} \citep{tan2020regularized} (see \cite{wang2020minimal}).
In fact, even imputing $\psi_1$ by averaging the predictions of least squares or ridge regressions  of $Y$ on $X$ fit on the treated units can be written in this form: they are weighted averages of observations $Y_i$ with
weights solving the problem above \eqref{eq_balancing-weights} for $\sigma^2=0$ and $\sigma^2 > 0$ respectively \citep{chattopadhyay2021implied, kallus2016generalized}.
% As we discuss below, there are many variants of these approaches, with different imbalance measures and different weights penalties.
%These methods all differ from the traditional modeling approach to IPW, which typically assesses maximum covariate imbalance as a post-hoc check on the estimated weights; the balancing approach instead incorporates the imbalance constraints directly into the estimation procedure.

The general case \eqref{eq_balancing-weights-general} has a qualitatively similar dual characterization, in which
the inverse propensity weights $\gamma(X_i) = \eta_{\dispersion}\{g(X_i)\}$
are parameterized in terms of an \emph{index function} $g(x)$, and the penalty is imposed on the index. For example, if the model were the set of linear functions with $\ell_1$-bounded coefficients $\{ \beta \cdot x : \norm{\beta}_1 \le 1\}$, this would be a generalized linear model $\gamma(X_i) = \eta_{\dispersion}(\beta \cdot X_i)$ with a penalty scaling with $\norm{\beta}_1$. In practice, $\dispersion$ is often chosen so that the \emph{link function} $\eta_{\dispersion}$ is the reciprocal of a logit link, so when balancing we are fitting a familiar generalized linear model for the propensity score with a sparsity-inducing penalty. See \citet{hirshberg2025bregman} for additional discussion of these dual characterizations.

%including the equivalence of the primal \eqref{eq_balancing-weights} and dual \eqref{eq:dual_l2_penalized_regression} %characterizations of the minimax weights and a dual characterization for the more general balancing weights %\eqref{eq_balancing-weights-general}.

This connection between the details in the specification of a balancing problem and
an implicit model for the inverse propensity score has been discussed at length,
often with a focus on the case that the model $\model$ is a linear one \citep[e.g.,][]{zhao2019covariate}.
These details are less salient with a fully nonparametric model $\model$, as it is not really possible to
get the link $\eta_{\dispersion}$ wrong when the space of functions spanned by our model contains good approximations to all functions $g$. For these models, irrespective of our choice of $\chi$,
the index function $\hat g_\chi$ solving the appropriate generalization of the dual \eqref{eq:dual_l2_penalized_regression} will simply
converge to $g_\chi = \eta_{\chi}^{-1} \circ \gammaipw$, so
$\eta_\chi \circ \hat g_\chi$ will converge to $\gammaipw$.%\footnote{\textcolor{red}{This might be a little less awkward if we did the whole discussion of the minimax special case first, and then duplicated the whole thing for the general case being more explicit about the dual.}}



% In other words, the modeling approach to IPW often assesses maximum covariate imbalance as a post-hoc check on the estimated weights;
% % Equation\ref{eq:sbw}
% the balancing approach instead incorporates the imbalance constraints directly into the estimation procedure.
% The optimization problem is then:
% \begin{equation} \label{eq:sbw}
%   \hgamma = \mathop{\argmin}_{\gamma} \left[\max_{j=1,\ldots,J} \left| %\frac{1}{n} \sum_{i=1}^n X_{ij} - \frac{1}{n}\sum_{i=1}^n W_i \gamma_i %X_{ij} \right|\right]^2 +  \lambda \sum_{W_i=1} \dispersion(\gamma_i).
% \end{equation}
% this approach instead directly uses covariate imbalance to estimate the weights themselves.
% Examples of Equation\ref{eq:sbw} with slightly different constraints on the weights and complexity measure $\dispersion$ include:
% the \emph{stable balancing weights} approach \citep{zubizarreta2015stable};
% the \emph{lasso minimum distance estimator} of the inverse propensity score \citep{chernozhukov2018automatic};
% the \emph{regularized calibrated propensity score estimator} \citep{tan2020regularized}; and
% \emph{entropy balancing} \citep{hainmueller2012balancing}.
% In fact, even linear regression of $Y \sim X$ among treated units can be written in this form \citep{chattopadhyay2021implied}.
% As we discuss below, there are many variants of these approaches, with different imbalance measures and different weights penalties.


% Thus, while both approaches
% The difference is what they are optimizing for the data at hand.
% The modeling approach optimizes the fit of a treatment probability model,
% and in doing so achieves some degree of balance in the covariates.
% Conversely, the balancing approach explicitly targets balance, and in doing so
% estimates the inverse probability weights.

% The loss function in the Lagrangian dual guarantee covariate balance, it also side-steps the issues with first estimating the propensity score and then inverting it by ensuring that the loss is on the correct scale.

% (setting aside the link function).
% In other words, the modeling approach explicitly estimates a regression for the propensity score $e(\cdot)$ in the population; the balancing approach \emph{implicitly} fits a (regularized non-parametric) least-squares regression of $X$ on the \emph{inverse} propensity score.
% minimizes a loss function that estimates the propensity score $e(\cdot)$ in the population, e.g., the empirical likelihood or the mean square error for the treatment indicator $(W_i - \hat{e}(X_i))^2$.




% , with a finite sample correction term arising from the level of imbalance of the true propensity score on the sample.
% (see Section \ref{sec:ipw}).
% the Lagrangian dual problem implicitly fits a (regularized non-parametric) least-squares regression of $X$ on the \emph{inverse} propensity score;
% From the theory of regularized least squares estimators, we can see that the weights will converge to (a projection of) the true inverse propensity scores.


% Comparing to the modeling approach here is instructive. The modeling approach explicitly minimizes a loss function that estimates the propensity score $e(\cdot)$ in the population, e.g., the empirical likelihood or the mean square error for the treatment indicator $(W_i - \hat{e}(X_i))^2$. On the other hand, the balancing approach \emph{implicitly} minimizes a loss function the estimates the \emph{inverse} propensity score $\frac{1}{e(X_i)}$. Not only does minimizing this implicit loss function guarantee covariate balance, it also side-steps the issues with first estimating the propensity score and then inverting it by ensuring that the loss is on the correct scale.

%%%%%%%%%%%%%%%%%%%%%%%%%%%
\section{Asymptotic behavior}
\label{sec:asymptotic}
% , and the conditions under which they are asymptotically Normal.
%
% We now give a brief overview of the asymptotic behavior of these weighting estimators.
% At a high level, a good estimator $\hat \psi_1$ is one that estimates $\psi_1$ precisely and also that allows us to quantify our uncertainty about $\psi_1$, e.g., via confidence intervals.
% In the context of augmented IPW estimators $\hat \psi_1^{\aug}$, as in Equation\ref{eq:aipw2}, this is generally satisfied if
% % And with substantial generality, we can quantify the uncertainty in $\hat \psi_1^{\aug}$ if
% there is negligible imbalance in the outcome regression error $\dm$. We consider two sets of sufficient conditions for achieving these goals. First, under considerable generality, Proposition \ref{proposition:clt} shows that $\hat{\psi}_1^{\aug}$ is asymptotically Normal.
% Second, if we also impose a stronger form of overlap, Proposition  \ref{proposition:optimality} shows that $\hat{\psi}_1^{\aug}$ is also efficient.
% For both approaches, the critical assumption is that the covariate imbalance is (asymptotically) negligible relative to the variance.
%
%

%We now summarize the asymptotic behavior of these weighting estimators to motivate our discussion of model choice in the sections following.

Typically, when imbalance in $\dm$ is asymptotically negligible relative to the estimator's standard error, $\hat{\psi}_1^{\aug}$ is asymptotically normal.
To be more precise, following \citet[see Corollary 6]{athey2018approximate}, when the
weights $\hgamma(X_i)$ are not too concentrated on any one observation $i$ and the noise $\varepsilon_i=Y_i-\mu_{W_i}(X_i)$ is not too heavy-tailed,
negligible imbalance justifies the use of the standard Wald-type confidence interval $\hat \psi_1^{\aug} \pm z_{\alpha/2} \sqrt{\hat V}$, where $\hat V$
is a natural variance estimator and $z_{\alpha/2}$ is a quantile of the standard normal distribution.
We state a result for estimating
$\tilde \psi_1 = \frac{1}{n}\sum_{i=1}^n \mu_1(X_i)$.\footnote{By focusing on  accuracy in estimating $\tilde \psi_1$, a sample analog of $\psi_1$, we drop the sampling variation term from the error decomposition in Equation~\ref{eq:weights_error_aug}. In a sense, because all we
know about the distribution of $X_i$ is summarized by the sample $X_1 \ldots X_n$, it is inevitable that every estimator of $\psi_1$ is
an estimate of $\tilde \psi_1$ as well, and a more precise one at that. See \citet{imbens2004nonparametric}.}

\begin{proposition}
 \label{proposition:clt}
 Let $\Gamma :=  \sqrt{n^{-1}\sum_{i=1}^n W_i \hgamma(X_i)^2}$ be the root-mean-squared weight. Suppose that each squared weight is negligible compared to their sum $n\Gamma^2$, and assume that the noise sequence $\varepsilon_1 \ldots \varepsilon_n$ is uniformly square-integrable conditional on the design. If the imbalance in $\dm$
(see Equation~\ref{eq:weights_error_aug}) is $o_p(n^{-1/2} \Gamma)$, then
 \begin{equation}
 \label{eq:asymptotic-linearity}
 \hat\psi^{\aug} -  \tilde \psi_1 = \frac{1}{n}\sum_{i=1}^n W_i \hgamma(X_i) \{ Y_i - \mu_1(X_i) \} + o_p(n^{-1/2}\Gamma).
 \end{equation}
 Furthermore, if $\hat \mu_1$ is a consistent estimator of $\mu_1$, then
 \begin{equation}
 \label{eq:asymptotic-normality}
  \frac{\hat\psi^{\aug} -  \tilde \psi_1}{\sqrt{\hat V}} \to N(0,1) \quad \text{for} \quad
 \hat V = \frac{1}{n^2}\sum_{i=1}^n W_i \hgamma(X_i)^2 \{Y_i - \hat \mu_1(X_i)\}^2.
 \end{equation}
If imbalance in $\mu_1$ is $o_p(n^{-1/2} \Gamma)$, the same results hold with $\hat \psi^{\text{ipw}}$ in place of $\hat\psi^{\text{aug}}$.
\end{proposition}
\noindent In Appendix~\ref{sec:appendix-asymptotic-normality} we include a more formal and general statement and a simple proof.
In essence, because the noise term dominates the error decomposition \eqref{eq:weights_error_aug} and is
% , conditional on the design $\set{(X_i,W_i) : i \le n}$,
a sum of conditionally independent mean-zero random variables, the estimate satisfies a central limit theorem.%\footnote{TODO: Say something about how we get normality even with poor overlap via self-normalization; big weights appear in numerator and denominator.}
%\avi{rest of this paragraph is hard to parse}
%\dah{yup, not ideal. I tried just giving the root mean squared weight a name. Better? Otherwise tabling this for now, as lost details in Eli's rework made it potentially misleading (sufficient overlap means that the inverse propensity weights have constant order mean square, but in low-but-sufficient-overlap cases probably a lot of model-based estimators don't, so it requires a little more precision) and I haven't been able to write something better.}
We will from this point on focus on the case that the root mean squared weight $\Gamma$ is constant order, so the condition that imbalance is $o_p(n^{-1/2} \Gamma)$ stated in Proposition~\ref{proposition:clt} reduces to it being $o_p(n^{-1/2})$. This is the case when we have sufficient overlap (Assumption~\ref{a:regular-overlap}) and our weights $\hgamma$ are a consistent estimator of $\gammaipw$; we will have more extreme weights when overlap is poor, and can retain a noise-dominated error decomposition with a proportionately more imbalance.


When we assume sufficient overlap, a clear story of optimality emerges.\footnote{Sufficient overlap rules out the possibility that rare treatment assignments have an undue influence on the estimand. It is equivalent to a continuity property: if two conditional mean outcome functions $\mu$ and $\tilde \mu$ are close in the sense that $\E \{\mu_W(X)-\tilde \mu_W(X)\}^2$
is small, then the corresponding treatment-specific means $\psi_1=\E[\mu_1(X)]$ and $\tilde\psi_1 = \E[\tilde \mu_1(X)]$ must be close as well. This can be seen using the equivalent characterization $\psi_1-\tilde \psi_1 = \E[\{W/e(X)\} \{\mu_{W}(X)-\tilde \mu_{W}(X)\}]$; the Cauchy-Schwarz inequality implies that this is bounded by the square root of
$\E\{W/e(X)\}^2 \E\{\mu_{W}(X)-\tilde \mu_{W}(X)\}^2=\E\{1/e(X)\} \E\{\mu_{W}(X)-\tilde \mu_{W}(X)\}^2.$
Without sufficient overlap or strong assumptions on the (semi)parametric form of $\mu_w(x)$, there does not exist a regular estimator of $\psi_1$ \citep{van1991differentiable}.}
{\addtocounter{assumption}{-1}%
   \renewcommand{\theassumption}{\arabic{assumption}$'$}
\begin{assumption}[Sufficient Overlap]
\label{a:regular-overlap}
$\E[1 / e(X)] < \infty.$
\addtocounter{theorem}{-1}
\end{assumption}}
When this holds, all \emph{regular} estimators are asymptotically equivalent (i.e.\ differ by an $o_p(n^{-1/2})$ term) and therefore have the best possible asymptotic variance (see, e.g., van der Vaart, \citeyear{van2000asymptotic}, Chapter 25). In essence, regular estimators are ones that work as expected even if our assumptions about the data's distribution are slightly incorrect. In particular, they are equivalent to an instance of $\hat\psi_1^{aug}$ with the ideal weights $\hgamma=\gammaipw$ and  outcome regression  $\hat \mu_1=\mu_1$.
We get a regular estimator when we achieve negligible imbalance using weights that balance a
sufficiently rich model $\model$. Here rich enough means that it can be inflated enough to contain an arbitrarily good approximation to any square integrable function:
for every function $m$ with $\norm{m}_{L_2(P)} < \infty$ and $\epsilon>0$, there is an approximation $\tilde m \in \alpha \model$ (i.e. $\tilde m$ with $\norm{\tilde m}_\model \le \alpha$) that satisfies $\norm{m-\tilde m}_{L_2(P)} < \epsilon$.
By using a model like this, we control the imbalance in every square integrable function $m$ to some degree (i.e. $o_p(1)$ imbalance), as the maximal imbalance over the inflated model $\alpha \model$ is $\alpha$ times the maximal imbalance over the model itself. This ensures that small errors in our modeling assumptions do not completely change our estimator's behavior. See Appendix~\ref{sec:appendix-asymptotic-normality} for details.

On an intuitive level, optimality also boils down to balance, but in a different function: the inverse-probability-weighted conditional variance $m(x)=\gammaipw(x)v_1(x)$ for $v_1(x)=\Var[Y_i \mid W_i=1, X_i]$.
We can see this by comparing the conditional variance of an estimator with arbitrary weights $\hgamma$ to the one with the ideal weights $\gammaipw$. Using the identity $a^2-b^2=(a-b)^2 + 2(a-b)b$, we can characterize it as proportional to the sum of a variance-weighted mean-squared distance between the two sets of weights and the \emph{difference in imbalance} we get in the function  $f=\gammaipw v_1$ with the two sets of weights.
\begin{equation}
\begin{aligned}
&n \times (V^{\hat\gamma} - V^{\gammaipw}) \\
&= \frac{1}{n}\sum_{i=1}^n W_i \hgamma(X_i)^2 v_1(X_i) - \frac{1}{n}\sum_{i=1}^n W_i \gammaipw(X_i)^2 v_1(X_i)
\\
&= \frac{1}{n}\sum_{i=1}^n W_i \{\hgamma(X_i)-\gammaipw(X_i)\}^2 v_1(X_i) +
\frac{2}{n}\sum_{i=1}^n W_i \{\hgamma(X_i)-\gammaipw(X_i)\} \gammaipw(X_i) v_1(X_i)
\end{aligned}
\end{equation}
where, adding and subtracting the unweighted average of $f$,
\begin{equation}
\begin{aligned}
&\frac{1}{n}\sum_{i=1}^n W_i \{\hgamma(X_i)-\gammaipw(X_i)\} \gammaipw(X_i) v_1(X_i) \\
&= \underbrace{\left\{\frac{1}{n}\sum_{i=1}^n W_i \hgamma(X_i) f(X_i) - \frac{1}{n}\sum_{i=1}^n f(X_i)\right\}}_{\text{imbalance in $f=\gammaipw v_1$ with estimated weights $\hgamma$}} \ - \  \underbrace{\left\{\frac{1}{n}\sum_{i=1}^n W_i \gammaipw(X_i) f(X_i) - \frac{1}{n}\sum_{i=1}^n f(X_i)\right\}}_\text{imbalance in $f=\gammaipw v_1$ with ideal weights $\gammaipw$}.
\end{aligned}
\end{equation}
When we use a model that is a sufficiently rich model in the sense described, both imbalance terms will be $o_p(1)$, so this variance difference effectively reduces to the mean squared distance term. This implies that we cannot improve on the variance we get with the weights $\gammaipw$ without sacrificing balance in the function $f=\gammaipw v_1$. If we reverse the roles of $\hat\gamma$ and $\gammaipw$ in the equation above, we see that if both sets of weights also balance the random function $f(x)=\hgamma(x) v_1(x)$, which follows from control of maximal imbalance together with a mild regularity condition on the weights (see Appendix~\ref{sec:appendix-asymptotic-normality}), then we cannot improve upon the variance achieved with $\hat\gamma$ either. Furthermore, because in this case the variance difference must converge to zero, it follows that the estimated weights $\hat\gamma$ must converge to the inverse probability weights $\gammaipw$.

%Skipitty: \textcolor{red}{I'm starting to get the sense that people don't know this and sometimes get excited about balancing weights having way lower variance than IPW. Is there a way to emphasize this? It's definitely a practical consideration. Can we refer to it there? -Skip}
We can reduce variance by using weights that balance a less rich model. In particular, this is possible when we use a model that we cannot inflate to approximate $\gammaipw v_1$, e.g. if our model contains only linear functions of $x$ and $\gammaipw v_1$ is nonlinear. But this comes with a significant risk: if $\hat \mu_1$ deviates from $\mu_1$ by as little as $1/\sqrt{n}$ in a direction our model cannot approximate, we will get non-negligible bias. We are gaining efficiency by using balancing weights that converge to a \emph{projection} of $\gammaipw$ onto the span of the functions in the model $\model$ rather than $\gammaipw$ itself, and these weights do not balance errors $\hat \mu_1 - \mu_1$ in directions outside this span. By using a model like this, we can get a lower variance than we'd get using the ideal weights $\gammaipw$, but we pay for it with a substantial loss of robustness \citep[Remark 3]{hirshberg2021augmented}.

%Skipitty: COMMENT MAYBE OK IN SMALL SAMPLES BUT NOT IN LARGE ONES [with caveats about inferential double robustness]
%\dah{this is the citation}

%We discuss regularity in greater detail in Appendix~\ref{sec:regularity}.

%%%%%%%%%%%%%%%%%%%%%%%%%%%
%%%%%%%%%%%%%%%%%%%%%%%%%%%
%%%%%%%%%%%%%%%%%%%%%%%%%%%
% \vspace{-.5cm}
\section{Balance measures and models}
\label{sec:balance_models}

%Skipitty: \dah{Revise: be clear that these choices matter, and we'll see in the application/simulation.}

%\avi{Vitor: Put a header paragraph back in here?} \\
%\dah{fwiw, the last paragraph of the section above was not in the draft Vitor read.}
%We now turn to the important question of what to balance.
% and how this is intimately connected to the assumptions on the model class $\model$.
%We begin with the simple setting with only binary covariates, and highlight the importance of limiting higher-order interactions, which is equivalent to imposing restrictions on the model class.
%We then consider the general setting, including continuous covariates,
%and relate restrictions on the model class with balancing certain functions of the covariates.
%Finally, we turn to the question of how well we can balance a given model in practice.

%%%%%%%%%%%%%%%%%%%%%%%%%%%
%%%%%%%%%%%%%%%%%%%%%%%%%%%
\vspace{-.3cm}
\subsection{Balancing binary covariates}
\label{sec:balancing-binary}

%%%%%%%%%%%%%%%%%%%%%%%%%%%
\subsubsection{Models with no interactions}
\label{sec_balancing_marginally}
In the balancing approach of Section~\ref{sec:bal_weights_intro},
we attempt to control imbalance in $\dm$ by
controlling the maximal imbalance
for candidates in a model $\model$.
We now discuss possible choices for this model and connections to balance criteria in the literature. We begin with the case where we have $p$ binary covariates $X_{i1},\ldots,X_{ip} \in \{0,1\}$, for example indicating employment status or high school education. One option is to  balance covariates marginally,
which is equivalent to balancing a linear model without interactions.
This approach is implicit in typical balance checks in observational studies that compare the fraction of units in the treated group with each attribute to the corresponding fraction in the whole sample. In the linear model that puts no
constraints on the coefficients $\beta$, controlling maximal imbalance requires exact marginal covariate balance.
\begin{equation}
\begin{aligned}
\imbalance_{\model}(\gamma) &= \begin{cases} 0 & \text{ if } \ \frac{1}{n}\sum_{i=1}^n X_{ij} - \frac{1}{n}\sum_{i=1}^n W_i\gamma(X_i) X_{ij}  = 0
\text{ for all } j \in 1 \ldots p \\
+\infty & \text{otherwise}
\end{cases} \\
\text{ for }\quad \model &= \set*{ m(x) = \beta \cdot x : \beta \in \R^p }.
\end{aligned}
\end{equation}
In a linear model that constrains the $\ell_1$ norm of the coefficients,  maximal imbalance is the largest maximal imbalance in any one covariate.
\begin{equation}
    \label{eq:imbalance_margin_cat}
    \begin{aligned}
   \imbalance_{\model}(\gamma) &= \max_{j \in 1 \ldots p} \abs*{\frac{1}{n}\sum_{i=1}^n X_{ij} - \frac{1}{n}\sum_{i=1}^n W_i\gamma(X_i) X_{ij} }  \\
   \text{ for }\quad \model &= \set*{ m(x) = \beta \cdot x : \norm{\beta}_1 \le 1 } \\
\end{aligned}
\end{equation}
Finding weights to balance marginal proportions has a long history in survey statistics, where it is known as \emph{raking} \citep[e.g.,][]{Deming1940, deville1992calibration}.
In causal inference such balance measures are commonly used to assess covariate balance across treatment groups \citep[e.g.,][]{cochran1965planning, rosenbaum1985constructing, rosenbaum2007minimum}. Generally, if the conditional mean outcome $\mu_1$ is a linear function $x \cdot \beta$ of the covariates, its bias conditional on the design $\mathcal{D} = \{ (W_1, X_1) \ldots (W_n,X_n) \}$ is
\begin{equation}
  \label{eq:bias_linear}
  \beta \cdot \left(\frac{1}{n}\sum_{i=1}^n X_i - \frac{1}{n}\sum_{i=1}^n W_i\hgamma(X_i) X_i\right).
\end{equation}
Thus, weights that exactly balance each covariate marginally will produce an unbiased estimator if $\mu_1(x)=x \cdot \beta$ for any coefficients $\beta$. If these coefficients satisfy an $\ell_1$ norm constraint $\norm{\beta}_{\ell_1} \le B$,
H\"older's inequality implies this conditional bias is bounded by $B$ times  $\imbalance_\model(\gamma)$ from Equation~\ref{eq:imbalance_margin_cat}. And more generally, if they satisfy an $\ell_q$ norm bound $\norm{\beta}_{\ell_q} \le B$ on the coefficients, H\"older's inequality implies that this bias is bounded by the $\ell_{q'}$ norm of the vector of all marginal imbalances where $q'$ is the H\"older conjugate of $q$. Henceforth, we will focus on models that constrain the $\ell_2$ norm of the coefficients, which corresponds to a maximal imbalance that is the square root of the sum of all squared marginal imbalances.





However, unless $\mu_1$ is linear, balancing the margins alone is insufficient to remove bias.
%To see this, let $\phi(x)$ be the vector of covariates that we are balancing; for now $\phi(x) = x$, but below we will consider other transformations of the covariates.
To see this, consider the decomposition of
$\mu_1$ into its best linear approximation $\tilde \mu_1(x) := \tilde\beta \cdot x$ and the residuals
$r(x) := \mu_1(x) - \tilde \mu_1(x)$.
Balancing margins leaves additional \emph{approximation error} based on the difference between the approximation $\tilde \mu_1$ and $\mu_1$ itself. Generally, the design-conditional bias is
\begin{equation}
  \label{eq:bias_basis}
  \tilde \beta \cdot \underbrace{\left(\frac{1}{n}\sum_{i=1}^n X_i - \frac{1}{n}\sum_{i=1}^n W_i\hgamma(X_i) X_i\right)}_{\text{marginal imbalance}} \;\;+\;\; \underbrace{\left(\frac{1}{n}\sum_{i=1}^n r(X_i) - \frac{1}{n}\sum_{i=1}^n W_i\hgamma(X_i) r(X_i)\right)}_{\text{imbalance in the approximation error}}.
\end{equation}
We can reduce approximation error by balancing a \emph{basis expansion} $\phi(x)$ that includes interactions; however, this leaves us with two competing modeling goals.
We would like to balance an expressive enough model to give a good approximation of $\mu_1(x)$ and also like to ensure that we are able to get reasonable balance for that model.

%%%%%%%%%%%%%%%%%%%%%%%%%%%
\subsubsection{Models with strong interactions}
\label{sec_balancing_all}

Balancing models with interactions is the best course of action.
For example, we might believe that the implications of being both employed \emph{and} a high school graduate is more than the sum of its parts, so we would want to balance the interaction between the two.
More formally, for our vector of $p$ binary covariates $X_{i} = (X_{i1}, \ldots, X_{ip})$ and multi-index vector $j = (j_1, ..., j_p) \in \{0,1\}^p$,
let $X_{i}^{j}$ be the covariate interaction $X_{i1}^{j_1} \cdot \ldots \cdot X_{ip}^{j_p}$, so the interaction of being employed and a high school graduate
corresponds to an index vector $j$ with ones in those two coordinates, say with $j_1=j_2=1$, and zeros elsewhere.  In this notation, letting sums over $j$ range over $\{0,1\}^p$ implicitly, when $\mu_1(x)=\sum_j \beta_j x^j$ design-conditional bias is
\begin{equation}
  \label{eq:bias_interactions}
  \sum_j \beta_j \left(\frac{1}{n}\sum_{i=1}^n X_i^j - \frac{1}{n}\sum_{i=1}^n W_i\hgamma(X_i) X_i^j\right).
\end{equation}
What must be done to control this bias is analogous to what we saw in our discussion of imbalance measures for linear-in-covariates models \eqref{eq:bias_linear}. To control bias without assuming any constraint on the coefficients $\beta$, we must have exact balance in each interaction $X_i^j$,
and H\"older's inequality implies bounds on this bias for coefficients satisfying $\ell_q$ norm bounds. For example,
\begin{equation}
    \label{eq:imbalance_interact}
    \begin{aligned}
   \imbalance_{\model}^2(\gamma) &=  \sum_{j} \abs*{
   \frac{1}{n} \sum_{i=1}^n X_{i}^{j} - \frac{1}{n}\sum_{i=1}^n W_i \gamma(X_i) X_{i}^{j}
   }^2 \\
   \quad \text{ for } \quad \model &= \set{ m(x) = \sum_{j} \beta_j x^j : \sum_{j} \beta_j^2 \le 1}.
   \end{aligned}
\end{equation}
For low-dimensional covariates, this type of balance measure is commonly used in sample surveys, with exact balance of fully-interacted models being called \emph{post-stratification}.
%proposed a related method based on least-squares in the context of sample surveys,}.%
% \avi{think we can remove one of the next two sentences}
% This is the case
% In particular, because the model class $\model$ in
% Equation~\ref{eq:imbalance_interact}
% % is general (in the sense that it
% spans every function of our binary covariates, setting the imbalance in Equation~\ref{eq:imbalance_interact} to zero implies perfect balance in all possible functions of the covariates.
%5In other words, when we have categorical covariates, we can remove bias for every possible conditional mean function $\mu_1$ by finding weights that \emph{exactly} balance every interaction --- so long as such weights exist.


%%%%%%%%%%%%%%%%%%%%%%%%%%%
\subsubsection{Models with weak interactions}

The challenge with balancing higher-order interactions is that the set of all interactions to balance grows exponentially in the number of covariates, making it infeasible to balance everything well.
Thus, the curse of dimensionality often prevents us from exactly balancing all higher-order interactions.
Even if we allow for approximate balance,
the imbalance measure in Equation~\ref{eq:imbalance_interact} gives the same importance to the marginal imbalance in a covariate and the imbalance in the $p^{th}$ order interaction $X_{i1} \cdot X_{i2} \cdot ... \cdot X_{ip}$.
This is appropriate if we think the effects of high-order interactions may be as large as the main effects, but models in which effects tend to be smaller for higher order interactions tend to be more plausible.
%Skipitty: \textcolor{red}{language: models ($\model$) aren't plausible or implausible. Did we mean $\mu$ where we said model here?}


To encode the assumption that main effects
are larger than the effects of first-order interactions, which are larger than those for second-order interactions, and so on, %For more on these ``deep'' interactions, see \citet{Ghitza2013}.
we can control the maximal imbalance over a model in which the coefficients are in an ellipse that narrows for higher-order interactions, so the coefficient on $X_i^j$ must be less than a corresponding radius $\lambda_j^{-1/2}$ that decreases as a function of the interaction order $\sum_{k=1}^p j_p$. This measure of imbalance is used, for example, in \citet[Equation 9]{benmichael2021_multical}.
\begin{equation}
    \label{eq:binary-rkhs}
\begin{aligned}
\imbalance^2_{\model}(\gamma) &= \sum_{j} \lambda_j \left|\frac{1}{n}\sum_{i=1}^n X_i^j - \frac{1}{n}\sum_{i=1}^nW_i \gamma(X_i) X_i^j\right|^2 \\
\quad \text{for} \quad \model &= \set*{m(x) = \sum_j \beta_{j} x^j : \sum_j \beta_j^2 \ /\ \lambda_j \le 1}.
\end{aligned}
\end{equation}

%It differs from the previous model \eqref{eq:imbalance_interact} in that
%the strength of an effect $\beta_j$ is penalized at a rate $1/\lambda_j$ that increases with the order of the corresponding interaction $x^j$.%\footnote{Here we refer loosely to the idea of strengthening an effect being `penalized' in sense
%that it increases the gauge $\norm{\cdot}_{\model}$, moving the function away from the center of the class. To be precise,
%the rate at which an effect is penalized in the sense meant here is the square root of the
%rate of increase of the squared gauge $\norm{\cdot}_{\model}$ with an increase in that effect.}

%This is where restrictions on the underlying model $\model$ start to play a more important role.
% For instance, a reasonable assumption in many settings is that the coefficients corresponding to ``main effects'' (i.e., balancing the covariates marginally)
%For instance,
Taking $\lambda_j=0$ corresponds to assuming away the presence of an effect of $x^j$, so we do not penalize imbalance in the corresponding interaction, whereas taking $\lambda_j=\infty$ corresponds to assuming  no limits on the effect, so we must balance the corresponding interaction
perfectly. In Section~\ref{sec_balancing_marginally} above, we discussed the case in which $\lambda_j$ is nonzero only for main effects.

%\eli{The discussion after this is confusing to me. Post stratification exaclty balances all higher order terms, and I thought the connection to fine balance was balancing the marginals (e.g. no interactions)}
%\dah{Jose has convinced me that none of these terms have model-independent meanings, so unless we want to get into the weeds, we just sort of have to pick a meaning and try to do it so we don't come off criticising people who were doing something application appropriate for not specifying what to do in complete generality}
%One common approach is to adjust exactly for main effects and low-order interactions and not at all for higher-order ones.
%, known as \emph{post-stratification} \citep{Deming1940}
%in sample surveys.
%Sometimes exact balance on low-order terms is combined with approximate balance on higher order ones by taking $\lambda_j=\infty$ only for small $\norm{j}_1$;
%this is called \emph{fine balance} in the literature on matching in causal inference \citep{rosenbaum2007minimum}.
% is refined balance one of these things?

We can get a qualitatively similar model by restricting our coefficients to any shape that is more elongated in some directions than others, e.g. a rectangle, which corresponds to a bound on  $\norm{ \lambda^{-1/2} \beta }_{\ell_q}$ for $q=\infty$. The ellipse is a common choice because, as with the use of $\ell_2$ norms in other contexts such as least squares vs. least absolute deviations, it simplifies computation and analysis.

\subsubsection{Nonbinary discrete and coarsened continuous covariates}

The discussion above can be applied to discrete-valued but non-binary covariates
using the common practice of \emph{one-hot encoding}, i.e., expanding a $k$-valued covariate $X_{ij}$ into a vector of $k-1$ binary covariates $1(X_{ij}=2) \ldots 1(X_{ij} = k)$. However, if these discrete covariate are coarsened versions of continuous covariates, we may be missing an important source of bias. Our bias characterization \eqref{eq:bias_interactions} is based on the assumption that $\mu_1$ is a function of the coarsened covariates
rather than the underlying continuous-valued ones. For this reason, it can be better to work directly with models for functions of continuous-valued covariates.

%\textcolor{red}{The 'our bias characterization...' sentence is a little unclear. Can we do better?}

%%%%%%%%%%%%%%%%%%%%%%%%%%%
%%%%%%%%%%%%%%%%%%%%%%%%%%%
% \vspace{-.5cm}
\subsection{Balancing continuous covariates}
\label{sec:balancing-continuous-covariates}


%\avi{header paragraph here?}

\subsubsection{The basis expansion perspective}

The key change when we work with continuous covariates is that interactions are not the only source of complexity: functions of even one continuous covariate can be complex and nonlinear. In essence, we face the curse of dimensionality even in one dimension, as non-parametric models require an \emph{infinite} basis. For example, the natural generalization of the previous model \eqref{eq:binary-rkhs} is a Taylor series model, in which we sum over all powers of each covariate ($j \in \mathbb{N}^p$), instead of the binary ones ($j \in \{0,1\}^p$). In this model, we can think of higher-order terms as being the ones for which the vector $j$ is large, e.g. both $x_1 \cdot \ldots \cdot x_4$ and $x_1^4$, whether they are interactions or powers of one covariate.
More generally, we can use a similar model specified in terms of any set of basis functions $\phi_1(x),\phi_2(x),\ldots$ and corresponding radii $\lambda_1^{-1/2},\lambda_2^{-1/2},\ldots$.
\begin{equation}
    \label{eq:imbalance_phi_infinite}
\begin{aligned}
    \imbalance^2_{\calM}(\gamma) &= \sum_j \lambda_j \abs*{\frac{1}{n}\sum_{i=1}^n \phi_j(X_i) - \frac{1}{n}\sum_{i=1}^n W_i \gamma(X_i) \phi_j(X_i)}^2 \qquad \text{ for } \\
    \model &= \set*{m(x) = \sum_j \beta_{j} \phi_j(x) : \sum_j \beta_j^2 \ / \ \lambda_j \le 1}.
\end{aligned}
\end{equation}
To better isolate first order effects from higher order ones in a sequential manner, it is common to use a basis of \emph{orthogonal polynomials}
instead of the Taylor series-inspired basis above; for example, the Hermite polynomials $1,x_1,x_2,x_1^2-1,x_2^2-1,x_1x_2$, written here for two covariates.
%This isolation, which makes quadratic terms `pure quadratic,' etc., can make specifying  $\lambda_j$ a bit more intuitive.
This isolation can make specifying  $\lambda_j$ a bit more intuitive.
Another option, motivated by Fourier series, is to use an orthogonal basis of cosines of different frequencies. We cannot balance cosines of all frequencies accurately just like we cannot balance all polynomials, so we use models that restrict effects associated with high frequency terms.

As before, taking $\lambda_j=0$ corresponds to assuming away any effect for $\phi_j(x)$, so  $\imbalance_\model(\gamma)$ is unaffected by imbalance in $\phi_j$, and taking $\lambda_j=\infty$ corresponds to assuming no limits on that effect and therefore requires exact balance in $\phi_j$.
%coordinate. \avi{do we need this last sentence? Do we mean that this corresponds to a linear model *without* any constraints on the coefficients?} The common practice of balancing the covariates marginally corresponds to a linear model $(m(x) = \beta \cdot x$), much as in Section~\ref{sec_balancing_marginally}.

\subsubsection{The function space perspective}
\label{sec:function-space-perspective}
While the basis expansion perspective can be useful when thinking about the balancing problem, specifying an infinite-dimensional model $\model$
directly in terms of basis functions and scale factors ($\phi_j, \lambda_j$) is hard to do well. We can also view things from another perspective,
focusing on the properties of the functions rather than their basis expansion.
The theory of reproducing kernel Hilbert spaces (RKHSes) provides one approach for doing so. Subject to some constraints, the model $\model$ in Equation~\ref{eq:imbalance_phi_infinite} is the unit ball $\{m(x) : \norm{m}_\model \le 1\}$ of an RKHS,
and furthermore, the unit ball of every RKHS can be described like this  \citep[see, e.g.,][Chapters 2 and 4]{cucker2007learning}.\footnote{Specifically, these constraints are that the basis $\phi_j$ is orthonormal with respect to some probability measure and the decay factors $\lambda_j$ have a finite sum.}  %\avi{one more sentence here defining what this means?}.



%\avi{just noting that most of the real estate for this subsection is devoted to Sobolev spaces (partly because the formulas take up a lot of room!), while zooming past RKHS terms. Want to make sure this is a deliberate choice.}\\
%\dah{it was. Most other RKHSes are very hard to interpret.}\\
One of the most common ways to characterize functions is their smoothness. For this, it is natural to work with Sobolev spaces: RKHSes defined in terms of averages of a function's derivatives. We can characterize the unit ball of a \emph{Sobolev norm} in terms of the coefficients of Fourier series expansions. For example, these two Sobolev norms for functions on the unit square $[0,1]^2$, which we can characterize equivalently in terms of the coefficients $\beta_j$ in the function's cosine basis expansion $m(x) = \sum_{j \in \mathbb{N}^2} \beta_j \cos(\pi j_1x_1)\cos(\pi j_2x_2)$.
\begin{align}
\norm{m}_{\text{iso}}^2
&=  \int_{[0,1]^2} \ \abs*{m(x_1,x_2)}^2 + \abs*{\frac{\partial }{\partial x_1}m(x_1,x_2)}^2 +
     \abs*{\frac{\partial }{\partial x_2}m(x_1,x_2)}^2  \label{eq:sobolev-derivative} \\
&= \sum_{j \in \mathbb{N}^2} \beta_j^2/\lambda_j \text{ with } 1/\lambda_j = 1+\pi^2j_1^2 + \pi^2j_2^2 \label{eq:sobolev-fourier}. \\
\norm{m}_{\text{dom}}^2
&=  \norm{m}_{\text{iso}}^2 + \int_{[0,1]^2} \abs*{\frac{\partial^2 }{\partial x_1 \partial x_2}m(x_1,x_2)}^2 \label{eq:ms-sobolev-derivative} \\
&= \sum_{j \in \mathbb{N}^2} \beta_j^2/\lambda_j \text{ with } 1/\lambda_j = 1+\pi^2j_1^2 + \pi^2j_2^2 + \pi^4j_1^2j_2^2 \label{eq:ms-sobolev-fourier}
\end{align}

%The maximal imbalance for its unit ball,  $\model = \set{m(x) : \norm{m} \le 1 }$, is
%\begin{equation}
% \label{eq:imbalance-sobolev}
%   \imbalance^2_{\calM}(\gamma) = \sum_j %\frac{\abs*{\frac{1}{n}\sum_{i=1}^n %\{1-W_i\gamma(X_i)\} \cos(j_1 X_{i1} + %j_2 X_{i2})}^2}{(1+j_1^2)(1+j_2)^2}.
%\end{equation}
%  \eli{Shouldn't we have a parallel structure here where we write the imbalance term and the model class? It would help to make this more concrete and ground the section, I think.}
% \dah{I wrote this and commented it out above. I think it's worse with it. Inviting readers to puzzle over what a sobolev-ball imbalance means in terms of basis functions is a bit off-message in a section that says you don't have to and probably shouldn't think about basis expansions all the time.}
In these models, our assumptions are encoded in the derivatives included in the norm.
If we expect $\mu_1$ to be relatively smooth, including higher order derivatives would encode this assumption.
If we believe that most of the complexity lies along the coordinate axes, i.e., that a function tends to vary more with one covariate than with a mixture of a few, including higher order mixed partial derivatives, as we do in $\norm{m}_{\text{dom}}$, encodes this assumption.
%\color{red}
Additive models take this idea to an extreme, entirely
forbidding variation that is not simply a sum of what we see along the axes.\footnote{We get a series expansion of the set of functions of the form $\sum_k m_k(x_k)$ with $\norm{m}_{\text{iso}} \le 1$ by making $1/\lambda_j=\infty$ for terms involving mixed partials rather than simply increasing it as we do for $\norm{m}_{\text{dom}}.$}

%The reason to use a smaller vs. larger model, in a regression context, is that we have to if we want to avoid overfitting...we will see in the next section that the same is essentially true when specifying a model to balance...



%These models, which are a sort of intermediate choice between additive models --- which we would get by making $\lambda_j=\infty$ for terms involving mixed partials ($\sum_k 1(j_k > 0) = 1$) and isotropic models like the unit ball of $\norm{\cdot}_{\text{iso}}$, are called \emph{dominating mixed smoothness} Sobolev models.

% We have talked about functions on $\R^2$ only to simplify notation: these concepts readily extend to higher dimensions.
% The topic of nonparametric models is huge...
% FIX:Huge topic ... see appendix Appendix~\ref{sec:more-on-models}.
So far, we have focused on functions defined on $\mathbb{R}^2$ for notational simplicity. However, these concepts naturally extend to higher-dimensional settings.
The subject of nonparametric models is vast; see Appendix~\ref{sec:more-on-models} for further discussion.


%\color{black}


\begin{remark}[Invariance and equivariance]
\label{sec:translation-invariance}

The average treatment effect $\tau = \psi_1 - \psi_0$ is translation invariant; it is not affected by the way the outcomes are centered. It aids our intuition to use estimators $\hat\psi_1$ and $\hat\psi_0$ with a corresponding \emph{equivariance property}: if our estimate $\hat\psi_w$ is $a$ when $\mu_w(x)=m(x)$, then it should be $a + b$ when $\mu_w(x)=m(x) + b$. The IPW and AIPW estimators have this property when the weights $\hat\gamma$ average to one, and it is common practice to ensure translation invariance by normalizing them to average to one, using $\hat\gamma(X_i) = \hat e(X_i)^{-1} / \{ n^{-1} \sum_{i=1}^n W_i \hat e(X_i)^{-1}\}$ instead of $\hat \gamma(X_i)=\hat e(X_i)^{-1}$.
This equivariance property arises naturally from our balance criterion when our model includes all translations of every function in it, i.e., if $m(\cdot) + t \in \model$ whenever $m \in \model$. In models like this, the imbalance is infinite if we do not have equivariance.
  \[
   \begin{aligned}
   &\imbalance_{\model}(\gamma) =
     \begin{cases}
       \imbalance_{\model_0}(\gamma) & \text{ if }\ \frac{1}{n}\sum_{i=1}^n W_i \gamma(X_i) = 1\\
       \infty &  \text{ if }\ \frac{1}{n}\sum_{i=1}^n W_i \gamma(X_i) \neq 1
   \end{cases}
   \quad\ \text{ for } \\ &\quad \model = \set*{ m(x) + t : m \in \model_0,\  t \in \R}.
  \end{aligned}
   \]
We can add these translations to the Sobolev models described above by simply dropping the term without derivatives from the corresponding norm, e.g. $\model = \{ m \ : \rho_{\text{iso}}(m) \le 1 \}$ for
\begin{align}
\rho_{\text{iso}}(m)^2
&=  \int_{[0,1]^2}  \abs*{\frac{\partial }{\partial x_1}m(x_1,x_2)}^2 +
     \abs*{\frac{\partial }{\partial x_2}m(x_1,x_2)}^2  \label{eq:sobolev-seminorm-derivative}  \\
&= \sum_{j \in \mathbb{N}^2} \beta_j^2/\lambda_j \text{ with } 1/\lambda_j = \pi^2j_1^2 + \pi^2j_2^2.  \label{eq:sobolev-seminorm-fourier}
\end{align}
\end{remark}

\subsection{Other considerations}
\label{sec:other-considerations}

Constraining the weights to be non-negative and average to one disallows \emph{extrapolation} outside of the convex hull of the data \citep{hirshberg2025bregman}. Intuitively, using weights that interpolate the data limits the extent of biases due to misspecification of the model $\mathcal{M}$ relative to what is possible with extrapolation. In particular, this ensures that the estimate is sample bounded, i.e.\ that it lies between the minimum and maximum observed values of the outcome.
This rules out nonsensical estimates when outcomes are known to be bounded, such as for binary and survival outcomes.

Such constraints can be enforced by choosing a weight complexity penalty $\dispersion$ in Equation~\ref{eq_balancing-weights-general} that blows up outside the allowed range. A quadratic penalty $\dispersion(\gamma) = \gamma^2$ on the domain $[0, \infty)$ enforces non-negativity. The negative entropy penalty $\dispersion(\gamma) = \gamma \log \gamma - \gamma$ enforces $\gamma \ge 0$, and $\dispersion(\gamma) = \gamma \log \gamma - \gamma + 1$ enforces $\gamma \ge 1$.

However, because balancing weights for rich models converge to the inverse propensity weights, and they exceed one, we tend to get negative weights when we are balancing a small---and likely to be misspecified---model.
When this is the case, using a richer model can be a better fix than a non-negativity constraint or, at least, a good complement.

The parameter $\sigma^2$ in the minimax weights \eqref{eq_balancing-weights} is often described as controlling a bias-variance trade-off: a small $\sigma^2$ emphasizes balance and therefore control of bias, and a large $\sigma^2$ emphasizes control of the magnitude of the weights and therefore of variance. The minimax interpretation of the weights discussed in Section~\ref{sec:bal_weights_intro} suggests choosing $\sigma^2$ to be roughly the conditional variance of the outcomes, $\Var[Y_i \mid X_i, W_i]$. This choice is supported by asymptotic theory \citep[e.g.,][]{hirshberg2021augmented, hirshberg2019minimax}, which suggests that taking $\sigma^2$ to be constant order is a relatively robust choice.

But $\sigma^2$ is fundamentally a signal-to-noise ratio, comparing imbalance in model-scale units to outcome noise. The minimax heuristic assumes a particular model scale; changing the units of $\model$ changes what $\sigma^2$ should be. Sometimes we have intuition about model scale. For example, in a medical study, a domain expert may say treatment groups are comparable in age if they differ by no more than a few months---that pins down the scale of the age component of $\model$.\footnote{It can be more intuitive to minimize the magnitude of the weights subject to a constraint on imbalance, which is equivalent. There exists some data-dependent choice $\sigma^2$ that yields a solution of the constrained problem for each bound on the imbalance and vice versa \citep[Chapter 5]{boyd2004convex}.}

For rich models, the effect of $\sigma^2$ is limited anyway: weights converge to the inverse propensity weights regardless (e.g., Equation~\ref{eq_balancing-weights-comparison}). Extreme choices matter in small samples, but there is a range of reasonable values. How to select this parameter automatically is an open question. Some approaches have been proposed \citep{kallus2020generalized, wang2020minimal, zhao2019covariate} but theoretical and empirical evidence supporting them is currently limited.

%%%%%%%%%%%%%%%%%%%%%%%%%%%
% \vspace{-.5cm}
\section{Degree of balance}
\label{sec:trade-offs}

The discussion above focuses on balancing increasingly large, expressive models. Why would we choose a smaller model over a larger one when the latter is more likely to be correctly specified? The answer lies in a fundamental trade-off: while larger models can capture more complex relationships, we cannot balance them as well, and negligible imbalance in the regression error (in the sense described in Section~\ref{sec:asymptotic}) is crucial for inference.\footnote{However, see \citet{armstrong2018finite} for a discussion of bias-aware confidence intervals.} In this section, we will focus on understanding when this happens.

\subsection{Comparing to the true inverse propensity weights}
\label{sec:plugin}

We first focus on the unaugmented estimator, so $\delta m = \mu_1$,
and assume our model is valid in the sense that $\mu_1 \in \model$; thus,
imbalance in $\mu_1$ is bounded by $\imbalance_{\model}(\hgamma)$. We can bound this by comparing it to $\imbalance_\model(\gammaipw)$, the maximal imbalance achieved by the  true inverse propensity weights. For example, because the minimax weights $\hgamma$ minimize the objective function in Equation~\ref{eq_balancing-weights},
\begin{equation}
\label{eq_balancing-weights-comparison}
\imbalance_{\model}^2(\hgamma) \;\;+\;\; \frac{\sigma^2}{n^2} \sum_{W_i=1}\hgamma(X_i)^2 \le
\imbalance_{\model}^2(\gammaipw) \;\;+\;\; \frac{\sigma^2}{n^2} \sum_{W_i=1}\gammaipw(X_i)^2.
\end{equation}
As \citet{hirshberg2021augmented} discuss, because the estimated weights $\hgamma$ converge to the true inverse propensity weights $\gammaipw$, the difference between the variance terms is asymptotically negligible. So up to a negligible term, $\imbalance_{\model}(\hgamma) \le \imbalance_\model(\gammaipw)$, and we can focus on characterizing the latter.

Because the inverse propensity weights have the population balance property~\eqref{eq_mipw-balance},
$\imbalance_{\model}(\gammaipw)$ is the maximum, over all $m \in \model$, of an average of iid mean-zero terms. Consider the model $\model = \{\beta \cdot \phi(x) : \|\beta\|_1 \leq 1\}$ for a basis expansion $\phi(x) \in \R^p$. The maximal imbalance is the imbalance in the least-balanced basis function:
\begin{equation}
  \label{eq:ipw_linf_imbalance}
  \imbalance_\model(\gammaipw) = \max_{j=1,\ldots,p} \left| \frac{1}{n}\sum_{i=1}^n \phi_j(X_i) -  W_i \gammaipw(X_i) \phi_j(X_i)\right|.
\end{equation}
By the central limit theorem, each individual imbalance is approximately normal with standard deviation of order $n^{-1/2}$. So the maximal imbalance behaves like the maximum of $p$ mean-zero normals, which is $O_p(\sqrt{\log(p)/n})$---depending only weakly on the number of basis functions. \citet{athey2018approximate} use this approach to analyze balancing weights for high dimensional linear models.

We can extend this to infinite-dimensional models. The maximum of $p$ normals grows like $\sqrt{\log p}$ only when they are sufficiently uncorrelated. If the normals become increasingly correlated as $p$ grows---recycling variation we've already seen---the growth stops. Models $\model$ that are ``small enough'' for this to happen, i.e.\ models for which $\imbalance_\model(\gammaipw)=O_p(n^{-1/2})$, are called Donsker classes in empirical process theory.\footnote{To be more precise, $\imbalance_\model(\gammaipw)=O_p(n^{-1/2})$ when $\model^{ipw} = \{ (1-W\gammaipw)m \ : \ m \in \model \}$ is Donsker. If the \emph{strict overlap} condition that $\gammaipw$ is bounded holds, this is true if $\model$ itself is Donsker, but it can be stronger otherwise. See \citet[Remark 4]{hirshberg2021augmented}.} One example, discussed in \citet{kallus2020generalized} and \citet{wong2018kernel}, is the unit ball of a Sobolev space with $s > p/2$ bounded derivatives.

\paragraph{Symmetrization and the connection to regression.}
The same measures of model complexity that govern the difficulty of nonparametric regression also govern balance. The intuition above---maximal imbalance as a maximum of correlated normals---can be made precise. Rather than reasoning asymptotically about what the individual imbalances converge to, we can work with a symmetrized version that \emph{conditionally} has the structure we want. Letting $\xi_1, \ldots, \xi_n$ be independent and identically distributed standard normals, independent of the data, we define the \emph{symmetrized imbalance}
\begin{equation}
  \label{eq:ipw_gen_imbalance_symmetrized}
  \imbalance^\xi_{\model}(\gammaipw) = \max_{m \in \model} \left| \frac{1}{n}\sum_{i=1}^n \xi_i \{ m(X_i) -  W_i \gammaipw(X_i) m(X_i) \}\right|.
\end{equation}
Conditional on the data, this is normally distributed. And it's a good approximation of imbalance itself:\footnote{With Gaussian multipliers $\xi_i$, the expected value of maximal imbalance is at most  $\sqrt{2\pi}\approx 2.5$ times larger than
its symmetrized counterpart's and at least $1/(4\sqrt{\log(n)})$ times as large---almost equivalent up to constant factors. The expectation of this symmetrized imbalance is called the Gaussian complexity of the inverse probability-weighted model
$\model^{ipw}$. Using random signs, i.e. independent $\xi_i$ taking on the values $\pm 1$ each with probability $1/2$,
we can get true equivalence---the ratio of these expectations is between $1/2$ and $2$.  With these multipliers, the expectation is called the Rademacher complexity of $\model^{ipw}$. See \citet{bartlett2002rademacher} for details.}
\begin{equation}
  \label{eq:symmetrization-bounds}
   \frac{1}{4\sqrt{\log(n)}} \le \frac{\E \imbalance_{\model}(\gammaipw)}{\E \imbalance^\xi_{\model}(\gammaipw)} \le \sqrt{2\pi}.
\end{equation}

The symmetrized imbalance measures how well functions in our model can correlate with noise. This is exactly the quantity that governs the difficulty of nonparametric regression. In least squares, the estimation error of $\hat\mu$ depends on how large the noise correlation
$\frac{1}{n}\sum_{i=1}^n \varepsilon_i \{ \mu(X_i) - m(X_i) \}$
can be relative to the squared error $\norm{\mu - m}_{L_2}^2$. To see this, write the loss difference between an arbitrary $m \in \model$ and the true regression function $\mu$:
\begin{equation}
\begin{aligned}
\label{eq:loss-difference}
\ell(m)
&=  \frac1n \sum_{i=1}^n \{Y_i - m(X_i) \}^2 - \frac1n\sum_{i=1}^n\{Y_i - \mu(X_i) \}^2 \\
&= \norm{\mu - m}_{L_2(\Pn)}^2 - \frac2n \sum_{i=1}^n \varepsilon_i \{ \mu(X_i)-m(X_i) \}
\end{aligned}
\end{equation}
where $\varepsilon_i = Y_i - \mu(X_i)$ is the regression error. The least squares estimator $\hat\mu$ minimizes the loss, so $\ell(\hat\mu) \le 0$. This means the squared error $\norm{\mu - \hat\mu}^2$ can be at most twice the noise correlation term. Turning this around: if we can show that the noise correlation is small for all $m$ at distance $r$ from $\mu$, we can rule out $\norm{\hat\mu - \mu} = r$. Concretely,
\begin{equation}
\label{eq:least-squares-fixed-point}
\norm{\hat \mu - \mu} \neq r \quad \text{if} \quad
\frac{r^2}{2} \ge  \max_{\substack{m \in \model \\ \norm{\mu - m}_{L_2(\Pn)} = r}} \frac1n\sum_{i=1}^n \varepsilon_i \{ \mu(X_i) - m(X_i) \}.
\end{equation}

\begin{figure}[htbp]
\centering
% Row 1: LS model (blue)
\makebox[0.31\textwidth]{\includegraphics[width=0.31\textwidth]{figure/geometry-handpicked-ls-1.pdf}}%
\hfill
\makebox[0.31\textwidth]{\includegraphics[width=0.31\textwidth]{figure/geometry-handpicked-ls-2.pdf}}%
\hfill
\makebox[0.31\textwidth]{\includegraphics[width=0.31\textwidth]{figure/geometry-handpicked-ls-3.pdf}}%

\vspace{2pt}

% Row 2: IPW model (vermillion)
\makebox[0.31\textwidth]{\includegraphics[width=0.31\textwidth]{figure/geometry-handpicked-ipw-1.pdf}}%
\hfill
\makebox[0.31\textwidth]{\includegraphics[width=0.31\textwidth]{figure/geometry-handpicked-ipw-2.pdf}}%
\hfill
\makebox[0.31\textwidth]{\includegraphics[width=0.31\textwidth]{figure/geometry-handpicked-ipw-3.pdf}}%

\vspace{-2pt}% Row 3: Shared crossing plot (full width)
\includegraphics[width=\textwidth]{figure/geometry-handpicked-crossing.pdf}
\caption{Geometry of the fixed-point argument for two models: the regression model $\model$ (blue, top row) and the IPW-reweighted model $\model^\ipw$ (vermillion, middle row). Both rows show the same noise draw; the models differ only in the tube's slope and width, reflecting the coordinate-wise rescaling by propensity score factors. As the ball shrinks (right to left), the maximal noise correlation shrinks too---this is equicontinuity. But in most models---linear ones are an exception---it shrinks at a rate slower than the ball's radius does. What happens is that, as our ball gets smaller, it does more and more of the work in constraining $m-\mu$. In the extreme case shown on the left, where the ball is contained in the model, the model is doing nothing at all. The bottom panel plots the expected maximal correlation (Gaussian width) as a function of the ball's radius $r$ against the curve $r^2/2$; the \emph{fixed point} where they cross bounds the regression error $\norm{\hat\mu-\mu}$ or the analogous balancing error. The narrower IPW model crosses earlier, reflecting the additional structure imposed by reweighting.}
\label{fig:width-geometry}
\end{figure}
\FloatBarrier

The right-hand side is exactly the same ``correlation with noise'' structure---just for functions $\mu - m$ in the model $\model$ rather than functions $(1 - W\gammaipw)m$ in the IPW-reweighted model $\model^{ipw} = \{(1 - W\gammaipw)m : m \in \model\}$ that appear in the same role in our symmetrized imbalance~\eqref{eq:ipw_gen_imbalance_symmetrized}. Models that are easy to estimate by least squares are easy to balance. And the notions of model size used in the context of least squares---Rademacher complexity, entropy integrals, effective dimension, the Donsker property---are the right notions for talking about balance because they are approximations to or summaries of the maximal correlation~\eqref{eq:ipw_gen_imbalance_symmetrized}.

\paragraph{Why $O_p(n^{-1/2})$ is not enough.}
Unfortunately, even for Donsker models, $O_p(n^{-1/2})$ imbalance is not sufficient for valid inference---we need $o_p(n^{-1/2})$. To see this, consider the extreme case of a model containing a single function: $\model = \{\mu_1\}$. Then
\begin{equation}
  \label{eq:single_function_imbalance}
  \imbalance_\model(\gammaipw) = \left|\frac{1}{n}\sum_{i=1}^n \{1 - W_i \gammaipw(X_i)\} \mu_1(X_i)\right|.
\end{equation}
By the central limit theorem, $\sqrt{n}\, \imbalance_\model(\gammaipw)$ converges not to zero but to the absolute value of a normal with variance $\var[\{1 - W_i \gammaipw(X_i)\} \mu_1(X_i)]$.\footnote{The IPW estimator with the true weights $\gammaipw$ is asymptotically unbiased and normal, because the imbalance term is perfectly centered: we get additional variance, but not bias. However, this delicate centering property is not preserved by estimated weights, even ones close to $\gammaipw$ in mean squared error.} The imbalance is $O_p(n^{-1/2})$, not $o_p(n^{-1/2})$, even when the model contains only the truth.

We see the same phenomenon in least squares. The global complexity bound
\begin{equation}
\label{eq:least-squares-fixed-point-naive}
\norm{\hat \mu - \mu} \neq r \quad \text{if} \quad
\frac{r^2}{2} \ge  \max_{m \in \model} \frac1n\sum_{i=1}^n \varepsilon_i \{ \mu(X_i) - m(X_i) \}
\end{equation}
gives $n^{-1/4}$ rates at best---even for a model containing a single function. The problem is that we are bounding the noise correlation over all of $\model$, not just over the small functions near $\mu$ that we actually care about. To do better, we need \emph{local} complexity measures like the one in~\eqref{eq:least-squares-fixed-point}, which restrict the maximum to functions at distance $r$ from the truth. The decay in the maximum as $r$ shrinks is called the \emph{equicontinuity} of the empirical process.

And this is exactly what we need for the augmented estimator. Because we care about imbalance in $\dm = \hat \mu_1 - \mu_1$, which shrinks as our regression estimate improves, equicontinuity---the fact that we balance small functions better than big ones---is what converts $O_p(n^{-1/2})$ imbalance over the full model into $o_p(n^{-1/2})$ imbalance in the regression error.

\subsection{The role of augmentation}
\label{sec:role-of-augmentation}
% While the ``plugin'' approach is insufficient to show that Condition \ref{cond:small_imbalance} holds with the weights alone,
We can do better with this approach when analyzing the augmented estimator we discuss in Section \ref{sec:aipw}.
The key idea is that rather than analyzing balance in a \emph{fixed} model, by augmenting the estimator we can restrict our attention to a \emph{shrinking} sequence of models.
%We can show that the imbalance achieved by the inverse propensity weights on this sequence of models is indeed negligible.

For intuition, we will again consider a model containing a single element. This time, however, we will use the model $\calM_n = \{\delta m_n\}$ that contains the regression error $\dm_n = \hat \mu_1 - \mu_1$ instead of $\mu_1$ itself. What has changed is that, assuming we use a good estimator $\hat \mu_1$, the element $\dm_n$ in this set will shrink as we obtain more observations.
It is convenient to think of $\hat \mu_1$ as cross-fit, i.e., estimated
on a sample other than $\{ (W_i,X_i,Y_i) : i \le n\}$. This allows us to think of $\dm$  as \emph{non-random} in the sense that it does not depend on this sample; formally, we do most of our analysis conditional on $\dm$.\footnote{This is not necessary for all arguments, e.g., it is not required by those of \citet{hirshberg2021augmented}.
However, it is increasingly common practice, has been shown to be help in some simulation studies, and
can be done without loss of efficiency asymptotically by averaging estimators that fit $\hat \mu_1$ on different subsamples
\citep[see, e.g.][]{chernozhukov2018double, schick1986asymptotically, zheng2011cross}.}
The maximal imbalance is
\[
  \imbalance_{\calM_n}(\gammaipw) = \left|\frac{1}{n}\sum_{i=1}^n W_i \gammaipw(X_i) \delta m_n(X_i)-\delta m_n(X_i)\right| \quad \text{ for } \quad \model_n = \set{ \delta m_n}.
\]
By the same argument we used for $\model=\set{\mu_1}$, if $\Var[\{1-W_i\gammaipw(X_i)\}\delta m_n(X)]$ converges to zero, then $\sqrt{n} \imbalance_{\calM_n}(\gammaipw)$ will converge to zero in probability and the estimator will achieve asymptotically negligible imbalance.
% in Condition \ref{cond:small_imbalance} will be satisfied.
Given mild assumptions, this quantity converges to zero if $\hat \mu_1$
is a consistent estimator.\footnotemark

\footnotetext{For example, sufficient overlap (Assumption~\ref{a:regular-overlap}) and boundedness of $\norm{\dm_n}_{\infty}$ suffice \citep[see Proof of Theorem 1]{hirshberg2021augmented}.}

To view this phenomenon more generally, suppose we have a valid model $\model_n$ with many elements, all of them shrinking like $\dm_n$ in the previous example. In particular, suppose that we have an `a priori model' $\model_0$ specifying sample-size-independent characteristics we expect of $\dm$, e.g. smoothness or approximate sparsity,
and that we know the rate of convergence of our estimator $\hat \mu_1$,
so at any given sample size we are able to rule out the possibility that our regression error $\norm{\dm_n}_{L_2(P)}$ exceeds some bound $r_n$
that shrinks to zero.\footnote{In discussion of the role of augmentation within the balancing tradition,
it has been common to focus on consistency in the strong sense that $\norm{\dm}_{\model} \to 0$ rather than in the mean-square sense $\norm{\dm}_{L_2(P)} \to 0$ that we use here. While simpler, that approach has some disadvantages. See Appendix~\ref{sec:strong-norm-consistency} for a discussion.} Incorporating this knowledge leaves us with the smaller model $\model_n = \set{m \in \model_0 : \norm{m}_{L_2(P)} \le r_n}$.
It is reasonable to ask whether this is enough that the maximal imbalance would still satisfy $\sqrt{n}\imbalance_{\model_n}(\gammaipw) \to 0$, which would imply negligible imbalance in
$\dm_n$ if our assumptions were valid in the sense that $\dm_n \in \model_0$ and $\norm{\dm_n}_{L_2(P)} \le r_n$ with probability tending to one as $n$ grows.
The answer depends on the size the a-priori model $\model_0$, but essentially a sufficient condition on $\model_0$ is that $\imbalance_{\model_0}(\gammaipw) = O_p(n^{-1/2})$.\footnotemark\

\footnotetext{To be precise, it is yes when $\calM_0$ and $\{f(w,x) = w\gammaipw(x)m(x) : m \in \calM_0\}$ are Donsker
and the maximal variance of a term for an individual element of $\model_n$, $\max_{\dm \in \model_n}\Var[(1-\gammaipw(X_i)W_i)\dm(X)]$,
shrinks to zero. Generalizing the one-element case discussed above,
sufficient overlap (Assumption~\ref{a:regular-overlap}) and boundedness of $\max_{m \in \model_0}\norm{m}_{\infty}$ suffice for the latter.
A slight generalization of this case is addressed in \citep[see Proof of Theorem 1]{hirshberg2021augmented}.
This phenomenon, the inheritance of convergence of $\sqrt{n}\imbalance_{\calM_n}(\gammaipw)$
from the convergence of the largest variance of a term $(1-\gammaipw(X_i)W_i)\dm(X)$ in the averages
we are maximizing over, is called the asymptotic equicontinuity of the empirical process,
and is the defining feature of Donsker classes \citep[e.g.,][Section 2.1.2]{van1996weak}.}

From this, it would appear that we need to know a bound on the regression error $\norm{\dm_n}_{L_2(P)}$ so that we can use the model $\model_n$ to solve for our balancing weights. However, \citet[see Section 2]{hirshberg2021augmented} show that this is unnecessary, as it works well enough to simply use the model $\model_0$ instead. In particular, they show that the minimax weights \eqref{eq_balancing-weights} for an a-priori model $\model_0$ of appropriate size achieve negligible imbalance over any sequence of submodels $\model_n = \{ \dm \in \model_0 : \norm{\dm}_{L_2(P)} \le r_n\}$ satisfying  regression error bounds $r_n$ tending to zero. Because this will include one for which $r_n \ge \norm{\dm_n}_{L_2(P)}$ with probability tending to one if one exists (i.e., if $\hat \mu_1$ is consistent), in practice all we need is an a-priori model $\model_0$ that is valid in the sense that
 $\dm \in \model_0$ with probability tending to one.
Concretely, \citet[Theorem 1]{hirshberg2021augmented} show that when we use minimax balancing weights for a Donsker a-priori model $\model_0$ and augment with a consistent estimator $\hat \mu_1$, the AIPW estimator is asymptotically linear with the efficient influence function---requiring only overlap and that $\dm \in \model_0$, with no assumptions on the propensity score.

This is an instance of the equicontinuity phenomenon we discussed above: imbalance over a shrinking sequence of models inherits the $o_p(n^{-1/2})$ rate from shrinking variance, even when the maximal imbalance over the fixed a-priori model is only $O_p(n^{-1/2})$.

One subtlety is that using a model $\model'$ for the outcome $\mu_1$ does not guarantee that $\dm = \mu_1 - \hat \mu_1$ is in $\model'$. For example, if $\model'$ consists of smooth functions and $\hat \mu_1$ is a regression tree, the difference $\dm$ may not be smooth even when $\mu_1$ is. This mismatch can cause augmentation to fail; see Section 3.2.3 of the arXiv v1 version of \citet{hirshberg2021augmented} for a simulation. A simple fix is to take $\model_0$ to be the convex hull of $\model' - \hat \mu_1 := \{m - \hat \mu_1 : m \in \model'\}$, which contains $\dm$ whenever $\mu_1 \in \model'$ \citep[Section 2.6]{hirshberg2021augmented}.

%For this model class, we estimate the model via $\hat \mu_1(x, 1) = \hat{\beta} \cdot \phi(x)$. By H\"{o}lder's inequality, we can bound the bias term in Equation\ref{eq:aug_bias} by
%\begin{equation}
%    \label{eq:aug_bias_linear}
%   \frac{1}{n}\sum_{i=1}^nW_i\hat{\gamma}_i (\hat{\beta} \cdot \phi(X_i) - \beta \cdot \phi(X_i)) - \frac{1}{n}\sum_{i=1}^n (\hat{\beta} \cdot \phi(X_i) - \beta \cdot \phi(X_i)) \leq  \|\hat{\beta} - \beta\|_1 \imbalance_{\Phi^1}(\hat{\gamma}).
%\end{equation}
%As we have seen in Section \ref{sec:plugin}, via the ``plugin'' approach we can show that the imbalance term is $O_p\left(\sqrt{\frac{\log p}{n}}\right)$. Therefore, as long as the coefficients are consistent in the one-norm, $\|\hat{\beta} - \beta\|_1 = o_p(1)$, we can be sure that the augmented estimator achieves the negligible imbalance Condition \ref{cond:small_imbalance}. This approach and variants of it have been used by \citet{athey2018approximate} for the sparse model class $\Phi^1$, \citet{wong2018kernel} and \citet{kallus2020generalized} for RKHSs, and \citet{benmichael2021_multical} for discrete covariates with interactions.



%To view this phenomenon more generally, recall the bias due to imbalance for the augmented estimator in Equation\ref{eq:weights_error_aug}
%\begin{equation}
%    \label{eq:aug_bias}
%    \frac{1}{n}\sum_{i=1}^nW_i\hat{\gamma}_i (\hat \mu_1(X_i, 1) - \mu_1(X_i)) - \frac{1}{n}\sum_{i=1}^n (\hat \mu_1(X_i, 1) - \mu_1(X_i)).
%\end{equation}
%Here the bias stems from imbalance in the error term $\delta(X) := \hat \mu_1(X, 1) - \mu_1(X)$ rather than imbalance in the true model $\mu_1$. We now consider several approaches that instead analyze imbalance in the errors.



%


%\paragraph{Combining regression error and imbalance.}
%One approach combines the behavior of the regression error $\hat \mu_1(\cdot,1) - \mu_1$ with the balancing behavior that $\imbalance_{\calM}(\gammaipw) = O_p(\frac{1}{\sqrt{n}})$, to show that the bias will be asymptotically negligible. As an example of this approach, again consider the set of approximately sparse functions, $\Phi^1 = \{\beta \cdot \phi(x) : \|\beta\|_1 \leq 1\}$. For this model class, we estimate the model via $\hat \mu_1(x, 1) = \hat{\beta} \cdot \phi(x)$. By H\"{o}lder's inequality, we can bound the bias term in Equation\ref{eq:aug_bias} by
%\begin{equation}
%    \label{eq:aug_bias_linear}
%   \frac{1}{n}\sum_{i=1}^nW_i\hat{\gamma}_i (\hat{\beta} \cdot \phi(X_i) - \beta \cdot \phi(X_i)) - \frac{1}{n}\sum_{i=1}^n (\hat{\beta} \cdot \phi(X_i) - \beta \cdot \phi(X_i)) \leq  \|\hat{\beta} - \beta\|_1 \imbalance_{\Phi^1}(\hat{\gamma}).
%\end{equation}
%As we have seen in Section \ref{sec:plugin}, via the ``plugin'' approach we can show that the imbalance term is $O_p\left(\sqrt{\frac{\log p}{n}}\right)$. Therefore, as long as the coefficients are consistent in the one-norm, $\|\hat{\beta} - \beta\|_1 = o_p(1)$, we can be sure that the augmented estimator achieves the negligible imbalance Condition \ref{cond:small_imbalance}. This approach and variants of it have been used by \citet{athey2018approximate} for the sparse model class $\Phi^1$, \citet{wong2018kernel} and \citet{kallus2020generalized} for RKHSs, and \citet{benmichael2021_multical} for discrete covariates with interactions.



%\paragraph{Balancing the regression error.}
%\avi{I found this difficult to follow. Just leaving it for now.}
%An alternative approach proposed by \citet{hirshberg2021augmented} instead finds weights that control the worst-case imbalance in a class of models for the regression \emph{error} $\dm(\cdot) \in \Delta$ rather than the model $\mu_1 \in \calM$. Denoting the rate of convergence in mean square error as $r_n$ --- $\E[(\hat \mu_1(X,1) - \mu_1(X))^2] \leq r_n^2$ --- then the error function $\delta(\cdot)$ is in the class of functions with variance bounded by $r_n^2$: $r_n L_2(P) := \{f : E[f(X)^2] \leq r_n^2\}$, as well as being in $\Delta$.
%Therefore, we consider imbalance in the intersection of these two classes $\Delta \cap r_nL_2(P)$.


%Importantly, this new augmented model class is substantially smaller than the original model class $\calM$. As long as $\hat \mu_1(\cdot,1)$ is a consistent estimator for $\mu_1$ and so $r_n \to 0$ as $n \to \infty$, then the augmented model class $\calM \cap r_n L_2(P)$ will shrink to be zero. In the previous section we have seen that in many cases we can establish that the worst-case imbalance over a fixed model class $\calF$ for the true propensity score is $\imbalance_{\calF}(\gammaipw) = O_p(n^{-1/2})$. Now by augmenting with a consistent estimator we can use the consistency to guarantee asymptotically negligible imbalance in the shrinking model class $\Delta \cap r_n L_2(P)$: $\imbalance_{\Delta \cap r_n L_2(P)}(\gammaipw) = o_p(n^{-1/2})$.
%Now the ``plugin'' argument in Section \ref{sec:plugin} is sufficient to show negligible imbalance for the augmented estimator.

%While we have considered finding weights that balance the worst-case regression error $\delta(\cdot) \in \Delta$, \citet{hirshberg2021augmented} show that it is in fact sufficient to find weights that control the worst-case imbalance in the model class $\calM$. In essence, balancing the model class $\calM$ is nearly optimal for the augmented model class $\Delta \cap r_n L_2(P)$ as well.
%Therefore, by finding weights that balance $\calM$ and augmenting with a consistent estimate of $\mu_1$, we can ensure asymptotically negligible imbalance under a wide array of model classes.

\subsection{Doing Better than the True Weights}%
\label{sec:doing-better}

The approaches above---plug-in comparison and augmentation---ultimately characterize when the AIPW estimator achieves behavior comparable to what we would get with the true inverse propensity weights $\gammaipw$. But estimated weights can do better. Even when one is using the true weights, as in a randomized experiment, neither IPW nor AIPW has any bias, but the AIPW estimator has smaller variance:
\begin{equation}
\label{eq:aipw-variance-true-weights}
\Var\!\left(\hat\psi_1^{aipw}\right) = \Var\!\left(\frac{1}{n}\sum_i W_i \gammaipw(X_i) \varepsilon_i\right) + \E\left[\left(\frac{1}{n}\sum_i W_i \gammaipw(X_i) \dm(X_i) - \frac{1}{n}\sum_i \dm(X_i)\right)^{\!2}\right].
\end{equation}
The first term is irreducible noise. The second is the mean-squared imbalance in $\dm$. IPW is the special case $\hat \mu_1 = 0$, so $\dm = \mu_1$. When $\hat \mu_1$ is a good estimator, $\dm$ is closer to zero than $\mu_1$ is, so AIPW has lower variance. This is the cost of imbalance of order $n^{-1/2}$: it inflates variance even when it causes no bias.

There is another path to low variance: estimated weights that achieve smaller imbalance than the true weights would. This may seem paradoxical---how can an estimate beat the truth?---but there is a simple explanation. With a $k$-dimensional model and $k \le n$, we can achieve exact balance: it is just linear algebra. The true propensity score does not know it needs to balance our particular basis, so it will typically have imbalance on the order of $n^{-1/2}$.

\citet{chan2016globally} formalize this for sieves, showing that calibration weights achieving exact balance on a growing basis attain semiparametric efficiency \citep[see also][]{fan2016improving, wang2020minimal}. For RKHS models, \citet{hirshberg2019minimax}, inspired by \citet{newey2018cross}, show that minimax balancing weights achieve $\imbalance_{\model}(\hgamma) = o_p(n^{-1/2})$ essentially without assumptions on $\gammaipw$. When the RKHS model is correctly specified, we do not need augmentation for asymptotic efficiency. One does not even need to be explicitly balancing to get this benefit. \citet{hirano2003efficient} showed that the IPW estimator using an estimated propensity score is more efficient than one using the true propensity score. The intuition is that the estimated propensity score is fit to the sample, so it achieves better balance on observed covariates than the population propensity score would.

The results above establish that balancing weights can beat the true weights, but they do not tell us by how much. These results are part of a long tradition of work on minimax estimation of linear functionals \citep{ibragimov1985nonparametric, donoho1991geometrizing, donoho1994statistical, armstrong2018optimal} that provides the tools to answer this. That literature studies problems like ours: estimating a linear functional $\psi_1 = \E[\mu_1(X)]$ when $\mu_1$ is known to lie in a convex function class. Its central object is the \emph{modulus of continuity},\footnote{This penalized formulation is equivalent by Lagrangian duality to the constrained formulation used by \citet{donoho1994statistical} and \citet{armstrong2018optimal}. See Appendix~\ref{app:modulus-equivalence}.}
\begin{equation}
\label{eq:modulus-of-continuity}
\omega(s) = \sup_{m \in \model} \left\{ \frac{1}{n}\sum_{i=1}^n m(X_i) - \frac{1}{2sn} \sum_{i=1}^n m(X_i)^2 \right\},
\end{equation}
which characterizes the worst-case bias achievable at a given variance level. The maximal imbalance of the minimax balancing weights has a clean characterization in terms of $\omega$. At the optimal bias-variance tradeoff parameterized by $s > 0$,
\begin{equation}
\label{eq:donoho-liu-bias}
\imbalance_\model(\hgamma) = \omega(s) - \omega'(s)\,s.
\end{equation}
When $\omega$ is approximately linear---$\omega(s) \approx c\,s$ for some constant $c$---the maximal imbalance is approximately zero. So bounding imbalance reduces to bounding the deviation of $\omega$ from linearity.

The deviation of $\omega$ from linearity is bounded by the local complexity of the model---the Gaussian width of $\model$ intersected with an $L_2$ ball of radius $r$, the same quantity that drove the equicontinuity argument in Section~\ref{sec:role-of-augmentation}. There, localization converted $O_p(n^{-1/2})$ imbalance over the full model into $o_p(n^{-1/2})$ imbalance in a shrinking neighborhood. Here, the same localization bounds how far the modulus can deviate from linearity. For any $r > 0$ satisfying the fixed-point inequality~\eqref{eq:least-squares-fixed-point}---the same inequality that controls the least squares error---the deviation is bounded by a constant times $r^2$ \citep{kong2025asymptotics, hirshberg2025bregman}, giving
\begin{equation}
\label{eq:balance-fixed-point}
\imbalance_\model(\hgamma) \le C\, r^2.
\end{equation}
This holds provided $\gammaipw$ can be approximated in the model at rates compatible with $r$: there must exist a function $g$ with $\norm{g}_\model \le C\, r^2 n / \sigma^2$ and $\norm{g - \gammaipw}_{L_2(P)} \le C\, r^2 \sqrt{n}/\sigma$, where $\sigma$ is the tuning parameter in~\eqref{eq_balancing-weights}. A smaller $\sigma$ or larger $n$ makes the penalty cheaper, allowing a richer comparison function. When $\model$ is nonparametric---its gauge-bounded functions are dense in $L_2(P)$---such approximations exist for a sequence $r \to 0$, so $\imbalance_\model(\hgamma) = o_p(n^{-1/2})$ without any assumption on $\gammaipw$ itself. We record this.

\begin{proposition}[Maximal imbalance of minimax balancing weights]
\label{prop:bias-balance}
Let $\hgamma$ be the minimax balancing weights~\eqref{eq_balancing-weights} with tuning parameter $\sigma$,
and let $r > 0$ satisfy the fixed-point inequality~\eqref{eq:least-squares-fixed-point}. If there exists a function $g$ with $\norm{g}_\model \le C\, r^2 n / \sigma^2$ and $\norm{g - \gammaipw}_{L_2(P)} \le C\, r^2 \sqrt{n}/\sigma$, then $\imbalance_\model(\hgamma) \le C'\, r^2$ for universal constants $C, C'$. If $\model$ is nonparametric, such approximations exist for $r \to 0$, so $\imbalance_\model(\hgamma) = o_p(n^{-1/2})$.
\end{proposition}

The measures of model complexity that appear in nonparametric regression---Rademacher complexity, entropy integrals, eigenvalue decay---are not merely analogous to the right measures for balance. They \emph{are} the right measures. The fixed-point radius $r$ that bounds $\norm{\hat\mu - \mu}$ also bounds the maximal imbalance, through the same mechanism: localization of an empirical process. This is the same quantity illustrated in Figure~\ref{fig:width-geometry}, where the crossing point determines $r$.

This route to bounding imbalance is fundamentally different from the comparison to $\gammaipw$ in Section~\ref{sec:plugin}. That approach---bounding $\imbalance_\model(\hgamma)$ via~\eqref{eq_balancing-weights-comparison}---requires $\hgamma$ to converge to $\gammaipw$, limiting how aggressively the tuning parameter can prioritize balance over variance. The modulus argument bounds imbalance directly from the geometry of the model, without decomposing it into ``true weights plus estimation error.'' Because convergence of the weights does not enter the argument, there is no restriction on the degree of undersmoothing. Less smoothing---a smaller $\sigma$---means better balance, because a smaller $\sigma$ makes the penalty cost $(\sigma^2/n)\norm{g}_\model^2$ cheaper, admitting richer comparison functions that approximate $\gammaipw$ more closely.

\FloatBarrier
\begin{figure}[h]
\centering
\includegraphics[width=0.48\textwidth]{figure/insample-vs-outsample-balance.pdf}%
\hfill%
\includegraphics[width=0.48\textwidth]{figure/insample-vs-outsample-balance-b.pdf}
\caption{\textit{Left}: In-sample versus out-of-sample imbalance as a function of the tuning parameter $\sigma^2$, for $n \in \{100, 200, 400\}$ (bold: $n = 200$). Balancing weights are fit on sample~A and evaluated on both sample~A (in-sample, blue) and an independent sample~B (out-of-sample, red). Gray lines mark IPW imbalance. Shaded regions indicate three regimes: undersmoothing, bias--variance tradeoff, and constant weights. The gap is the in-sample advantage exploited by Proposition~\ref{prop:bias-balance}. \textit{Right}: Regime widths as a function of $n$. Each stick spans from the crossing $\sigma^2$ (where in-sample and out-of-sample curves cross) to the near-plateau $\sigma^2$, across three propensity shapes. The bias--variance regime widens systematically with~$n$.}
\label{fig:insample-outsample}
\end{figure}

But this is an \emph{in-sample} phenomenon. If we fit weights on one sample and evaluate balance on another, we lose the cancellation that makes aggressive tuning work. Figure~\ref{fig:insample-outsample} illustrates this: in-sample imbalance decreases monotonically as $\sigma$ shrinks, while out-of-sample imbalance first decreases then increases as the weights overfit the training sample. The gap between the two curves is precisely the in-sample advantage of balancing weights---the advantage that cross-fitting, which evaluates weights on held-out data, cannot access.

\begin{remark}[Concrete rates in RKHS models]
\label{remark:rkhs-rates}
For an RKHS $\calH$ with eigenvalues $\lambda_j$ and eigenfunctions $\phi_j$, the approximation condition in Proposition~\ref{prop:bias-balance} depends on how smooth $\gammaipw$ is relative to the RKHS. Write $\gammaipw = \sum_j c_j \phi_j$ in the eigenbasis of the kernel integral operator. If $\sum_j c_j^2 / \lambda_j^\kappa < \infty$ for some $\kappa > 0$---a regularity condition requiring the Fourier coefficients of $\gammaipw$ to decay fast enough relative to the eigenvalues---then the bias satisfies
\[
  \imbalance_\model(\hgamma) \lesssim \left(\frac{\sigma}{\sqrt{n}}\right)^{1+\kappa}
\]
\citep[Appendix~B]{hirshberg2019minimax}. This is $o(\sigma/\sqrt{n})$ for any $\kappa > 0$, so the approximation condition holds automatically. The exponent $\kappa$ measures regularity beyond membership in the RKHS: $\kappa = 1$ means $\gammaipw \in \calH$, while $\kappa < 1$ allows rougher functions that are not in the RKHS itself but can still be approximated at a polynomial rate.
\end{remark}

\begin{remark}[Augmentation and convex models]
\label{remark:augmentation-convex}
Where does this leave augmentation? For convex models, the minimax linear theory tells us that augmentation is unnecessary for near-optimality, and the results above show this extends to asymptotic efficiency with no restrictions on undersmoothing. The largest benefits of augmentation appear in nonconvex models, such as sparse models, where linear estimators are no longer near-optimal \citep{cai2004minimax}. In simulations using sparse models, augmentation cuts root mean squared error by half or more \citep[Section~6, Setup~2]{hirshberg2021augmented}.
\end{remark}

\subsection{The inner product perspective}%
\label{sec:inner-product}

The preceding subsections analyze bias through maximal imbalance---a worst-case measure over the model. Here we take a complementary perspective, working with the following decomposition into an inner product of errors and the imbalance in the regression error achieved by the true weights:
\begin{equation}
    \label{eq:aug_bias_dr}
    \frac{1}{n}\sum_{i=1}^nW_i \delta\gamma(X_i) \dm(X_i) \;+\; \frac{1}{n}\sum_{i=1}^n \left\{W_i \gammaipw(X_i) - 1\right\}\dm(X_i) \quad \text{ where } \quad  \delta\gamma = \hat{\gamma}  - \gammaipw.
\end{equation}
The second term in Equation~\ref{eq:aug_bias_dr} is the imbalance in the error function $\dm$ with the true
inverse propensity weights $\gammaipw$; as we discuss above, this will be $o_p(n^{-1/2})$ if $\hat \mu_1$ is consistent.
So we can focus on the first term, the inner product of the errors $\dgamma$ and $\dm$. The conditional bias is negligible \emph{if and only if} this inner product is. This can happen for two reasons: the factors can be small in magnitude, or they can be nearly orthogonal \citep[Section 1.4]{hirshberg2021augmented}.

The magnitude path uses Cauchy--Schwarz to bound the inner product by the product of $L_2$ norms. If this product, $\norm{\delta\gamma}_{L_2(\Pn)}\norm{\dm}_{L_2(\Pn)}$, is $o_p(n^{-1/2})$, the conditional bias is negligible. This is \emph{rate double robustness}: we need not estimate either $\gamma$ or $\mu_1$ well, so long as we estimate one well enough to compensate for the other. But the magnitude path has limitations. If both nuisances converge at the same rate, that rate must be faster than $n^{-1/4}$ for their product to be $o_p(n^{-1/2})$.

The balancing arguments in Sections~\ref{sec:plugin} and~\ref{sec:role-of-augmentation} take a different path. Rather than bounding the magnitudes of $\dgamma$ and $\dm$ separately, they bound the inner product directly. When $\hgamma$ balances a model containing $\dm$, the imbalance bound gives $\langle \dgamma, \dm \rangle_{L_2(\Pn)} = o_p(n^{-1/2})$. Because an inner product can only be small if the factors are small or nearly orthogonal, this tells us that if neither nuisance estimate is converging fast, they must be approximately orthogonal.

The inner product perspective has an advantage that a uniform, worst-case approach to bounding imbalance does not: it can exploit instance-specific structure in how $\dgamma$ and $\dm$ interact. When weights or the outcome model are estimated by least squares onto a basis, regularized regression in an RKHS, or the lasso, the first-order conditions make estimation errors orthogonal to the fitted subspace. If the other nuisance's error lies in or near that subspace, this sharpens the Cauchy--Schwarz bound \citep{vermeulen2015bias,tan2010bounded,hirshberg2019minimax,bradic2019sparsity}. This is the mechanism underlying the RKHS results of \citet{hirshberg2019minimax} discussed in Section~\ref{sec:doing-better}.

\begin{remark}[Optimal double robustness]
\label{remark:optimal-dr}
This sharper analysis is what makes \emph{optimal double robustness} possible: efficient inference whenever \emph{either} $\gammaipw$ or $\mu_1$ is smooth enough to be estimable at faster than $n^{-1/4}$ rate---a requirement that cannot be improved \citep{mukherjee2017semiparametric, robins2009semiparametric}. While the estimators known to achieve this are relatively complex, \citet{newey2018cross} show that AIPW with specialized nuisance estimators can come close, and recent work extends this to balancing weights \citep{hirshberg2019minimax, kennedy2020optimal}.
\end{remark}

\begin{remark}[Regression for orthogonality]
\label{remark:regression-for-orthogonality}
\citet{wang2020debiased} observe that this logic can be reversed: rather than choosing weights orthogonal to outcome model errors, we can choose the outcome regression so that its error is orthogonal to weight errors. This yields $\sqrt{n}$-consistent estimation under assumptions on the propensity score rather than the outcome model. See Appendix~\ref{app:regression-orthogonality} for a sketch in our notation.
\end{remark}


% \vspace{-.5cm}
\section{Other estimands}
\label{sec:other-estimands}

Thus far, we have considered estimating the average treatment effect $\tau = \psi_1 - \psi_0$, focusing on estimating $\psi_1$ in particular.
These ideas and this framework, however, can be used for many other quantities of interest.
In particular, they generalize
% This is not limited to estimating the treatment specific mean $\psi_1$. It can be generalized
to estimate
many other functions $\estimand(m)$ of the conditional mean outcome $\mu_w(x) = \E[Y \mid W=w,\ X=x]$, where the treatment $w$ need not be discrete-valued. This includes many generalizations of the average treatment effect. For example,
\begin{align}
%&\estimand(m) = \E[ \mu_1(X) - \mu_0(X)]                 &&\text{average treatment effect}; \\
&\estimand(m) = \E\sqb*{\frac{\partial}{\partial w} m_{w} (X) \mid_{w=W} }, &&\text{average effect of an infinitesimal dose increase}; \\
&\estimand(m) = \int \{\mu_1(x) - \mu_0(x)\} p(x)dx, &&\text{population-targeted average treatment effect}. \label{eq:targeted-ate}
\end{align}
All of these are instances of a  general estimand that can be estimated with essentially the same approach
\citep[see, e.g,][]{chernozhukov2018biased, hirshberg2021augmented}.
\begin{equation}
\label{eq:amle-general}
\estimand(m) = \E[ h(W_i, X_i, m) ]  \quad \text{ where } \quad h(x,w,m) \ \text{ is linear in }\ m.
\end{equation}
Above, where we have focused on the case that $\estimand(m) = \psi_1 = \E[\mu_1(X_i)]$, the inverse propensity weights $\gamma_\estimand(w,x)=w\gammaipw(x)$ have played an essential role. In general, this role is played by the \emph{Riesz representer} of the function $\estimand$, which is the solution $\gamma_{\estimand}(w, x)$ to a balance condition that generalizes the one defining the inverse propensity weights \eqref{eq_mipw-balance}:
\begin{equation}
\label{eq:riesz-representer-general}
\E[h(W_i,X_i, f)] = \E[ \gamma_{\estimand}(W_i, X_i)f(W_i,X_i)] \ \text{ for all }\  f \ \text{ with }\ \E[ f(W_i, X_i)^2 ] < \infty.
\end{equation}
The solution $\gamma_\estimand$ is guaranteed to exist and be square integrable when $\estimand$ is \emph{mean-square-continuous}. For $\estimand(m)=\psi_1$, mean-square continuity is %,as discussed in Section~\ref{sec:framework},
 equivalent to sufficient overlap.


We can mechanically derive generalized forms of everything we discuss above: the IPW and AIPW estimators, imbalance-focused decompositions of their error, appropriate forms of maximal imbalance, and estimators for balancing weights.
To do this, we substitute
$\gamma_{\estimand}(W_i,X_i)$ for $W_i\gammaipw(X_i)$ and then
$h(W_i,X_i,m)$ for $\mu_1(X_i)$ in Equations~\ref{eq_general_ipw}-\ref{eq_balancing-weights-general} above.%
%
% \avi{there is a lot in this sentence. can you break it apart?}
% Generalized forms of the IPW and AIPW estimators, imbalance-focused decompositions of their error, and appropriate forms of maximal imbalance
% and
% therefore balancing weights
% can be derived by substituting
% $\gamma_{\estimand}(W_i,X_i)$ for $W_i\gammaipw(X_i)$ and then
% $h(W_i,X_i,m)$ for $\mu_1(X_i)$ in Equations~\ref{eq_general_ipw}-\ref{eq_balancing-weights-general} above.
\footnote{Order is important when making these substitutions only because we have followed convention in writing $W_i \mu_1(X_i)$ where the equivalent expression $W_i \mu_{W_i}(X_i)$ is perhaps more natural.}
All of our discussion above generalizes, whether focused more on maximal imbalance
as in Sections~\ref{sec:balancing-binary}-\ref{sec:doing-better}
\citep{hirshberg2021augmented} or on estimation of the inverse propensity score as in Section~\ref{sec:inner-product}  \citep{chernozhukov2018biased}.

For the resulting IPW and AIPW estimators $\hat\psi^{\text{ipw}} = \frac{1}{n}\sum_{i=1}^n  \hat\gamma_i(W_i,X_i) Y_i$ and $\hat\psi^{\text{aug}} = \frac{1}{n}\sum_{i=1}^n h(W_i,X_i,\hat \mu) + \hat\gamma(W_i,X_i) \{ Y_i - \hat \mu(W_i,X_i) \}$, the relevant notion of imbalance is
\begin{equation}
    \label{eq:generalizations}
    \begin{aligned}
    \imbalance_{h,\calM}(\hgamma) &= \max_{m \in \calM} \abs*{\frac{1}{n}\sum_{i=1}^n \hgamma(W_i, X_i) m(W_i, X_i) - \frac{1}{n}\sum_{i=1}^n h(W_i,X_i, m)}.
    \end{aligned}
\end{equation}
\cite{hirshberg2021augmented} describes the behavior of $\hat\psi^{\text{aug}}$ with weights generalizing the minimax weights \eqref{eq_balancing-weights} described above. The proposition below, which generalizes Proposition~\ref{proposition:clt}, describes the behavior of the estimators $\hat\psi^{\text{ipw}}$ and $\hat\psi^{\text{aipw}}$ with arbitrary design-based weights.
\begin{proposition}
\label{proposition:clt-general}
Let $\Gamma :=  \sqrt{n^{-1}\sum_{i=1}^n \hgamma(W_i, X_i)^2}$ be the root-mean-squared weight. Suppose that each squared weight is negligible compared to their sum $n\Gamma^2$, and assume that the noise sequence $\varepsilon_1 \ldots \varepsilon_n$ is uniformly square-integrable conditional on the design. If the imbalance in the regression error $\delta m = m - \hat \mu$ is $o_p(n^{-1/2} \Gamma)$,
 \begin{equation}
 \label{eq:asymptotic-linearity-amle}
 \hat\psi^{\aug} -  \tilde\psi = \frac{1}{n}\sum_{i=1}^n\hgamma(W_i,X_i) \{ Y_i - m(W_i,X_i) \} + o_p(n^{-1/2}\Gamma) \quad \text{ for } \quad \tilde\psi = \frac{1}{n}\sum_{i=1}^n h(W_i,X_i, m).
 \end{equation}
 Furthermore, if $\hat \mu$ is a consistent estimator of $\mu$, then
 \begin{equation}
 \label{eq:asymptotic-normality-amle}
  \frac{\hat\psi^{\aug} -  \tilde \psi}{\sqrt{\hat V}} \to N(0,1) \quad \text{for} \quad
 \hat V = \frac{1}{n^2}\sum_{i=1}^n \hgamma(W_i, X_i)^2 \{Y_i - \hat \mu(W_i, X_i)\}^2.
 \end{equation}
If imbalance in $m$ itself is $o_p(n^{-1/2} \Gamma)$, the same results hold with $\hat \psi^{\text{ipw}}$ in place of $\hat\psi^{\text{aug}}$.
\end{proposition}
Some common estimands, such as the average treatment effect on the treated (ATT), are not directly included in the general class defined by \eqref{eq:amle-general}, but can be expressed as the quotient of something that is in this class and something that is trivial to estimate. For example, the ATT can be written as $\E[m(1,X) - m(0,X) \mid W=1] = \E[W \{ m(1,X) - m(0,X) \}] / \E[W]$. To estimate it, we can start with an estimator  $\hat\psi^{ipw}$ or $\hat\psi^{aug}$ for the numerator, then divide by the natural estimator of the denominator, $\bar W = \frac{1}{n}\sum_{i=1}^n W_i$, so division by $n$ in the formula for $\hat\psi^{ipw}$ or $\hat\psi^{aug}$ is replaced with division by $N_1 = n\bar W = \sum_{i=1}^n W_i$.

Furthermore, a small tweak allows us to estimate some functions $\estimand(m)$ of the conditional mean outcome that are not mean-square continuous, like the conditional average treatment effect (CATE) $\tau(x_0) = \mu_1(x_0)-\mu_0(x_0)$ at a point $x_0$. The CATE is an instance of the population-targeted average treatment effect \eqref{eq:targeted-ate} in which the target population's covariate distribution $p(x)$ is a spike $\delta_{x_0}(x)$ at $x=x_0$.
Here, the lack of overlap of this spiky target distribution $\delta_{x_0}$ with the distribution of the observed covariates causes the function not to be continuous.
% here, the lack of continuity arises from the lack of overlap of this spiky target distribution $\delta_{x_0}$ with the distribution of the observed covariates.
However, we can approximate the spike with a smoother density,
% We can approximate it \avi{this distribution?} by a continuous function by approximating the spike by a smoother density,
% e.g.,
taking $p(x)$ to be the density of a normal distribution with mean $x_0$ and covariance $h^2 I$ for small \emph{bandwidth} $h$, which will yield good approximation of $\tau(x_0)$ if $\tau$ is continuous at $x_0$. A refined version of this approach can be used for other estimands like the dose-response curve for a continuous-valued treatment, $\mu(w_0)=\E[m_{w_0}(X_i)]=\E[\int m_{w}(X_i) \delta_{w_0}(w)dw]$,
which fix some aspects of the treatment or population and average over others \citep[e.g.,][]{chernozhukov2018global, kennedy2017non}.

Other parameters identified by moment conditions can be estimated by a similar approach as well, including quantile treatment effects or average treatment effects identified by instrumental variables.
In essence, if a parameter $\theta_0$ satisfies the moment condition $\E[h(W_i, X_i, m, \theta)] = 0$
for any function $m$ that we can estimate, and we can propose a valid model $\model$ for this estimator's error  $\dm = \hat \mu - m$, then for any given $\theta$ we can form a generalized AIPW estimator $\hat\estimand_{\theta}$ of $\estimand_{\theta}(m) = \E[h(W_i,X_i,m,\theta)]$. To estimate $\theta_0$, all we must then do is solve $\hat \estimand_{\theta} \approx 0$.
See \citet{chernozhukov2016locally, singh2019biased}.



%%%%%%%%%%%%%%%%%%%%%%%%%%%
%%%%%%%%%%%%%%%%%%%%%%%%%%%
%%%%%%%%%%%%%%%%%%%%%%%%%%%

% \vspace{-.5cm}
\section{Matching for Balance}
\label{sec:matching}

%These emthods look for weights that satisfy a certain contraint on max imbalance and these reoudning results suggest the existence of a solution where the max imbalance is op(1/rootn). Connection to card sbw. This hasent been stated formally. but this result implies that. Not only implies this but implies that the weights don't have to be massive. How much is loost in this rounding The excesss variance induced in this procedure from the rounding paper. Have a sentence on this. These other papers are min variance so they are taking care of that. Assuming that hte variance of the weights is the same at every x. Corolally 2 would. Int consttraint SBW fits betters with the rest of the paper so...

The analysis of balance we have developed applies to any weights, integer-valued or otherwise: imbalance, augmentation, asymptotic normality, and the empirical process theory connecting model complexity to balance all carry over unchanged. What changes when we move to matching is how the weights are constructed. The optimization methods we have discussed---stable balancing weights, kernel balancing, entropy balancing---solve continuous problems and produce real-valued solutions. Matching uses integer weights, and this comes with both interpretability and efficiency advantages. Because unmatched controls have weight exactly zero, the cost of a study is driven by the number of controls followed up, not the number available. In longitudinal retrospective cohort studies, for instance, matching selects which control subjects to follow up \citep{stuart2008best}, and the savings can be the primary determinant of study cost. The question is whether we can reap the cost savings of integer weights without sacrificing the inferential guarantees of continuous ones.

The answer is essentially yes. We show that any good set of continuous weights can be rounded to integers with negligible change in imbalance (Proposition~\ref{prop:rounding}), yielding a matching estimator that is asymptotically unbiased and normal. This is both a usable method and a guarantee that extant matching-for-balance approaches like cardinality matching \citep{zubizarreta2014matching} and generalized optimal matching \citep{kallus2020generalized} can achieve the same inferential guarantees. All results are proven in Appendix~\ref{app:rounding}.

To state results precisely, we work with the matched difference in means estimator of the average treatment effect on the treated (e.g., \citealt{rosenbaum1985constructing}, \citealt{abadie2006large}, \citealt{zubizarreta2014matching}), which can be interpreted as a special case of the inverse probability weighting estimator $\hat\psi^{IPW}=\frac{1}{N_1} \sum_{i=1}^n W_i Y_i - \overset{\neg}{W}_i \theta_i Y_i$. Here we have written our estimator in terms of the rescaled weights $\theta_i=\gamma_i / \bar W$ appropriate when dividing by $N_1$ and written $\overset{\neg}W_i$ for $1-W_i$.
In this context, we write the design-conditional bias and the associated imbalance, which bounds the magnitude of the bias when $m \in \model$, as follows.
\begin{equation}
\label{eq:att-imbalance}
\begin{aligned}
\text{bias} &=  \frac{1}{N_1}\sum_{i=1}^n W_i m(0,X_i) - \frac{1}{N_1}\sum_{i=1}^n \overset{\neg}{W_i}\theta_i m(0, X_i) \\
\imbalance_{h,\model}(\theta) &= \max_{m \in \model} \abs*{ \frac{1}{N_1}\sum_{i=1}^n W_i m(0,X_i) - \frac{1}{N_1}\sum_{i=1}^n \overset{\neg}{W_i}\theta_i m(0, X_i)}.
\end{aligned}
\end{equation}
In the matched difference in means estimator, each control unit has an integer weight $\theta_i \in \{0, 1, \ldots\}$, describing the number of treated units to which it is matched.

\begin{remark}[Nearest-neighbor matching and the Lipschitz model]
Nearest-neighbor matching is, in fact, a matching-for-balance method: it is the $\sigma^2=0$ case of the minimax weights~\eqref{eq_balancing-weights} for the Lipschitz model $\model = \{f : \|f\|_{\text{Lip}(\rho)} \le \kappa\}$ \citep[Theorem 3]{kallus2020generalized}. The problem is that the Lipschitz model is large. The conditional bias of the nearest-neighbor matching estimator is $O_p(n^{-1/d})$ \citep[Theorem 1]{abadie2006large}, which is $o_p(n^{-1/2})$ only when $d \le 2$. In the language of Section~\ref{sec:trade-offs}, the Gaussian width of the Lipschitz ball grows too quickly with dimension for the fixed point to yield $\sqrt{n}$ rates. Matching with a smaller model---one whose complexity is controlled by stronger smoothness or sparsity assumptions---avoids this barrier; see the discussion in Section~\ref{sec:trade-offs}.
\end{remark}

A simple way to turn a weighting estimator into a matching one is through \emph{Bernoulli rounding}. We shave off the fractional parts $\hat r_i = \hat\theta_i - \floor{\hat\theta_i}$ of our weights and replace them with independent Bernoullis $B_i \sim \text{Bernoulli}(\hat r_i)$ with means matching what we have shaved. This gives us a \emph{randomized estimator} with integer weights $\floor{\hat\theta_i} + B_i$. The procedure is unbiased, as the expectation of the estimator is unchanged, but it increases the variance solely to obtain integral weights. By the law of total variance,
\begin{equation}
\label{eq:bernoulli-rounding-variance}
\begin{aligned}
\Var\left(\frac{1}{N_1} \sum_i \overset{\neg}{W}_i (\floor{\hat\theta_i} + B_i) Y_i\right)
&= \E\left(\frac{1}{N_1^2} \sum_i \overset{\neg}{W}_i \, \hat r_i(1 - \hat r_i) \, Y_i^2 \right) + \Var\left(\frac{1}{N_1} \sum_i \overset{\neg}{W}_i \hat\theta_i Y_i\right).
\end{aligned}
\end{equation}
This is not the irreducible cost of matching. We can do better. This is what matching costs: the variance of any negligibly-biased estimator using deterministic integer weights $\tilde\theta_i$ is
\begin{equation}
\label{eq:rounding-variance-deterministic}
\begin{aligned}
\Var\left(\frac{1}{N_1} \sum_i \overset{\neg}{W}_i \tilde\theta_i Y_i\right) &= \E\left(\frac{1}{N_1^2} \sum_i \overset{\neg}{W}_i \, \hat r_i(1 - \hat r_i) \, \Var\{Y_i \mid W_i, X_i\} \right) + \Var\left(\frac{1}{N_1} \sum_i \overset{\neg}{W}_i \hat\theta_i Y_i\right),
\end{aligned}
\end{equation}
where $\tilde\theta_i$ denotes any deterministic integer rounding of the real-valued weights $\hat\theta_i$, given mild conditions on $\hat\theta$ (see Proposition~\ref{prop:rounding-efficiency}). The difference between these two variances---call it $\Delta V$---is the mean-squared \emph{change in imbalance in the conditional mean $\mu$} from random rounding:
\begin{equation}
\label{eq:rounding-excess}
\begin{aligned}
\Delta V
&= \E\left(\frac{1}{N_1^2} \sum_i \overset{\neg}{W}_i \, \hat r_i(1 - \hat r_i) \, \mu(X_i)^2\right)
= \E\left[\left\{\frac{1}{N_1} \sum_i \overset{\neg}{W}_i (B_i - \hat r_i) \mu(X_i)\right\}^2\right].
\end{aligned}
\end{equation}
This is the same phenomenon that makes AIPW preferable to IPW with the true weights $\gammaipw$: mean-zero imbalance costs nothing in bias but inflates variance. Here, it is the \emph{change} in $\mu$-imbalance from rounding, rather than the $\mu$-imbalance itself. To avoid this cost, we need to round with negligible change in imbalance. Given a model $\model$ for $\mu$, we can ask for negligible change in imbalance by solving
this integer program, which minimizes the rounding-induced change uniformly over $m\in\model$.
\begin{equation}
\label{eq:round}
  \round{\hat\theta} = \argmin_{\substack{\theta\in \mathbb{Z}^n \\ \norm{\theta - \hat\theta}_\infty \leq 1 \\ \frac{1}{N_1} \sum_{i=1}^{n} \overset{\neg}{W}_i \theta_i = 1}} \roundingerror_{\model, \hat\theta}(\theta) \quad \text{for} \quad \roundingerror_{\model, \hat\theta}(\theta) = \max_{m\in\model} \abs*{\frac{1}{N_1}\sum_{i=1}^n \overset{\neg}{W}_i (\theta_i - \hat\theta_i) m(X_i)}.
\end{equation}
We show that, for models with bounded uniform entropy integral (BUEI)\footnotemark, the change in imbalance we get with this rounding scheme
is uniformly negligible. To keep things simple, we shall suppose we have started with a translation-invariant estimator. This means that our weights average to one and therefore can be rounded to integers with no change in sum.

This, of course, has implications for bias as well as variance. By the triangle inequality, the maximal imbalance with the rounded weights is no larger than the sum of our maximal rounding error and our maximal imbalance with our original weights.
\[ \imbalance_{h, \model}(\round{\hat\theta}) \le \imbalance_{h, \model}(\hat\theta) + \roundingerror_{\model, \htheta}(\round{\hat\theta}). \]
That means that if we start with weights achieving uniformly negligible imbalance over correctly specified model ($\mu \in \model$), the matching estimator we get by rounding has uniformly negligible imbalance too, and as a result we get asymptotics like we discussed in Proposition~\ref{proposition:clt} above. Inference is easy---the $t$-statistic based on our matched difference in means is asymptotically standard normal.

\footnotetext{Having a bounded uniform entropy integral is probably the most frequently-used sufficient condition for $\model$ to be Donsker. See \citep[Chapter 2.5]{van1996weak}}.

\begin{proposition}
\label{prop:rounding}
Suppose $\model$ has bounded uniform entropy integral. For any weights $\hat \theta$ with $\frac{1}{N_1}\sum_{i=1}^n \hat\theta_i = 1$,
the maximal rounding error $\roundingerror_{\model, \hat\theta}(\round{\hat\theta})$ resulting from our balance-preserving rounding scheme \eqref{eq:round} is  $o_p(n^{-1/2})$.
\end{proposition}


We prove it by working with an apparently simpler quantity. The rounding error in \eqref{eq:round} is a scaled version of what is called \emph{linear discrepancy} in the combinatorics and optimization literature: $\roundingerror_\model = N_1^{-1} \lindisc_\model$. A classical result in discrepancy theory \citep{lovasz1986discrepancy} shows that linear discrepancy is bounded by \emph{hereditary discrepancy}, i.e.,
\begin{equation}
\label{eq:hdisc}
\lindisc \le \hdisc \quad \text{where} \quad \hdisc(T) = \max_{I \subset \{1 \ldots n\}}\min_{\epsilon \in \{-1,1\}^n}\max_{t \in T} \abs*{\sum_{i\in I} \epsilon_i t_i}
\end{equation}
is the worst-case imbalance when we optimally two-color any subset of our units.

Rounding weights to integers is, at heart, a coloring problem: we're deciding which units to ``round up'' ($\epsilon_i = +1$) and which to ``round down'' ($\epsilon_i = -1$), trying to keep every function in our model balanced. The ``hereditary'' part---the max over all subsets $I$---is crucial. Ordinary discrepancy (for the full set) can be small by lucky cancellation, but hereditary discrepancy is small only when the structure genuinely allows good colorings.

To build intuition, consider a model containing all constant functions, so that perfectly balanced weights must satisfy $\sum_{i \in I} \epsilon_i = 0$. If the subset $I$ has even size $|I| = 2k$, this means assigning $k$ units to each color---a matched pair design. The inner minimization then asks: what is the best balance we can achieve when we optimally match the $2k$ units? And hereditary discrepancy asks: what is this worst-case optimal balance, over all subsets $I$ and all functions $m \in \model$?

Two caveats. First, discrepancy is defined in terms of sums $\sum_{i \in I} \epsilon_i m(X_i)$, not averages, so the connection to matched differences in means involves a factor of $|I|/2$. Second, the constraint $\sum_{i \in I} \epsilon_i = 0$ requires $|I|$ even; for odd-sized subsets, one ``group'' is larger by one, so discrepancy for any model containing all constants is infinite. For the same reason, models like this are not actually BUEI and do not satisfy the conditions of Proposition~\ref{prop:rounding}.

Why does rounding reduce to coloring? The idea is to round recursively on binary expansions: approximate each weight $q_i \in [0,1]$ to $M$ binary digits, then round the last digit optimally (choosing $\pm 2^{-M}$ to minimize imbalance), then the second-to-last, etc. Each step is a coloring problem on a subset---the units whose $k$th digit is 1---and costs at most $2^{-k}$ times the hereditary discrepancy. Summing the geometric series gives a bound of $\sum_{k=1}^M 2^{-k} \hdisc \to \hdisc$ as $M \to \infty$.

We show in Appendix~\ref{app:rounding}, following the approach of \citet{mendelson2011discrepancy}, that $\hdisc(\model) = o_p(n^{1/2})$ for BUEI classes. Because $\roundingerror \le N_1^{-1} \lindisc \le N_1^{-1} \hdisc$, this means that $\roundingerror_\model = o_p(n^{-1/2})$.

The cost of matching can be characterized precisely. For any regular matching estimator,
\begin{equation}
\label{eq:matching-cost}
\Var(\hat\psi^{\textup{match}}) - V^{\textup{eff}} \ge \frac{1}{N_1^2}\sum_{i=1}^n \overset{\neg}{W}_i\, r(X_i)\{1-r(X_i)\}\,v_0(X_i) + o_p(N_1^{-1}),
\end{equation}
where $V^{\textup{eff}}$ is the variance of the efficient estimator and $r(x) = \theta(x) - \floor{\theta(x)}$ is the fractional part of the propensity odds $\theta(x) = e(x)/\{1-e(x)\}$.

\begin{proposition}[Lower bound]
\label{prop:matching-cost}
Suppose $v_0(x) = \Var[Y \mid W=0, X=x]$ is bounded. For any \emph{regular} matching estimator---one whose integer weights achieve $o_p(1)$ imbalance in every $f \in L^2(P)$---the variance satisfies the lower bound~\eqref{eq:matching-cost}.
\end{proposition}

\begin{proposition}[Achievability]
\label{prop:matching-achievable}
Under the conditions of Proposition~\ref{prop:matching-cost}, balance-preserving rounding~\eqref{eq:round} achieves equality in~\eqref{eq:matching-cost}:
\[
\Var(\hat\psi^{\textup{match}}) - V^{\textup{eff}} = \frac{1}{N_1^2}\sum_{i=1}^n \overset{\neg}{W}_i\, r(X_i)\{1-r(X_i)\}\,v_0(X_i) + o_p(N_1^{-1}).
\]
\end{proposition}

\begin{corollary}
\label{cor:rounding}
Let $\Gamma :=  \sqrt{N_1^{-1}\sum_{i=1}^n \overset{\neg}W_i \htheta(X_i)^2}$ be the root-mean-squared weight. Suppose that each squared weight is negligible compared to their sum $n\Gamma^2$, and assume that the noise sequence $\varepsilon_1 \ldots \varepsilon_n$ is uniformly square-integrable conditional on the design.
If $m$ is in a BUEI model $\model$ and $\imbalance_{h,\model}(\htheta) = o_p(n^{-1/2})$,
then $\hat\psi^{match}=\frac{1}{N_1}\sum_{i=1}^n W_iY_i - \overset{\neg}W_i \round{\htheta}_i Y_i$, where $\round{\htheta}$ are the uniform-balance-preserving rounded weights \eqref{eq:round}, satisfies
\begin{equation}
\label{eq:asymptotic-linearity-rounded}
\hat\psi^{match}-\tilde\psi(m) = \frac{1}{N_1} \sum_{i=1}^{n}(W_i + \overset{\neg} W_i \round{\htheta}_i)\{ Y_i - m(W_i,X_i) \} + o_p(n^{-1/2}).
\end{equation}
Furthermore, if $\hat \mu$ is a consistent estimator of $\mu$,
 \begin{equation}
 \label{eq:asymptotic-normality-rounded}
  \frac{\hat\psi^{match} - \tilde\estimand(m)}{\sqrt{\hat V}} \to N(0,1) \quad \text{for} \quad
 \hat V = \frac{1}{N_1^2}\sum_{i=1}^n \{W_i + \overset{\neg}{W}_i (\round{\htheta}_i)^2\} \{Y_i - \hat \mu(W_i, X_i)\}^2.
\end{equation}
\end{corollary}

\begin{remark}
\label{rem:matching-variance}
The most interesting case is when the fractional part is the weight itself (i.e.\ the weight is in $[0,1]$), as rounding the whole weight to zero is what provides us the opportunity to avoid follow-up. This will be the case when the treatment probability $e(X_i)$ is well under $1/2$, as the weights $\hat\theta(X_i)$ estimate the inverse probability weights $e(X_i)/\{1-e(X_i)\}$. When this is true for all or almost all units, our rounded weighting estimator will have all or almost all weights in $\{0,1\}$ and the same asymptotic distribution as a simple difference in means in a randomized experiment in which each of $2N_1$ units received treatment with probability $1/2$. Relative to weighting the full set of $N_0$ control units, this comes at the cost described in \eqref{eq:bernoulli-rounding-variance}. But in cases where follow-up is costly, it is natural to compare its efficiency to other estimators requiring the same amount of follow-up, e.g.\ an IPW estimator based on our first $N_1$ control observations. Relative to that estimator, Proposition~\ref{prop:matching-efficiency} in Appendix~\ref{sec:rounding-and-efficiency} shows that our matching approach reduces asymptotic variance by a factor of $(\rho+1)/\{\rho + \mu_2/\mu_1^2\}$, where $\rho = v_1/v_0$ is the ratio of treated to control noise variance and $\mu_k = \E[\theta(X)^k \mid W=0]$. For equal noise ($\rho=1$) and small propensity ($e(X) \ll 1/2$, so $\theta \approx e$), this simplifies to $1/\{1+\Var(e(X)\mid W=0)/(2\E[e(X)\mid W=0]^2)\}$.
\end{remark}

 %We use this approach to establish the following characterization of the difference in variance between the matching estimator from Corollary~\ref{cor:rounding} and the IPW estimator it is derived from and show that the aforementioned equivalence to a simple difference in means in an experiment is essentially a special case of this bound.
%\begin{proposition}
%If $\model$ is nonparametric in the sense used in Section~\ref{sec:asymptotic}, and the conditions of Corollary~\ref{cor:rounding} are satisfied, then letting $\hat r(X_i) = \hat\theta(X_i) - \floor{\hat\theta(X_i)}$ be the fractional part of our weights,
%\begin{equation*}
%\label{eq:rounding-variance-bound}
%N_1 \times \left\{ \Var(\hat\psi^{match}) - \Var(\hat\psi^{ipw}) \right\} = \frac{1}{N_1}\sum_{i=1}^n \overset{\neg}{W}_i \hat r(X_i) \{ 1-\hat r(X_i) \} \Var[Y_i \mid W_i,X_i] + o_p(1).
%\end{equation*}
%\end{proposition}



%\footnotetext{That's probably too much detail. We need some way of saying `models like the ones we've been considering'. Or maybe at the end of the models section we could say something about sobolev models satisfying complexity conditions like donskerity and bounded uniform entropy integral. Maybe look at some footnote in AMLE re: how to talk about the relationship between this and Donskerity.}

This rounding approach, while reasonable, is not necessarily the best way to get an asymptotically unbiased and normal matching estimator.
Rather than finding non-integer weights that give us negligible imbalance and then rounding them in a way that controls the rounding-induced change in imbalance, we could simply find integer-valued weights that give us negligible imbalance.
We could do this, for example, by solving the following integer program to minimize the estimator's conditional variance subject to an imbalance budget $B_n \ll n^{-1/2}$,
\begin{equation}
\label{eq:integer-sbw}
\round{\hat\theta} = \argmin_{\substack{\theta\in\mathbb{Z}^n \ \text{ with } \\ \imbalance_{h, \model}(\theta) \leq B_n}} \  \frac1n\sum_{i=1}^n \overset{\neg}{W_i} \theta_i^2
\end{equation}
This integer-constrained balancing weights approach, which was called \emph{generalized optimal matching} by \citet{kallus2016generalized},
and the related \emph{cardinality matching} approach of \citet{zubizarreta2014matching}, can be understood using Proposition~\ref{prop:rounding} as well. For models in which a non-integer-constrained version of \eqref{eq:integer-sbw} achieves uniformly negligible imbalance---for example, RKHS models \citep[see][and Section~\ref{sec:doing-better}]{hirshberg2019minimax}---Corollary~\ref{cor:rounding} implies that there is a feasible solution and the result is, in fact, an optimal matching estimator.

\begin{corollary}
\label{cor:integer-sbw}
If $m$ is in an RKHS model $\model$ with bounded continuous kernel, then for an appropriately chosen sequence of budgets $B_n \ll n^{-1/2}$, the integer program \eqref{eq:integer-sbw} has a solution with probability tending to one and the corresponding difference in means estimator \smash{$\hat\psi^{match}=\frac{1}{N_1}\sum_{i=1}^n W_iY_i - \overset{\neg}W_i \round{\htheta}_i Y_i$} satisfies the asymptotic expansions \eqref{eq:asymptotic-linearity-rounded} and \eqref{eq:asymptotic-normality-rounded}, with variance achieving equality in the matching cost bound~\eqref{eq:matching-cost}.
\end{corollary}

The results in this section show that matching for balance makes statistical inference straightforward and can be more cost-effective than asymptotically efficient estimators by using the responses of a smaller number of units in the analysis.
%But if efficiency while controlling the number of responses we must record is all we are after, this approach is not necessarily optimal.
However, if our sole concern is efficiency while strictly limiting the number of responses we need to record, this approach may not be optimal.
We could, for example, replace the integrality constraint on our balancing weights \eqref{eq:integer-sbw} with a constraint on the number of nonzero weights. Sparse weighting ideas like these are an interesting area for future work.

%Remarks about.
%\begin{enumerate}
%\item how matching is a way to get sparsity but not necessarily the best way. Sparse weighting ideas have yet to be explored in the literature.
%\item We can characterize this budget in the appendix.
%\item imbalance-constrained vs. imbalance-penalized form used here to ensure that discretization increases variance rather than bias.
%\item Maybe? How we haven't constrained our weights to be non-negative in \eqref{cor:integer-sbw} because that's not compatible w/ arguments from the retargeted mean paper, but if we did have an $o_p(n^{-1/2})$
%bound on maximal imbalance with non-negative weights, the same result would follow for nonnegative integer weights.
%\end{enumerate}


%Finally, we construct the estimator $\tilde\psi = \frac{1}{n}\sum_{i=1}^{n}(1-W_i)\tilde\gamma_i Y_i$.
%As the following theorem states, this rounded estimator is asymptotically normal with nearly efficient variance.

%The theorem says that the weighting estimator $\hat\psi$ is asymptotically efficient, with variance $\hat V=N_1^{-2}\sum_i (1-W_i)\hat\gamma_i^2 \sigma^2(X_i)$.
%Thus, the variance of the matching estimator $\tilde\psi$ is only larger than the efficient variance by $N_1^{-2}\sum_i (1-W_i) (\tilde\gamma_i - \hgamma_i)^2 \sigma^2(X_i)$.
%This excess variance can be at least $1/4$, when $\gamma_i = \lfloor \gamma_i\rfloor = 1/2$.
%When $\mathcal{F}$ is symmetric, the excess variance will not be larger than $1/2$ by symmetry.

%We only need a consistent estimator of $\tilde V$ to perform inference, as opposed to knowledge of $\sigma(x)$. In particular, if we have a consistent estimator $\hat \mu_0(x)$ of $\mu_0(x)$, by Slutsky's theorem we can use $N_1^{-2}\sum_i (1-W_i)(\hat\gamma_i^2 + (\tilde\gamma_i - \hgamma_i)^2) (Y_i - \hat \mu_0(X_i))^2$ in place of $\tilde V$, and obtain asymptotically valid confidence intervals.

%\footnotetext{ . } %TODO: replace with the following or an updated version
%\footnotetext{\cite{bradic2019minimax} propose an augmented IPW estimator
%that is in fact optimally double robust, albeit in another sense
%common in the high-dimensional statistics literature. Assuming that
%$\gammaipw$ and $\mu_1$ are linear in a high dimensional basis $\phi(x)$
%that satisfies a restricted eigenvalue condition, they show that it is
%rate-optimal under the weakest possible conditions on the \emph{sparsity} of $\gammaipw$ and $\mu_1$
%as opposed to their smoothness. What differentiates this setting is that, as they show, the minimal conditions are
%indeed just those that permit faster-than-$n^{-1/4}$ rate of either $\mu_1$ of $\gammaipw$;
%no combination of sparsity conditions that does not do so is sufficient. As a result,
%it is essentially sufficient to combine the argument of Section~\ref{sec:role-of-augmentation}
%with a `reversed' version, like that of %\citet{wang2020debiased},
%that applies when $\gammaipw$ is sparser than $\mu_1$. \dah{TODO: Check with Yinchu or Jelena that this characterization is cool.}}

%This approach principally differs in what assumptions we place on the weights and outcome model, and how they are estimated. The double robust approach requires that both the propensity score and outcome model are sufficiently well-behaved so that we can estimate them at, e.g., a rate of $n^{-1/4}$, and the sample splitting prevents over-fitting. In contrast, the balancing approach  ``doubles down'' on the outcome model. As \citet{hirshberg2021augmented} discuss, while the requirements on the outcome model remain the same, by targeting the weights towards the outcome model we can achieve an efficient estimator even if it is not possible to estimate the propensity score well. Furthermore, rather than protecting against over-fitting the propensity score, the balancing approach undersmooths by prioritizing the bias from imbalance. Overall, we view the connections between these approaches --- and the role of sample splitting versus in-sample balance in particular --- to be an important area of future research.

%%%%%%%%%%%%%%%%%%%%%%%%%%%
%%%%%%%%%%%%%%%%%%%%%%%%%%%
%%%%%%%%%%%%%%%%%%%%%%%%%%%


%%%%%%%%%%%%%%%%%%%%%%%%%%%
%%%%%%%%%%%%%%%%%%%%%%%%%%%
%\subsection{Checking balance}
%\label{sec:}
%
%Standard practice in observational studies is checking covariate balance before and after adjustments.
%With the constrained formulation, as long as the weighting problem is feasible, balance is satisfied by construction.
%An infeasibility certificate can be helpful information to know that a form of adjustment is unattainable for the data at hand without relaxing some of the constraints or dropping observations.
%With the relaxed formulation, the weighting problem will always have a solution, but it is possible that balance is unsatisfactory.
%In view of this, it is advisable to always check covariate balance after weighting. Furthermore, it is advisable to check balance not only for the conjectured model class $\mathcal{M}$ but for other models that go beyond it to have a sense of the robustness of the adjustments.

%%%%%%%%%%%%%%%%%%%%%%%%%%%
%%%%%%%%%%%%%%%%%%%%%%%%%%%
%%%%%%%%%%%%%%%%%%%%%%%%%%%


\fi % end \ifmainbody

%\clearpage
\singlespacing
\ifmainbody
\ifarxiv
  \printbibliography
\else
  \bibliography{review}
\fi
\fi

\ifappendix
\clearpage
\appendix

\ifarxiv\begin{refsection}\fi

\section{Modulus of continuity: penalized and constrained forms}
\label{app:modulus-equivalence}

The penalized modulus~\eqref{eq:modulus-of-continuity} used in Section~\ref{sec:doing-better} is related by Lagrangian duality to the constrained formulation of \citet{donoho1994statistical} and \citet{armstrong2018optimal}. The constrained modulus is
\begin{equation}
\label{eq:modulus-constrained}
\bar\omega(\delta) = \sup\left\{\left|\frac{1}{n}\sum_{i=1}^n m_1(X_i) - \frac{1}{n}\sum_{i=1}^n m_2(X_i)\right| : m_1, m_2 \in \model,\, \norm{m_1 - m_2}_{L_2(\Pn)} \le \delta\right\},
\end{equation}
where $\norm{f}_{L_2(\Pn)}^2 = n^{-1}\sum_{i=1}^n f(X_i)^2$. The two formulations are related by the Fenchel--Legendre transform: $\omega(s) = \sup_{\delta \ge 0}\{\bar\omega(\delta) - \delta^2/(2s)\}$, or equivalently $\bar\omega(\delta) = \inf_{s > 0}\{\omega(s) + \delta^2/(2s)\}$, with the optimal $s$ and $\delta$ linked by $\delta = s\,\bar\omega'(\delta)$.

\section{Asymptotics}
\label{sec:appendix-asymptotic-normality}

We state our asymptotic results in the general AMLE framework of \citet{hirshberg2021augmented}, then specialize to the treatment-specific mean as a corollary. For a functional $\psi(m) = n^{-1}\sum_{i=1}^n h(W_i, X_i, m)$, the AIPW estimator is
\[
\hat\psi^{\aug} = \frac{1}{n}\sum_{i=1}^n h(W_i, X_i,\hat \mu) + \hgamma(W_i, X_i)\{Y_i - \hat \mu(W_i, X_i)\}.
\]

\begin{propositionprime}{\ref{proposition:clt-general}}
\label{proposition:appendix-clt-general}
Let $\Gamma = (n^{-1}\sum_{i=1}^n \hgamma(W_i, X_i)^2)^{1/2}$ be the root-mean-squared weight. Suppose that the noise $\varepsilon_i = Y_i - m(W_i, X_i)$ is well behaved in the sense that its conditional variance $\nu(W_i, X_i) := \E[\varepsilon_i^2 \mid W_i, X_i]$ is bounded away from zero and that it satisfies the uniform integrability condition
$\max_{i \le n} \E[\varepsilon_i^2 1(\varepsilon_i^2 \ge M) \mid W_1,X_1 \ldots W_n, X_n] \to_p 0$ as $M \to \infty$. Suppose that imbalance in $\dm = m - \hat \mu$ is negligible, i.e.\ $o_p(n^{-1/2}\Gamma)$, that the weights are infinitesimal in the sense that $\max_{i \le n} \hgamma(W_i, X_i)^2 = o_p(n\Gamma^2)$, and that $\hat \mu$ is consistent in the weighted $L_2$ sense
\begin{equation}
\label{eq:weighted-l2-consistency}
\sum_{i=1}^n \hgamma(W_i, X_i)^2 \{\hat \mu(W_i, X_i) - m(W_i, X_i)\}^2 = o_p(n\Gamma^2).
\end{equation}
Then
\[
\frac{\hat\psi^{\aug} - \tilde\psi}{(\hat V)^{1/2}} \to N(0,1) \quad \text{where} \quad \tilde\psi = \frac{1}{n}\sum_{i=1}^n h(W_i, X_i, m), \quad \hat V = \frac{1}{n^2}\sum_{i=1}^n \hgamma(W_i, X_i)^2 \{Y_i - \hat \mu(W_i, X_i)\}^2.
\]
\end{propositionprime}

The weighted $L_2$ consistency condition \eqref{eq:weighted-l2-consistency} is implied by sup-norm consistency, i.e.\ $\max_{i \le n} |\hat \mu(W_i, X_i) - m(W_i, X_i)| = o_p(1)$. It is also satisfied naturally by IPW regression estimators, where $\hat \mu$ is fit by weighted least squares using the weights $\hgamma$---in this case, the weighted residuals are orthogonal to the fitted values by construction.

\begin{propositionprime}{\ref{proposition:clt}}
\label{proposition:appendix-clt}
Let $\Gamma = (n^{-1}\sum_{i=1}^n W_i \hgamma(X_i)^2)^{1/2}$ be the root-mean-squared weight. Suppose that the noise $\varepsilon_i = Y_i - \mu_1(X_i)$ is well behaved in the sense that its conditional variance $\nu(X_i) = \E[\varepsilon_i^2 \mid X_i, W_i=1]$ is bounded away from zero and satisfies the uniform integrability condition $\max_{i \le n} \E[\varepsilon_i^2 1(\varepsilon_i^2 \ge M) \mid W_1, X_1 \ldots W_n, X_n] \to_p 0$ as $M \to \infty$. Suppose that imbalance in $\delta\mu_1 = \hat \mu_1 - \mu_1$ is negligible, i.e.\ $o_p(n^{-1/2}\Gamma)$, that the weights are infinitesimal in the sense that $\max_{i \le n} W_i \hgamma(X_i)^2 = o_p(n\Gamma^2)$, and that $\hat \mu_1$ is consistent in the weighted $L_2$ sense
\[
\sum_{i=1}^n W_i \hgamma(X_i)^2 \{\hat \mu_1(X_i) - \mu_1(X_i)\}^2 = o_p(n\Gamma^2).
\]
Then
\[
\frac{\hat\psi_1^{\aug} - \tilde\psi_1}{(\hat V)^{1/2}} \to N(0,1) \quad \text{where} \quad \tilde\psi_1 = \frac{1}{n}\sum_{i=1}^n \mu_1(X_i), \quad \hat V = \frac{1}{n^2}\sum_{i=1}^n W_i \hgamma(X_i)^2 \{Y_i - \hat \mu_1(X_i)\}^2.
\]
\end{propositionprime}

%\[ \sum_{i=1}^n \E[ \xi_{i,n}^2 1(\abs{\xi_{i,n}} \ge \epsilon \mid \mathcal{D}_n ] =
%   \sum_{i=1}^n \frac{W_i \hgamma^2(X_i) E[ \xi_{i,n}^2 1(\abs{\xi_{i,n}} \ge \epsilon \mid \mathcal{D}_n ] / ... \]

\begin{proof}[Proof of Propositions \ref{proposition:clt-general}$'$ and \ref{proposition:clt}$'$]
Proposition~\ref{proposition:clt}$'$ is an instance of Proposition~\ref{proposition:clt-general}$'$ with $h(W,X,m) = m(1,X)$ and $\hgamma(W,X) = W\hgamma(X)$, so it suffices to prove Proposition~\ref{proposition:clt-general}$'$. This proof has two steps. We'll show that $Z_n=(\hat\psi^{\aug} - \tilde\psi)/V_n^{1/2}$ is asymptotically standard normal for $V_n = n^{-2}\sum_{i=1}^n \hgamma(W_i, X_i)^2 \nu(W_i, X_i)$, then we'll show that $\hat V / V_n$ converges to one in probability, so Slutsky's theorem implies our claim. For the first step, we work with the decomposition
\[
\hat\psi^{\aug} - \tilde\psi =
\sqb*{ \frac{1}{n}\sum_{i=1}^n \hgamma(W_i, X_i) \dm(W_i, X_i) - \frac{1}{n}\sum_{i=1}^n \{h(W_i, X_i,\hat \mu) - h(W_i, X_i, m)\} } +
\sqb*{\frac{1}{n}\sum_{i=1}^n  \hgamma(W_i, X_i)\varepsilon_i  }.
\]
The first term is the imbalance in $\dm$, which is $o_p(n^{-1/2}\Gamma)$ by assumption and therefore, because $\nu$ is bounded away from zero, negligible relative to the denominator $V_n^{1/2}$, which exceeds $n^{-1}\Gamma \underline{\nu}^{1/2}$ for $\underline{\nu}=\min_{w,x}\nu(w,x)$.
It follows that $(\hat\psi^{\aug} - \tilde\psi)/V_n^{1/2}$ converges in probability to $V_n^{-1/2}$ times the second term.
Conditional on the design $\mathcal{D}_n=\set{W_i, X_i : i \le n}$, this is a standardized sum of independent mean zero random variables
$\xi_{i,n}= \hgamma(W_i, X_i)\varepsilon_i / (n^2 V_n)^{1/2}$, so
the central limit theorem implies it converges to a standard normal random variable if the sum is not dominated by too few terms.
To be more precise, letting $\hat\gamma$
refer to the weights estimated at sample size $n$ when defining $\xi_{i,n}$, this is a martingale difference array adapted to the $\sigma$-algebras $\mathcal{F}_{i,n}$ generated by $\mathcal{D}_n \cup \set{\varepsilon_1 \ldots \varepsilon_{i}}$, it is standardized in the sense that its sum of conditional variances $\sum_{i=1}^n \E[\xi_{i,n}^2 \mid \mathcal{F}_{i-1,n}]$ is one, and this conditional Lindeberg condition suffices to ensure its asymptotic normality \citep[Corollary 3.1]{hall2014martingale}:
\[ \sum_{i=1}^n \E[\xi_{i,n}^2 1(\xi_{i,n}^2 \ge \epsilon) \mid \mathcal{F}_{i-1,n}] \to_p 0 \quad \text{ for all } \epsilon > 0. \]
To show that our assumptions imply this Lindeberg condition, we decompose it and use H\"older's inequality to bound it as follows. Writing $\hgamma_i = \hgamma(W_i, X_i)$ and $\nu_i = \nu(W_i, X_i)$ for brevity,
\begin{align*}
    \sum_{i=1}^n \E[\xi_{i,n}^2 1(\xi_{i,n}^2 \ge \epsilon) \mid \mathcal{D}_n] &=  \frac{1}{n^2 V_n}\sum_{i=1}^n \hgamma_i^2 \E[\varepsilon_i^2 1(\varepsilon_i^2 \ge M) \mid \mathcal{D}_n] \\& +
    \frac{1}{n^2 V_n}\sum_{i=1}^n \hgamma_i^2 \E[\varepsilon_i^2 1(M \ge \varepsilon_i^2 \ge n^2 V_n) \mid \mathcal{D}_n] \\
&\le \frac{1}{n^2 V_n}\sum_{i=1}^n \hgamma_i^2 \cdot  \max_{i \le n}\E[\varepsilon_i^2 1(\varepsilon_i^2 \ge M) \mid \mathcal{D}_n]\\
&+ \frac{1}{n^2 V_n}\sum_{i=1}^n \hgamma_i^2 \cdot  M\max_{i \le n}1(\hgamma_i^2 M \ge n^2 V_n).
\end{align*}
Both terms in this bound have the same leading factor, the ratio of $n^{-2}\sum_{i=1}^n \hgamma_i^2$ and $n^{-2}\sum_{i=1}^n \hgamma_i^2 \nu_i$, which is $O_p(1)$ because $\nu$ is bounded away from zero. Thus, it suffices to show that the sum of the second factors is $o_p(1)$. Our uniform integrability assumption implies that for large enough $M$, the probability that the first term exceeds any threshold $\delta>0$ is arbitrarily small.
Moreover, for this $M$, the second term tends to zero with probability tending to one, as given our assumption that $\nu$ is bounded away from zero, the claim that
$\max_{i \le n}\hgamma_i^2 / \sum_{i \le n}\hgamma_i^2 \nu_i \to_p 0$
reduces to our assumption that the weights are infinitesimal. This concludes the first step of our proof.

For the second step, it is equivalent to show that $(\hat V - V_n)/V_n \to 0$ in probability. Writing $\dm_i = \hat \mu(W_i, X_i) - m(W_i, X_i)$, we work with the decomposition
\begin{align*}
\frac{\hat V - V_n}{V_n}
&= \frac{1}{n^2}\sum_{i=1}^n \alpha_i [ \{\varepsilon_i - \dm_i\}^2 - \nu_i] \ \text{ for } \ \alpha_i = \frac{\hgamma_i^2}{V_n} \\
&= \frac{1}{n^2} \sum_{i=1}^n \alpha_i \dm_i^2 - \frac{2}{n^2} \sum_{i=1}^{n} \alpha_i \varepsilon_i \dm_i +
\frac{1}{n^2} \sum_{i=1}^n \alpha_i \{\varepsilon_i^2 - \nu_i\}.
\end{align*}
We will show that each term converges to zero. Note that $n^{-2}\sum_{i=1}^n \alpha_i = \Gamma^2/V_n$ is, as discussed above, bounded because $\nu$ is bounded away from zero.
\begin{enumerate}
\item
The first term equals $(n^2 V_n)^{-1} \sum_{i=1}^n \hgamma_i^2 \dm_i^2$. The numerator is $o_p(n\Gamma^2)$ by our weighted $L_2$ consistency assumption \eqref{eq:weighted-l2-consistency}, and the denominator is $\Theta_p(n\Gamma^2)$ because $\nu$ is bounded away from zero. So this term is $o_p(1)$.

\item For the cross-term, we use the Cauchy-Schwarz inequality:
\[
\abs*{\frac{1}{n^2}\sum_{i=1}^n \alpha_i \varepsilon_i \dm_i}^2 \le \frac{1}{n^2}\sum_{i=1}^n \alpha_i \varepsilon_i^2 \cdot \frac{1}{n^2}\sum_{i=1}^n \alpha_i \dm_i^2.
\]
The first factor on the right-hand side is $(n^2 V_n)^{-1}\sum_{i=1}^n \hgamma_i^2 \varepsilon_i^2$, which converges in probability to one by step 3 below. The second factor is the first term in our decomposition, which we have just shown is $o_p(1)$. So the cross-term is $o_p(1)$.
\item The third term is difference between the sum of squares $\sum_{i=1}^n \xi_{i,n}^2$ and the sum of conditional variances $\sum_{i=1}^n \E[\xi_{i,n}^2 \mid \mathcal{F}_{i-1,n}]$ of the standardized martingale difference array $\xi_{i,n}$, and the aforementioned conditional Lindeberg condition is sufficient to imply that the two converge in probability \citep[Theorem 2.23]{hall2014martingale}.
\end{enumerate}
\end{proof}

%\subsection{Asymptotic linearity and regularity}
%\label{sec:regularity}

%\section{Computation}
%\label{sec:computation}
%\subsection{Solving the optimization problem}
%\label{sec:solving}

%A computational advantage of the balancing weights is that it is a convex quadratic program under certain conditions \citep{Boyd2004}.
% , at least for certain choices of model class $\model$, dispersion function $\dispersion(\cdot)$, and monotone function $h_\lambda(\cdot)$ \citep{boyd2004convex}.
% Above we have mainly focused on a formulation of this problem that simultaneously minimizes measures of covariate imbalances and dispersion of the weights.
% As we discussed in Section \ref{sec:primal}, equivalent formulations minimize one of these measures subject to constraints on the other.
% For example, \cite{zubizarreta2015stable} minimizes the variance of the weights subject to approximate balance constraints, whereas \cite{santacatterina2018optimal} minimize imbalances subject to a dispersion threshold.
% While these formulations can have different implications for practice, they are equivalent and admit solution algorithms that can find the optimal weights in polynomial time.
%
%Traditionally, quadratic programs are solved using active-set or interior-point methods (see, e.g., Chapter 16 of \citealt{nocedal2006numerical}).
%Usually, active-set methods perform better than interior-point methods when the number of covariates is large or warm starts are required.\footnote{A major drawback of active-set methods is that the worst-case complexity is an exponential function of the number of the number of constraints.}
%Otherwise, interior-point methods often outperform active-set methods in terms of scalability and speed.
%Neither approach, however, scales well to very large problems.
%However, a major drawback of active-set methods is that the worst-case complexity is an exponential function of the number of the number of constraints.
%On the other hand, interior-point algorithms run in polynomial time, but they do not scale well to very large problems.

%A more recent solution is based on the alternating direction method of multipliers (ADMM; \citealt{boyd2011distributed}).
%In particular, \cite{stellato2020osqp} provide an implementation of ADMM tailored for quadratic problems that is robust and scalable to very large instances of these problems.
%Their technique is implemented in the open-source operator splitting solver for quadratic programs (OSQP).
%In a recent empirical study of weighting methods for causal inference in large data sets, \cite{kim2021weighting} illustrate the performance of OSQP and other open-source and commercial solvers such as quadprog and MOSEK, respectively.
%In simulated data sets with one million observations, OSQP was able to provide up to an order of magnitude of improvement in computational time over the second best performing solver, MOSEK, whereas quadprog failed to provide a solution in these instances.
%In instances with 10,000 observations (the largest sample size for which quadprog could could find a solution), the computational time improvement of OSQP over quadprog was of two orders of magnitude.
%In general, OSQP also showed to be faster in finding a solution than modeling approaches to weighting including the packages \texttt{admm.lasso} and \texttt{ranger} in R, and comparable to \texttt{glm}.


\section{More on Models}
\label{sec:more-on-models}


\subsection{Models defined in terms of $\ell_q$ and $L_q$ norms}
\label{sec:lq-models}
In Section~\ref{sec:balance_models}, we have discussed models defined in terms of $\ell_2$ norms on coefficients (\ref{eq:bias_linear}-\ref{eq:imbalance_phi_infinite}). It is common to use models defined analogously in terms of the $\ell_1$ norm, or another $\ell_q$ norm, on the coefficients. For example, the
largest imbalance in any coordinate $X_{ij}$ is the maximal imbalance for a variant of the model \eqref{eq:bias_linear} using the $\ell_1$ norm and
the sum of absolute imbalances in the coordinates arises from a model based on the $\ell_{\infty}$ norm.
In models with a low-dimensional basis, such as models without interactions, these are roughly equivalent, in the sense that the unit balls of the $\ell_1$, $\ell_2$, and $\ell_{\infty}$ norms on $\R^p$ satisfy $p^{-1} B_{\infty} \subseteq p^{-1/2} B_2 \subseteq B_1 \subseteq B_2 \subseteq B_{\infty}$.
For models requiring a high dimensional basis ($p$ large), the same relationship does not spell equivalence. Geometrically, $B_1$ is a `diamond' inscribed in $B_2$, a `round ball', inscribed in $B_{\infty}$,
the `cube' with side-length $2$. In $\R^2$ these are of increasing but similar size, whereas in high dimensional spaces almost all the volume of $B_{\infty}$ is in the corners, outside $B_2$, and similarly almost all the volume of $B_2$ is outside of $B_1$. In fact, the size of $B_1$ is, in a meaningful sense, almost independent of dimension $p$. This fact lies at the heart of the tractability of some problems in high dimensional statistics and their connection to $\ell_1$ regularization, including the problem of balancing high dimensional linear models \citep[e.g.,][]{athey2018approximate, tan2020regularized}.

The RKHS model discussed above \eqref{eq:imbalance_phi_infinite} is based on a \emph{weighted} $\ell_2$ norm on infinitely many scaled coefficients $\beta_j / \lambda_j^{1/2}$.
When the scale factors $\lambda_j^{1/2}$ decrease sufficiently fast, the corresponding coefficients $\beta_j$ must decay rapidly, and as a result this model behaves like one defined in terms of a low-dimensional basis. In particular, models defined analogously in terms of other $\ell_q$ norms on $\beta_j / \lambda_j^{1/2}$ are of similar size.

Similarly, generalizing the Sobolev balls we have discussed above, based on $L_2$ norms on derivatives \eqref{eq:sobolev-derivative}, are Sobolev balls defined analogously in terms of other $L_q$ norms. These are of comparable size for different $q$ if the functions in them are sufficiently smooth, e.g., $L_2$ and $L_\infty$ Sobolev balls of functions on $\R^p$ are of similar size if we include derivatives of order greater than $p/2$ in the norm \citep[e.g.,][Corollary 4]{nickl2007bracketing}. %Although $L_q$ Sobolev balls do not, in general, have equivalent characterizations in terms of Fourier coefficients like the $L_2$ ones do \eqref{eq:sobolev-fourier}, coefficient decay
%as discussed above can still offer a sense of what is going on.
%of the similarity of models defined in terms of different $\ell_q$ norms on decaying coefficients gives a sense of  what is going on, as
%for $q > 2$ (e.g. $q=\infty$) they \emph{contain} balls defined in terms of the $\ell_{q'}$ norm on the scaled Fourier coefficients $\beta_j / \lambda_j^{1/2}$, where $q'$ (e.g. $q'=1$ for $q=\infty$) is the H\"older conjugate of $q$ \citep[]{ something}.

\subsection{Models defined in terms of variation norms}
\label{sec:variation-models}
To check balance in covariate distributions, researchers will sometimes compare histograms of the covariates for the re-weighted treated units to those of the population, often summarizing this comparison by the marginal Kolmogorov-Smirnov (KS) statistics.
This measure of imbalance corresponds to an \emph{additive model} for the outcome, $\mu_1(x)=\sum_{j=1}^d m_{1j}(x_j)$, with a component function $m_{1j}(x_j)$ for the $j$th covariate. In particular, the maximal KS statistic is the maximal imbalance for
an additive model with a bounded summed componentwise \emph{total variation},
$\model = \set{ \sum_{j=1}^d f_j(x_j) : \sum_{j=1}^d \norm{f_j}_{TV} \le 1}$ where, for differentiable functions, $\norm{f_j}_{TV} = \sum_j \int \abs{f_j'(x)}dx$.\footnote{See \cite{fang2021multivariate} for a thorough discussion of total variation.
One useful interpretation is that $\norm{f}_{TV}$ is the vertical distance a pen travels when drawing the function $f$: for differentiable functions, it is $\int \abs{f'(x)}dx$
and the seminorm associated with the model is $\sum_j \abs{f_j'(x)}dx$.}
This is a more general approach to controlling bias than comparing covariate means, but it fails to control potential bias arising from imbalance in covariate interactions. This additive model is, however, essentially contained in a very natural model that does include interactions: the unit ball of the space of functions with bounded \emph{sectional variation} \citep{bibaut2019fast, fang2021multivariate}. The norm for this space is similar to the mixed-smoothness Sobolev norm from Equation~\ref{eq:ms-sobolev-derivative}, depending on the same sum of derivatives, but evaluates it via mix of $L_1$-type averaging and evaluation at zero, so it does not coincide with any Sobolev norm exactly \citep[see Equations 23 and 25 in][]{fang2021multivariate}.

\subsection{Sobolev and Reproducing Kernel Hilbert Spaces}
\label{sec:sobolev-spaces}
The particular norms we state in Section~\ref{sec:function-space-perspective}'s Equations~\ref{eq:sobolev-derivative} and \ref{eq:ms-sobolev-derivative} are the norms of the first-order isotropic and dominating mixed smoothness Sobolev spaces of bivariate functions on the cube $[0, 1]^2$ \citep[e.g.,][Section 2.1]{kuhn2015approximation}. Sometimes this is described as a norm on functions on $[-1,1]^2$ that periodic with period $2$ in all coordinates, and in Equation~\ref{eq:sobolev-derivative} the average over the uniform distribution on
$[0,1]^2$ is replaced with an average over the uniform distribution over $[-1,1]^2$; the arbitrary functions on $[0,1]^2$ can be extended to such functions on $[-1,1]^2$ via `reflection', letting $f(\pm x_1, \pm x_2)=f(x_1, x_2)$. For these `reflections', the norm in Equation~\ref{eq:sobolev-derivative} and the variant that averages over $[-1,1]^2$ match. If instead of the cube $[0,1]^2$, we define our space on all of $\R^2$ by averaging our derivatives over the normal distribution, we get a RKHS with a norm characterized in terms of a Hermite polynomial rather than a cosine basis.


Choosing as the model the unit ball $\calB_\calH$ of an RKHS $\calH$ is computationally convenient due to the so-called \emph{kernel trick}, which can be used to reduce optimization problems involving infinite series (as in \ref{eq:imbalance_phi_infinite}) to finite dimensional problems. What lies behind it is the so-called \emph{reproducing property}: for each point $x$, there is an element $K(x,\cdot) \in \calH$ with the property that we can evaluate functions $f \in \calH$ at the point $x$ by taking their inner product with $K(x,\cdot)$, i.e., with the property that $f(x)=\langle K(x,\cdot), f\rangle$ for all $f \in \calH$. As a result, if we take regularized optimization problems over $f \in \calH$ and rewrite each instance of $f(X_i)$ as an inner product with $K(X_i,\cdot)$, it often becomes clear that the solution must lie in the span of $K(X_1,\cdot) \ldots K(X_n, \cdot)$. Thus, knowing the solution has the form $f(\cdot)=\sum_{i=1}^n \alpha_i K(X_i, \cdot)$ for $\alpha \in \R^n$, we can solve our optimization problem by optimizing over a finite-dimensional vector $\R^n$ instead of an infinite-dimensional function $f \in \calH$.
This is summarized in so-called \emph{representer theorems}. Taking this approach, we can simplify the maximal imbalance over the ball $\calB_\calH$ in terms of its kernel $K$ as follows.

\begin{equation}
\label{eq:imbalance_rkhs}
\imbalance_{\calB_\calH}(\gamma)^2 = \frac{1}{n^2}\sum_{j=1}^n \sum_{i=1}^n K(X_i,X_j) - 2 W_i K(X_i,X_j) \gamma_i +  W_iW_j K(X_i,X_j) \gamma_i\gamma_j.
\end{equation}
\noindent This is helpful when we know the kernel of our space $\calH$.
See, e.g., \citet{novak2018reproducing} for the kernels associated with
certain Sobolev spaces and techniques for calculating kernels more generally,
and \citet[Section 2.3]{hirshberg2019minimax} and references therein for additional detail on RKHSes more generally.


Balancing weights for Sobolev space models and RKHS models more generally are discussed in \citet{Hazlett2019}, \citet{hirshberg2019minimax}, \citet{kallus2020generalized}, \citet{singh2020reproducing}, and \citet{wong2018kernel}.


% \vspace{-.5cm}
\section{Augmentation and strong notions of convergence}
\label{sec:strong-norm-consistency}

In Section~\ref{sec:role-of-augmentation}, we discussed the impact of augmentation on
imbalance by focusing on mean-square consistency of $\hat \mu$ and the equicontinuity of
the imbalance process $f \to n^{-1}\sum_{i=1}^n f(X_i) - W_i \gammaipw(X_i) f(X_i)$.

An argument that appeared earlier  \citep[e.g., in][]{athey2018approximate, kallus2020generalized, wong2018kernel}
relied on the assumption that $\hat \mu$ was consistent in a stronger sense:
instead of requiring the mean squared error $\norm{\dm}_{L_2(P)}$ to vanish, it required the gauge of the error $\norm{\dm}_{\model}$ to do so.
In other words, it required that $\dm$ be in a vanishingly small scaled-down version of $\model$, $\model_n = r_n \model$ for $r_n \to 0$;
in the case of the model $\model = \set{ \beta \cdot x : \norm{\beta}_1 \le 1}$ \eqref{eq:ipw_linf_imbalance}, this means that we must use an estimate of the form
$\hat \mu_1(x) = \hat\beta \cdot \phi(x)$, and the coefficients must be consistent in the $\ell_1$ norm, as $\norm{\hat \mu_1-\mu_1}_{\model} = \norm{\hat\beta-\beta}_1$.
When we have consistency in this sense, $\model_n$ is a valid model for $\dm$,
and the maximal imbalance $\imbalance_{\model_n}(\hgamma) = r_n \imbalance_{\model}(\hgamma)$ is negligible
if $\imbalance_{\model}(\hgamma)$ is $O_p(n^{-1/2})$;
if instead $\imbalance_{\model}(\hgamma)$ is $O_p(\sqrt{\log(p) / n})$ as discussed above for this model,
it suffices to have gauge (e.g., $\ell_1$) consistency with an $o_p(\sqrt{1/\log(p)})$ rate.
While this argument is perhaps easier to internalize, it is often preferable to use the one above,
as  this stronger form of consistency can distract us with thinking about extraneous aspects of the data.
For example, multicollinearity of the basis functions $\phi_j$
has no effect on our ability to estimate the conditional mean \emph{function} $\mu_1(x)=\beta \cdot \phi(x)$,
but it can make it impossible to accurately estimate its \emph{coefficients} $\beta$
\citep[see, e.g.,][]{chatterjee2013assumptionless}.

\section{Rounding}
\label{app:rounding}

\subsection{Rounding and Imbalance}

In this section, we will prove the rounding error bound from Section~\ref{sec:matching}. We will let $P_{\calX_I} \calF$ denote the set of \emph{coordinate projections} of the functions in $\calF$ on the indices $i_1 \ldots i_{\abs{I}}$ in $I$, i.e., $P_{\calX_I}\calF = \{ [f(X_{i_1}) \ldots f(X_{i_{\abs{I}}})] \ : \ f \in \calF \}$, and $P_{\calX}$ the same for the whole set of indices $I=\{1 \ldots n\}$. And we define the \emph{linear discrepancy} and \emph{hereditary discrepancy} of a set of vectors $T \subset \R^n$ as follows.

\begin{equation}
\begin{aligned}
  \label{eq:discrepancy}
  \lindisc(T) &= \max_{q \in [0,1]}\min_{r \in \{0,1\}}\max_{t \in T}\abs*{\sum_{i=1}^n (q_i-r_i) t_i} \\
  \hdisc(T) &= \max_{I \subset \{1 \ldots n\}}\min_{\epsilon \in \{-1,1\}^n}\max_{t \in T} \abs*{\sum_{i\in I} \epsilon_i t_i}
\end{aligned}
\end{equation}

Linear discrepancy bounds the degree to which balance in the vectors in $T$ is impacted
when we optimally round a vector of weights $q \in [0,1]^n$ into binary weights $r \in \{0,1\}$.
Hereditary discrepancy bounds the imbalance we get in those vectors when we optimally
color subsets of our units $\pm 1$, i.e. when we split subsets of our units into two groups and
compare the averages of the vectors in $T$. \citet[Corollary 1]{lovasz1986discrepancy} shows that rounding
is no harder than this kind of coloring, i.e. that
\begin{equation}
\label{eq:lindisc-hdisc}
\lindisc(T) \le \hdisc(T).
\end{equation}
The formal proof we will leave to \citet{lovasz1986discrepancy}, but for the
sake of offering a complete argument, we will sketch one. The idea is to think of a rounding process that
works recursively, term-by-term  and `from back to front', on the binary expansions $q_i = \sum_{j=0}^{\infty} 2^{-j} q_{ij}$
of the elements $q_i \in [0,1]$ of the vector we are rounding, and use hereditary discrepancy bound the changes induced
by our rounding process at each step.
\begin{proof}[Proof Sketch for Equation~\ref{eq:lindisc-hdisc}]
Having approximated the elements $q_i$ to arbitrary precision using an $M$-digit expansion
$q_i^M=\sum_{j=0}^M q_{ij} 2^{-j}$,
we'll first show that we can find $M-1$ digit approximations $q_i^{M-1}$ for which
$\max_{t \in T} \sum_{i=1}^n (q_i^{M-1}-q_i^M)t_i$ is small, then similarly
find $M-2$ digit approximations to $q_i^{M-1}$, etc, until we have elements $q_i^0$ in $\{0,1\}$.
To approximate the $k$-digit expansions by $k-1$-digit ones, we'll find the set of coordinates
$I_k$ where its last digit $q^k_{ik}$ is one and add $\pm 2^{-k}$. Doing this optimally, the
change in imbalance $\max_{f \in \F} \sum_{i = 1}^n (q_i^{M-1}-q_i^M)t_i$ is at most $2^{-k}\hdisc(T)$.
And we can bound the magnitude of the total change in imbalance over our $k$ rounding steps by
summing our bounds at each level, giving the bound $\sum_{k=1}^{M} 2^{-k} \hdisc(T) \le \sum_{k=1}^{\infty} 2^{-k} \hdisc(T)= \hdisc(T)$.
\end{proof}

What this tells us is that we can bound our rounding-induced change in imbalance in Equation~\ref{eq:round} by the (appropriately scaled) hereditary discrepancy for the vector of coordinate projections of the functions in
$\F$, i.e., by $n^{-1}\hdisc(P_{\calX} \F)$. To show that this $o_p(n^{-1/2})$, it suffices to show that $\hdisc(P_{\calX} \F) = o_p(n^{1/2})$. We will do this for classes $\F$ of bounded functions with bounded uniform entropy integral, first proving a finite sample bound then showing that bound is $o_p(n^{1/2})$.

\begin{lemma}
\label{lemm:bdd_ent_class}
Let $\F$ be a class of functions that is star-shaped around zero with $\max_{f \in \F}\norm{f}_{\infty} \le b$
and let $r > 0$ satisfy the Rademacher complexity fixed point inequality $2b^3 R_n ( \{f \in \F : P f^2 \le r^2\} ) \le r^2$.
With probability $1-\exp(-nr^2/8b^2)$, for every $I \subseteq \{ 1 \ldots n \}$ and $s \ge 0$,
$$
\begin{aligned}
\inf_{\epsilon \in \{-1,1\}^n} \frac{1}{\sqrt{n}} \max_{f \in \calF}\abs*{\sum_{i \in I} \epsilon_i f(X_i)}
&\le c_1b\left\{ s \sqrt{\frac{\abs{I}}{n}} \expo\left(-c_2 2^{2^{s+1}}\right) + n^{-1/2} 2^{2^{s+1}} 2^{(s+1)/2} \right\} \\
&+c_3\int_0^{2\sqrt{\frac{\log(\abs{I})}{\log(4/3)}}\max(\delta,2r)}  \sup_{Q} \sqrt{\log\calN(\calF, \norm{\cdot}_{L_2(Q)}, u\norm{F}_{L_2(Q)}}) du \\
&\qfor  \delta \ \text{ satisfying }  \ \log_2 \calN(\calF, \norm{\cdot}_{L_2(P)}, \delta) \le 2^{s+1}.
\end{aligned}
$$
Here $F$ is any function satisfying $F(x) \ge \abs{f(x)}$ for all $f \in \F$, $\calN(\calF, \norm{\cdot}, \varepsilon)$
is $\F$'s $\norm{\cdot}$ $\epsilon$-covering number, $R_n(\mathcal{G})$ is the Rademacher complexity of the class $\mathcal{G}$,
$\abs{I}$ denotes the cardinality of the set $I$,
$c_1 \ldots c_3$ are universal constants, and $\expo(x)=\exp(x)/\{1-\exp(x)\}$.
\end{lemma}
Because this bound is increasing in $\abs{I}$, replacing $\abs{I}$ with $n$ throughout gives a single bound that holds for all $I$, i.e. a bound on hereditary discrepancy. We summarize its asymptotic implications as follows.
\begin{corollary}
\label{cor:hdisc-negligible}
If $\F$ is a class of bounded functions with bounded uniform entropy integral that is star-shaped around zero,
$\hdisc(\P_{\calX} \F) = o_p(n^{1/2})$.
\end{corollary}


These results are similar to discrepancy bounds in \citet[Theorems 5.1 and A]{mendelson2011discrepancy}
for \emph{subgaussian classes} $\F$, i.e. classes of functions on which the $L_2(P)$ and subgaussian norms are equivalent.
Our proofs are essentially simplified versions of those results, as some of the complexity in \citet[Theorem 5.1]{mendelson2011discrepancy}
is in its comparison of empirical and population two-norms, and that is substantially simplified by the uniformity of the
bounds we get when working with classes with bounded uniform entropy integral.

\begin{lemma}
\label{lem:empirical-and-population-norm-equivalence}
Let $\F$ be a class of functions that is star-shaped around zero and with $\max_{f \in \F}\norm{f}_{\infty} \le b$
and let $r > 0$ satisfy the Rademacher complexity fixed point inequality $2b^3 R_n ( \{f \in \F : \Pn f^2 \le r^2\} ) \le r^2$.
With probability $1-\exp(-nr^2/8b^2)$, all functions $f \in \F$ satisfy $\norm{f}_{L_2(\Pn)} \le 2\max(\norm{f}_{L_2(\P)}, 2r)$.
\end{lemma}

\begin{proof}[Proof of Lemma~\ref{lemm:bdd_ent_class}]
To show that the bound from Lemma~\ref{lemm:bdd_ent_class} is satisfied, we will show that there exists a sequence $\epsilon \in \{-1,1\}^n$ that satisfies the bound. We will do this iteratively, constructing $\epsilon$ as a sum $\sum_{j' \le j_0} \eta^j$ of `partial assignments' $\eta^j \in \{-1,0,1\}^n$. At each step $j$, we
update only the `unassigned' indices, i.e. we find a vector $\eta^j$ with the property that
$\eta^j_i = 0$ unless $\eta_i^{j'}=0$ for all $j' < j$. As a result, the sum $\sum_{j' \le j} \eta^{j'}$ is itself a partial assignment for all $j$. And we will repeat this process until we have no zeros left, i.e., until $\sum_{j' \le j} \eta^{j'} \in \{-1,1\}^n$; that
$j$ will be $j_0$. Our approach, following \citet{mendelson2011discrepancy}
uses the entropy method from discrepancy theory --- indirectly through \citet[Corollary 3.4]{mendelson2011discrepancy} --- in combination with chaining.

\textbf{Notation.} Throughout, we'll write $c_1$, $c_2$, \ldots to denote universal constants.\footnote{Constants will be renamed for use the lemma's statement. $c_1$ here does not correspond to $c_1$ there, etc.} We'll write $I$ for a subset of the indices $\{1 \ldots n\}$ and $\abs{I}$ for its cardinality.
In our iterative process, $I$ will be the set of indices for which we have not yet assigned $\eta_i$ to $\pm 1$. $\pi_0, \pi_1, \ldots$ will be an admissible sequence of maps from $\F$ to itself, i.e. a sequence of functions for which the image $\pi_s(T)$ is a discrete set taking on at most $2^{2^s}$ values that are nested in the sense that $\pi_{s}(f)=\pi_{s}(g)$ implies $\pi_{s-1}(f)=\pi_{s-1}(g)$. We will choose these maps later, but we can think of $\pi_1(f), \pi_2(f), \ldots$ as a sequence of approximations to $f$ at increasing `resolution $s$', where we say we are at resolution $s$ if we allow our approximations to take on $2^{2^s}$ values.


We will keep track of the unassigned indices with a set $I^j=\{i \in I \ : \sum_{j'<j}\epsilon^{j'} = 0\}$. We start with $I^0 = I$ and proceed inductively. Suppose we are $j \ge 0$ steps into this process and our set of unassigned indices is $I^j$. \citet[Corollary 3.4]{mendelson2011discrepancy} says that there exists a vector $\eta^j \in \{-1,0,1\}^{n}$ with (i) $\eta_i^j = 0$ for $i \notin I^j$, (ii) $\eta_{I^j}$ having between $\abs{I^j}/4$ and $3\abs{I^j}/4$ zeros, and (iii) for all $t \in T$,
\begin{equation}
\label{eq:rounding-one-step}
\begin{aligned}
\abs*{\sum_{i \in I} \eta_i^j f(X_i)} \leq \sum_{s\geq 1} Q^j_s\norm*{P_{\calX_{I^j}}\{\pi_s(f) - \pi_{s-1}(f)\}}_2
\quad \text{for} \quad
Q^j_s &= c_1 \begin{cases}
  \exp(-c_2 \sqrt{\abs{I^j}}), &\text{if } s<s_{j},\\
  1, &\text{if } s = s_{j},\\
  2^{s/2}, &\text{if } s>s_{j}
\end{cases} \\
\quad \text{and} \quad
s_{j} &= \max\cb*{s \ : 2^{2^{s+1}}\leq c_3\abs{I^j} }.
\end{aligned}
\end{equation}
Because the number of zeros $\abs{I^j}$ at each step is at least $1/4$ and at most $3/4$ the number $\abs{I^{j-1}}$ we had in the previous step, $\abs{I^j} \in [ (1/4)^j \abs{I}, \ (3/4)^j \abs{I}]$ and we will have a complete assignment $\epsilon = \sum_{j=0}^{j_0} \eta^j \in \{-1,1\}^k$ for $j_0$ satisfying $(3/4)^{j_0} \abs{I} \le 1$, i.e. for $j_0 \le \log(1/\abs{I})/\log(3/4) = \log(\abs{I})/\log(4/3)$.
It follows from Equation~\ref{eq:rounding-one-step} and the triangle inequality that this complete assignment $\epsilon$ satisfies this bound:
\begin{equation}
\label{eq:sum-over-stages-discrepancy-bound}
\begin{aligned}
\max_{f \in \calF}\abs*{\sum_{i \in I} \epsilon_i f(X_i)}
&= \max_{f \in \calF}\abs*{\sum_{i \in I} \p*{\sum_{j=0}^{j_0} \eta_i^{j}} f(X_i)} \\
&\le  \max_{f \in \calF}\sum_{j=0}^{j_0}\abs*{\sum_{i \in I} \eta_i^{j} f(X_i)} \\
&\le \max_{f \in \calF} \sum_{j=0}^{j_0} \sum_{s\geq 1}  Q^{j}_s\norm*{P_{\calX_{I^j}}\{\pi_s(f) - \pi_{s-1}(f)\}}_2 \\
\end{aligned}
\end{equation}

We will bound this by choosing some `critical resolution' $s_\star$ and splitting our sum over $s$ into two parts.
We will break it into a high-resolution ($s > s_\star$) part where we will use the bound $Q_j^s \le 2^{s/2}$ and use the same chaining argument we would if $\varepsilon_1 \ldots \varepsilon_k$ were chosen by flipping a coin, and a low-resolution ($s \le s_\star$) part where we will take advantage of $Q_j^s$ often being substantially smaller. The parts will be separately bounded uniformly over $\F$.
$$
\begin{aligned}
& \max_{f \in \calF}\abs*{\sum_{i \in I} \epsilon_i f(X_i)} \le  T_{\le} + T_{>} \quad \text{for} \\
&T_{\le} = \max_{f \in \calF} \sum_{s \le s_\star} \sum_{j=0}^{j_0} Q^{j}_s \norm*{P_{\calX_{I^j}}\{\pi_s(f) - \pi_{s-1}(f)\}}_2 \\
&T_{>}   = \max_{f \in \calF} \sum_{s > s_\star} 2^{s/2} \sum_{j=0}^{j_0}  \norm*{P_{\calX_{I^j}}\{\pi_s(f) - \pi_{s-1}(f)\}}_2.
\end{aligned}
$$


\noindent\textbf{The Low-Resolution Part.}

On the low-resolution part, we will use the bound
$$
\begin{aligned}
\norm{P_{\calX_{I_j}}\{ \pi_s(f) - \pi_{s-1}(f) \}}_2
&\le \sqrt{\abs{I_j}}\norm{\pi_s(f) - \pi_{s-1}(f)}_{\infty} \\
&\le \sqrt{\abs{I_j}}\diam_{\infty}(\calF) \qfor
\diam_{\infty}(\calF) = \max_{g,h\in \calF}\norm{g-h}_{\infty}.
\end{aligned}
$$
Making this substitution, we can bound the low-resolution part of our discrepancy as follows.
$$
T_{\le} \le \diam_{\infty}(\calF)  \sum_{s \le s_\star} \sum_{j=0}^{j_0} Q^{j}_s \sqrt{\abs{I_j}}.
$$
What remains is to bound this sum. To simplify our handling of the `resolution weights' $Q_s^j$, we will group terms where $s_j$ takes on the same value and therefore the three cases of $Q^j_s$ occur at the same values of $s$, defining $U_\ell = \{ j  \ : \ s_j = \ell \}$. In these terms, we can rewrite our bound as follows, then bound the values of $Q_s^j$ uniformly over $U_\ell$.
\begin{equation}
  \begin{aligned}
  \sum_{s \le s_\star} \sum_{j=0}^{j_0} Q^{j}_s \sqrt{\abs{I_j}}
  &= \sum_{s=1}^{s_\star} \sum_{\ell=1}^{s_n} \sum_{j \in U_\ell} Q_s^j \sqrt{\abs{I^j}} \\
  &= \sum_{s=1}^{s_\star} \sum_{\ell=1}^{s_n} \bar Q_s^\ell \sum_{j \in U_\ell}  \sqrt{\abs{I^j}}  \\
  &\text{ where } \ \bar Q^\ell_s = c_1
  \begin{cases}
  \exp(-c_2 c_3^{-1/2} 2^{2^{\ell}}), &\text{if } s<\ell,\\
  1, &\text{if } s = \ell,\\
  2^{s/2}, &\text{if } s>\ell
  \end{cases}
  \end{aligned}
\end{equation}
Here our bound for on $Q_j^s$ for $s < \ell$ follows from the definition of $s_j$ as the largest value of $s$ (and therefore \emph{a} value of $s$) for which $c_3\abs{I^j} \ge 2^{2^{s+1}}$:
$$
\sqrt{\abs{I^j} }\ge c_3^{-1/2} \sqrt{2^{2^{s_j+1}}} = c_3^{-1/2} 2^{2^{s_j}} \ \ \text{and consequently} \ \
\exp(-c_2\sqrt{\abs{I_j}}) = Q_s^j \le \bar Q^\ell_s = \exp(-c_2 c_3^{-1/2} 2^{2^{\ell}}).
$$
Because $\abs{I_j}$ decays geometrically ($\abs{I_{j+1}} \le 3/4\abs{I_j}$ and therefore $\abs{I_{j+k}} \le (3/4)^k \abs{I_{j}}$), we can bound the sum of $\sqrt{\abs{I_j}}$ over $U_\ell$ by a multiple of its largest term. Letting $j_{\ell}$ be smallest elements of $U_\ell$, we have
$$
\begin{aligned}
\sum_{j \in U_{\ell}} \abs{I_j}^{1/2}
&= \sum_{k=0}^{\abs{U_{\ell}}-1} \abs{I_{j_\ell + k}}^{1/2} \\
&\le \sum_{k=0}^{\abs{U_{\ell}}-1} \sqrt{3/4}^{k} \sqrt{\abs{I_{j_\ell}}} \\
&\le \sqrt{\abs{I_{j_\ell}}}\sum_{k=0}^{\infty} \sqrt{3/4}^{k} = \sqrt{\abs{I_{j_\ell}}} \frac{1}{1-\sqrt{3/4}}
\end{aligned}
$$
Furthermore, as $s_j=\ell$ if and only if $c_3^{-1}2^{2^{\ell+1}} \le \abs{I_j} < c_3^{-1}2^{2^{\ell+2}}$, we have the bound $\sqrt{\abs{I_{j_\ell}}} \le \min(c_3^{-1/2} 2^{2^{\ell+1}}, \sqrt{n\abs{I}})$. It follows that
$$
\begin{aligned}
T_{\le}
&\le \diam_{\infty}(\calF) \sum_{s=1}^{s_\star} \sum_{\ell=1}^{s_k}  2^{2^{\ell+1}} \bar Q^\ell_s \\
&\le \frac{c_1 \diam_{\infty}(\calF) }{1-\sqrt{3/4}}  \sum_{s=1}^{s_\star} \sum_{\ell=1}^{s_{\abs{I}}}
 \begin{cases}
  \sqrt{\abs{I}} \exp(-c_2 c_3^{-1/2} 2^{2^\ell}) & \text{for} \ \ell > s_\star \\
  c_3^{-1/2} 2^{2^{\ell+1}} 2^{s/2}                          & \text{for} \ \ell \le s_\star \\
\end{cases}
\end{aligned}
$$
We split this into a sum over $\ell > s_\star$ and one over $\ell \le s_\star$ and bound the sums as follows. In the first, because the summand does not depend on $s$, we get $s_\star$ copies of the sum over $\ell$.
$$
\begin{aligned}
\sum_{s=1}^{s_\star} \sum_{\ell=s_\star+1}^{s_{\abs{I}}} \sqrt{\abs{I}} \exp(-c_2 c_3^{-1/2} 2^{2^\ell})
&=
\sqrt{\abs{I}} s_\star  \sum_{\ell=s_\star+1}^{s_{\abs{I}}} \exp(-c_2 c_3^{-1/2} 2^{2^\ell}) \\
&\le \sqrt{\abs{I}} s_\star \frac{\exp(-c_2 c_3^{-1/2} 2^{2^{s_\star+1}})}{1-\exp(-c_2 c_3^{-1/2} 2^{2^{s_\star+1}})} \\
\end{aligned}
$$
Here the last bound follows by bounding the sum over $\ell$
by the sum of a geometric series, observing that the rate of
decrease is at least the sum's first term:
$$
\begin{aligned}
\frac{\exp(-c_2 c_3^{-1/2} 2^{2^{x+1}})}{\exp(-c_2 c_3^{-1/2} 2^{2^{x}})}
&= \exp\left\{-c_2 c_3^{-1/2} (2^{2^{x+1}} - 2^{2^{x}})\right\} \\
&= \exp\left\{-c_2 c_3^{-1/2} 2^{2^{x}}(2^{2^{x}} - 1)\right\}  \\
&\le \exp(-c_2 c_3^{-1/2} 2^{2^{x}}).
\end{aligned}
$$

In the second, we can use the same geometric series approach to bound our sums over $s$ and then over $\ell$.
$$
\begin{aligned}
\sum_{s=1}^{s_\star} \sum_{\ell=0}^{s_\star} 2^{s/2} 2^{2^{\ell+1}}
&= \sum_{\ell=0}^{s_{\star}} 2^{2^{\ell+1}} 2^{s_\star/2}\sum_{k=0}^{s_\star-1} 2^{-k/2} \\
&\le \sum_{\ell=0}^{s_{\star}} 2^{2^{\ell+1}} 2^{s_\star/2}\sum_{k=0}^{\infty} \sqrt{2}^{-k} \\
&= \sum_{\ell=0}^{s_{\star}} 2^{2^{\ell+1}} 2^{s_\star/2}/\left(1-\sqrt{1/2}\right) \\
&\le 2^{2^{s_\star+1}} 2^{s_\star/2} \cdot 2/\left(1-\sqrt{1/2}\right).
\end{aligned}
$$
Here we've used the following bound on the rate of decrease of our series $2^{2^x}$ as we count down from $\ell=s_\star$ to $\ell=0$
to justify replacing the sum with two times its largest term.
$$
\frac{2^{2^x}}{2^{2^{x+1}}} = 2^{2^{x}(1-2)}=2^{-2^x} \le 1/2 \qfor x \ge 0.
$$

Summarizing, we have the following bound, in which $\expo(x)=e^x /(1-e^x)$
$$
\begin{aligned}
T_{\le}
&\le \frac{c_1 \diam_{\infty}(\calF) }{1-\sqrt{3/4}}
\left\{ \sqrt{\abs{I}} s_\star \expo\left(-c_2 c_3^{-1/2} 2^{2^{s_\star+1}}\right) + c_5 2^{2^{s_\star+1}} 2^{(s_\star+1)/2} \right\} \\
&\qfor c_5 = c_3^{-1/2} \cdot 2/(1-\sqrt{1/2}).
\end{aligned}
$$

\noindent\textbf{The High-Resolution Part.}
Using the Cauchy-Schwarz inequality, we can bound the sum of two-norms that appears in
\eqref{eq:sum-over-stages-discrepancy-bound}
by a weighted two-norm, which we will interpret as a multiple of a weighted empirical norm $\norm{\cdot}_{L_2^w(I)}$ with weights proportional to the number of rounds each unit went without being assigned.
$$
\begin{aligned}
\sum_{j=0}^{j_0} \norm*{P_{\calX_{I^j}} g }_2
&\le\sqrt{j_0+1}\sqrt{ \sum_{j=0}^{j_0} \norm*{P_{\calX_{I^j}} g }^2 } \\
&=\sqrt{\sum_{i \in I} \tilde w_i g(X_i)^2} \ \ \text{for} \ \  \tilde w_i =  \sum_{j=0}^{j_0} \mathbf{1}(i \in I^j) \\
&=\sqrt{nW}\sqrt{\frac{1}{\abs{I}}\sum_{i=1}^n w_i g(X_i)^2 } \ \ \text{ for } \ \
W = \frac{1}{\abs{I}}\sum_{i=1}^n \tilde w_i \ \ \text{for} \ \  w_i = \frac{\tilde w_i}{W}.
\end{aligned}
$$
Here $W$ is no larger than $4$ (and at least $1$). This follows from the geometric decrease in the number of unassigned units.
$$
W=\frac{1}{\abs{I}}\sum_{j=0}^{j_0} \abs{I_j} \le \sum_{j=0}^{j_0} (3/4)^j \le \sum_{j=0}^{\infty} (3/4)^j =  1/(1-3/4)=4.
$$
It follows that each weight is between $1/4$ and $j_0/4 \le 4\log(n)/\log(4/3)$ and therefore that this norm is equivalent to the unweighted empirical norm $\norm{g}_{L_2(I)}$ up to log factors. In particular $1/2 \le \norm{g}_w/\norm{g}_{L_2(I)} \le \sqrt{\log(\abs{I})/\log(4/3)}$. Thus,

$$
\begin{aligned}
T_{>}
&\le 2\sqrt{\abs{I}} \max_{f \in \calF} \sum_{s = s_\star+1}^{\infty} 2^{s/2}\norm*{\pi_s(f) - \pi_{s-1}(f)}_{L_2(w)} \ \text{ where } \ \frac{\norm{g}_{L_2^w(I)}} {\norm{g}_{L_2(I)}} \in \left[\frac{1}{2}, \ \sqrt{\frac{\log(\abs{I})}{\log(4/3)}} \right] \\
%&\le 2\sqrt{\frac{k\log(k)}{\log(4/3)}} \max_{f \in \calF} \sum_{s > s_\star} 2^{s/2}\norm*{\pi_s(f) - \pi_{s-1}(f)}_{L_2(P_k)}.
\end{aligned}
$$
We will bound this in terms of the \emph{entropy numbers}
$$
e_s(\calF, \norm{\cdot}_{L_2^w(I)}) := \inf_{\calF_s \subseteq \calF, \abs{\calF_s} \le 2^{2^s}} \sup_{f \in \calF} \inf_{f' \in \calF_s}\norm{f-f'}_{L_2^w(I)}.
$$
By choosing $\pi_s(f)$ to be the closest point to $f$ in the associated optimal cover $\calF_s$,
the bound $\norm*{\pi_s(f) - \pi_{s-1}(f)}_{L_2^w(I)} \le e_s(\calF, \norm{\cdot}_{L_2^w(I)})$ holds uniformly over $f \in \F$. In those terms,
$$
T_{>} \le 2\sqrt{n}\sum_{s = s_\star+1}^{\infty} 2^{s/2} e_s(\calF, \norm{\cdot}_{L_2^w(I)}).
$$
This is a discrete version of Dudley's entropy bound, from which it is a well-known exercise to derive the more common integral version via integral approximation \citep[e.g.][Remarks after Proposition 2.5.1]{talagrand2014upper}. And we can bound this in terms of a uniform entropy integral, first making the change of variables $u=t/\norm{F}_{L_2^w(I)}$ where $F$ is an envelope function for our class $\F$, then maximizing over any set of probability measures $Q$ containing $L_2^w(I)$ \citep[see e.g. the proof of][Theorem 2.5.2]{van1996weak}.
\begin{equation}
\begin{aligned}
T_{>}
&\le c_6 \sqrt{\abs{I}}\int_0^{e_{s_\star + 1}(\calF, \norm{\cdot}_{L_2^w(I)})} \sqrt{\log\calN(\calF, \norm{\cdot}_{L_2^w(I)}, t)} dt \\
&= c_6   \sqrt{\abs{I}}\norm{F}_{L_2^w(I)} \int_0^{e_{s_\star + 1}(\calF, \norm{\cdot}_{L_2^w(I)})/\norm{F}_{L_2^w(I)}} \sqrt{\log\calN(\calF, \norm{\cdot}_{L_2^w(I)}, u\norm{F}_{L_2^w(I)}}) du \\
&\le c_6 \sqrt{\abs{I}}\int_0^{e_{s_\star + 1}(\calF, \norm{\cdot}_{L_2^w(I)})/\norm{F}_{L_2^w(I)}} \sup_{Q} \sqrt{\log\calN(\calF, \norm{\cdot}_{L_2(Q)}, u\norm{F}_{L_2(Q)}}) du.
\end{aligned}
\end{equation}

We conclude the section by replacing our integral's upper limit with a deterministic bound. Because bounds on entropy numbers follow from bounds on norms, i.e. $e_s(\F, \norm{\cdot}_a) \le \min\{c e_s(\F, \norm{\cdot}_b), \  d\} $ if $\norm{f}_a \le \min\{c\norm{f}_b, \ d\}$ for all $f \in \F$, it suffices to find a deterministic upper bound on $\norm{\cdot}_{L_2^w(I)}$. To  do this, we combine our aforementioned bound $\norm{\cdot}_{L_2^w(I)} \le \sqrt{\log(\abs{I})/\log(4/3)}\norm{\cdot}_{L_2(I)}$ with (i)
the bound $\norm{\cdot}_{L_2(I)} = \frac{1}{\sqrt{\abs{I}}}\norm{\P_{\calX_I} \cdot}_2 \le \frac{1}{\sqrt{\abs{I}}}\norm{\P_{\calX} \cdot}_2 = \sqrt{n/\abs{I}}\norm{\cdot}_{L_2(\{1 \ldots n\})}$ and (ii) Lemma~\ref{lem:empirical-and-population-norm-equivalence}'s bound on the ratio $\norm{\cdot}_{L_2(\{1 \ldots n\})}/\norm{\cdot}_{L_2(P)}$. Together,  these imply that on an event of probability $1-\exp(-nr^2/8b^2)$,
$$
\begin{aligned}
&\norm{f}_{L_2^w(I)} \le 2\sqrt{\log(\abs{I})/\log(4/3)}\max\{\norm{f}_{L_2(P)}, 2r\} \ \text{for all $f \in \F$ and therefore} \\
&e_s(\F, \norm{\cdot}_{L_2(w)}) \le 2\sqrt{\log(\abs{I})/\log(4/3)}\max\{ e_s(\F, \norm{\cdot}_{L_2(P)}), 2r\} \ \ \text{for all $s$.}
\end{aligned}
$$
Making this substitution in our bound above for $s=s_\star+1$ implies that
$$
\begin{aligned}
&T_{>} \le c_6 \sqrt{n}\int_0^{2\sqrt{\frac{\log(\abs{I})}{\log(4/3)}}\max(\delta,2r)}  \sup_{Q} \sqrt{\log\calN(\calF, \norm{\cdot}_{L_2(Q)}, u\norm{F}_{L_2(Q)}}) du \\
&\qfor \delta \ge e_{s_\star+1}(\F, \norm{\cdot}_{L_2(P)}) \ \ \text{i.e.} \ \  \delta \ \text{ satisfying }  \ \log_2 \calN(\calF, \norm{\cdot}_{L_2(P)}, \delta) \le 2^{s_\star+1}.
\end{aligned}
$$

Summing our bounds on $T_\le$ and $T_>$ and dividing by $\sqrt{n}$ gives us the claimed bound. As our only probabilistic bound is on a quantity
that does not involve $I$, our bound holds simultaneously for all $I \subseteq \{ 1 \ldots n \}$ on the aforementioned event.

To simplify our statement of the bound, we replace the infinity norm diameter with a multiple ($2$) of our infinity norm bound $b$,
rename $s_\star$ to $s$, and simplify and rename our universal constants.
\end{proof}


\begin{proof}[Proof of Corollary~\ref{cor:hdisc-negligible}]
Substituting the upper bound $n$ for $\abs{I}$ in the discrepancy bounds Lemma~\ref{lemm:bdd_ent_class} gives us this finite sample bound on hereditary discrepancy. Letting $F$ be any envelope function for $\F$, in terms of the uniform entropy integral
$$
J(x,\F,F) := \int_0^{x} \sup_{Q} \sqrt{\log\calN(\F, \norm{\cdot}_{L_2(Q)}, u\norm{F}_{L_2(Q)}}) du,
$$
we have the bound
$$
\begin{aligned}
\frac{1}{\sqrt{n}}\hdisc(\P_{\calX} \F) &\le c_1b\left\{ s \expo\left(-c_2 2^{2^{s+1}}\right) + n^{-1/2} 2^{2^{s+1}} 2^{(s+1)/2} \right\} \\
&+c_3 J\left(2\sqrt{\frac{\log(n)}{\log(4/3)}}\max(\delta,2r), \F, F\right) \\
&\qfor
\delta \ \text{ satisfying }  \ \log_2 \calN(\calF, \norm{\cdot}_{L_2(P)}, \delta) \le 2^{s+1} \qand b \ge \max_{f \in \F}\norm{f}_\infty.
\end{aligned}
$$
This bound holds for all $s$ with probability $1-\exp(-nr^2/8b^2)$ for any $r$ satisfying the inequality $r^2 \ge 2b^3 R_n( \{f \in \F \ : \ \norm{f}_{L_2(P)} \le r \})$. To establish that $\hdisc(\P_{\calX} \F)$ is $o_p(n^{-1/2})$
we will show that we can choose $s$ and $r$ as functions of $n$ so that bound above is $o(1)$ and holds with probability tending to one.

We can take $r$ to be on the order of $n^{-1/4}$. A standard chaining argument \citet[see e.g.][Theorem 2.1.41]{van1996weak}), gives us a bound on this Rademacher complexity in terms of our uniform entropy integral: for some universal constant $c$,
$$
R_n(\{f \in \F : P f^2 \le r^2\}) \le c n^{-1/2} \norm{F}_{L_2(P)} J(1,\F, F).
$$
Thus, our inequality is satisfied for $r^2 = cn^{-1/2} \norm{F}_{L_2(P)} J(1,\F, F) / (2b^3)$, i.e. for $r = c_{\F} n^{-1/4}$ for some constant depending on our class $\F$.

To choose $s$, observe that because our uniform entropy integral is bounded,
it must be the case that it vanishes when we integrate over arbitrarily small neighborhoods of zero, i.e. that $J(\delta', \calF, F) \to 0$ as $\delta' \to 0$. Thus, the integral term in our bound tends to zero if its upper limit $\delta'=\sqrt{\frac{\log(n)}{\log(4/3)}}\max(\delta,2r) \to 0$ which, given our characterization of $r$, happens if $\sqrt{\log(n)}\delta \to 0$. Considering the other terms as well, our bound will tend to zero if the following conditions on $s$ are satisfied for some sequence $s_n$.
\begin{enumerate}
  \item $s_n \expo(-c_2 2^{2^{s_n+1}})$ is $o(1)$.
  \item $2^{2^{s_n+1}} 2^{(s_n+1)/2} = o(\sqrt{n})$.
  \item There exists $\delta_n = \delta(s_n)$ satisfying
  $\sqrt{\log(n)}\delta_n \to 0$ and $\log_2 \calN(\calF, \norm{\cdot}_{L_2(P)}, \delta) \le 2^{s_n+1}$.
\end{enumerate}

Let $\rho$ be any number in $(0,1/2)$.
We will choose $s_n$ to be the integer $s$ for which $2^{2^{s+1}} \le n^{\rho} < 2^{2^{s+2}}$. With this choice, condition (1) holds because it does for any sequence $s_n$ tending to $\infty$. And
condition (2) holds by almost by construction, as $2^{(s_n+1)/2} / n^{1/2-\rho} \to 0$ as $s \to \infty$. This leaves condition (3).

Because $\F$ has bounded uniform entropy integral, we know that the integrand,
$\sup_{Q}$ \ $\sqrt{\log\calN(\calF, \norm{\cdot}_{L_2(Q)}, u\norm{F}_{L_2(Q)}})$, is $o(1/u)$ as $u \to 0$; if it were not, our integral would diverge. This means that for any constant $c$, we can choose $\delta$ small enough that $\log\calN(\calF, \norm{\cdot}_{L_2(P)}, \delta) < c/\delta^2$ and therefore $\log\calN(\calF, \norm{\cdot}_{L_2(P)}, \delta) < 2^{s_n+1}$ if
$c < \delta^2 2^{s_n+1}$. Having chosen $s_n$ so that $n^{\rho/2} \le 2^{s_n+1}$, it follows that $\delta_n = c^{1/2}n^{-\rho'/4}$ for $\rho' < \rho$ works and, in addition, satisfies $\sqrt{\log(n)}\delta_n \to 0$, so condition (3) holds as well.
\end{proof}


\begin{proof}[Proof of Lemma~\ref{lem:empirical-and-population-norm-equivalence}]
We will use \citet[Theorem 3.3]{bartlett2005local} which, applied to the set of functions $\F^2=\{f^2 : f \in \F \}$ with $K \downarrow 1$ and $B=b^2$, gives the probability  $1-e^x$ bound
$$
\Pn f^2 \le 2\P f^2 + 6b^{-2} s^2 + 16b^2x/n \qif s^2 \ge b^2R_n\left(\{f^2  \ : \ f \in \F : \P f^4 \le s^2\}\right).
$$
We will use the boundedness of $\F$ to simplify our characterization of $r$. It implies that $f^4 \le b^2f^2$
and consequently that $\{f \in \F : \P f^4 \le s^2\} \subseteq \{f \in \F : \P f^2 \le b^2s^2\}$ and
therefore, because $\phi(x)=x^2$ has Lipschitz constant $2b$ on $x \in [0,b]$, via the contraction principle for Rademacher processes \citep[e.g.][Theorem A.6]{bartlett2005local}, that
$$
\begin{aligned}
b^2 R_n \left( \{f^2 : f \in \F, \ \Pn f^4 \le s^2\} \right)
&\le b^2 R_n \left(\{f^2 : f \in \F, \ \Pn f^2 \le b^2 s^2\}\right) \\
&\le b^2 \cdot 2b R_n \left(\{f \in \F, \ \Pn f^2 \le b^2 s^2\}\right).
\end{aligned}
$$
Thus, letting $r=b s$, we have
$$
\Pn f^2 \le 2\P f^2 + 6r^2+ 16b^2x/n \qif 2b^3 R_n\left( \{ f : \Pn f^2 \le r^2\} \right) \le r^2.
$$
Taking $x=nr^2/8b^2$ gives us the probability $1-\exp(-nr^2/8b^2)$ bound
$$
\Pn f^2 \le 2\P f^2 + 8 r^2 \le \max(4\P f^2, 16b^{-2}r^2)
$$
and taking square roots gives the bound we claimed.
\end{proof}


\subsection{Rounding and Efficiency}
\label{sec:rounding-and-efficiency}

Here we analyze the variance cost of rounding in more detail. Let $v_0(x) = \Var(Y_i \mid W_i=0, X_i=x)$ denote the conditional variance among controls, and let $\hat r(x) = \hat\theta(x) - \floor{\hat\theta(x)}$ denote the fractional part of the weight. Define the rounding error functional
$$
R(f) = \frac{1}{N_1}\sum_{i=1}^n \overset{\neg}{W}_i \{\hat\theta(X_i) - \round{\hat\theta}(X_i)\} f(X_i).
$$
The excess variance from rounding can be decomposed as follows.
\begin{equation}
\label{eq:rounding-variance-excess}
\begin{aligned}
N_1 \times (V^{\round{\hat\theta}} - V^{\hat\theta})
&= \frac{1}{N_1}\sum_{i=1}^n \overset{\neg}{W}_i \{\round{\hat\theta}_i - \hat\theta(X_i)\}^2 v_0(X_i) + 2R(\hat\theta v_0) \\
&= \frac{1}{N_1}\sum_{i=1}^n \overset{\neg}{W}_i \hat r(X_i)\{1 - \hat r(X_i)\} v_0(X_i) + R\{(1-2\hat r)v_0\} + 2R(\hat\theta v_0).
\end{aligned}
\end{equation}
The first term is the variance from independent Bernoulli rounding. The remaining terms involve the rounding error functional $R(\cdot)$ applied to functions that may not be in our model $\model$.

\begin{proposition}
\label{prop:rounding-efficiency}
Suppose $\model$ has bounded uniform entropy integral and $\roundingerror_{\model,\hat\theta}(\round{\hat\theta}) = o_p(n^{-1/2})$. If $\norm{\hat\theta v_0}_\model = O_p(n^{1/2})$ and $\norm{(1-2\hat r)v_0}_\model = O_p(n^{1/2})$, then $R(\hat\theta v_0) = o_p(1)$ and $R\{(1-2\hat r)v_0\} = o_p(1)$.
\end{proposition}

\begin{proof}
By Proposition~\ref{prop:rounding}, the rounding error is $o_p(n^{-1/2})$ uniformly over the unit ball of $\model$. For $f$ with $\norm{f}_\model = O_p(n^{1/2})$, we can write $f = \norm{f}_\model \cdot (f/\norm{f}_\model)$ where $f/\norm{f}_\model$ is in the unit ball. Because the uniform bound holds for all elements of the unit ball simultaneously, it applies to the random element $f/\norm{f}_\model$. Thus
$$
R(f) = \norm{f}_\model \cdot R(f/\norm{f}_\model) = O_p(n^{1/2}) \cdot o_p(n^{-1/2}) = o_p(1).
$$
\end{proof}

Like our analysis of the efficiency of weighting in Section~\ref{sec:asymptotic}, the conditions of Proposition~\ref{prop:rounding-efficiency} do not require the functions $\hat\theta v_0$ and $(1-2\hat r) v_0$ to be in $\model$ itself---only that they can be approximated by sequences $m_n \in \model$ with $\norm{m_n}_\model = O_p(n^{1/2})$. For nonparametric models, this holds whenever these functions are square integrable.

For RKHS models, we can verify these conditions using the results of \citet{hirshberg2019minimax}.

\begin{corollary}
\label{cor:rounding-efficiency-rkhs}
Let $\model$ be the unit ball of an RKHS $\calH$ that is dense in $L^2(P)$, and let $\hat\theta$ be derived from the minimax linear weights of \citet{hirshberg2019minimax}. If $\hat\theta v_0$ and $(1-2\hat r)v_0$ are in $L^2(P)$, then $R(\hat\theta v_0) = o_p(1)$ and $R\{(1-2\hat r)v_0\} = o_p(1)$.
\end{corollary}

\begin{proof}
The proof has two parts. First, we use results from \citet{hirshberg2019minimax} to establish that the rounding error $R(\cdot)$ is $o_p(n^{-1/2})$ uniformly over the unit ball of $\calH$. Second, we show that this implies $R(f) = o_p(1)$ for any $f \in L^2(P)$.

\paragraph{Part 1: Rounding error over the model.}
By \citet[Theorem 2]{hirshberg2019minimax}, the minimax linear weights $\hat\gamma$ achieve $o_p(n^{-1/2})$ imbalance uniformly over $\model$. Combined with Proposition~\ref{prop:rounding}, which shows the rounding scheme \eqref{eq:round} preserves this rate, we have
\[
\sup_{m \in \model} |R(m)| = o_p(n^{-1/2}).
\]

\paragraph{Part 2: Extension to $L^2(P)$.}
We now show $R(f) = o_p(1)$ for any $f \in L^2(P)$. The idea is to approximate $f$ by elements of $\calH$ whose norms grow slowly enough that we can apply Part 1.

Because $\calH$ is dense in $L^2(P)$, there exists a sequence $g_1, g_2, \ldots$ in $\calH$ with $\norm{g_k - f}_{L^2(P)} \to 0$. We construct a sequence $f_n \in \calH$ with $\norm{f_n - f}_{L^2(P)} \to 0$ and $\norm{f_n}_\calH = O(n^{1/2})$ by setting $k(n) = \max\{k : \norm{g_k}_\calH \le n^{1/2}\}$ and $f_n = g_{k(n)}$. Because each $g_k$ has finite RKHS norm, $k(n) \to \infty$ as $n \to \infty$, so $\norm{f_n - f}_{L^2(P)} = \norm{g_{k(n)} - f}_{L^2(P)} \to 0$.

Now we bound $R(f)$ by the triangle inequality:
\[
|R(f)| \le |R(f_n)| + |R(f - f_n)|.
\]
For the first term, because $f_n / \norm{f_n}_\calH$ is in the unit ball of $\calH$, Part 1 gives
\[
|R(f_n)| = \norm{f_n}_\calH \cdot |R(f_n / \norm{f_n}_\calH)| = O(n^{1/2}) \cdot o_p(n^{-1/2}) = o_p(1).
\]
For the second term, we use the fact that $|\hat\theta - \round{\hat\theta}| \le 1$:
\[
|R(f - f_n)| \le \frac{1}{N_1}\sum_{i=1}^n \overset{\neg}{W}_i |f(X_i) - f_n(X_i)|.
\]
This is a weighted empirical average of $|f - f_n|$. By the law of large numbers, it converges in probability to $\E[|f(X) - f_n(X)|]$, which is bounded by $\norm{f - f_n}_{L^2(P)} \to 0$.

Combining these bounds, $R(f) = o_p(1)$.
\end{proof}

The main term in \eqref{eq:rounding-variance-excess} is an average of $\hat r(1-\hat r) v_0$. Combined with the $\hat\theta^2 v_0$ term from the variance of the unrounded estimator, the total variance is proportional to $\{\hat\theta^2 + \hat r(1-\hat r)\} v_0$. When weights are close to binary ($\hat\theta \approx \hat r$), this simplifies to $\hat r v_0$---larger than $\hat r^2 v_0$ by a factor of $1/\hat r$, but we only need outcomes from $N_1 \approx n\E[\hat r]$ individuals rather than all $n$ controls.

\begin{proof}[Proof of Corollary~\ref{cor:rounding}]
By Proposition~\ref{prop:rounding}, the rounding scheme \eqref{eq:round} satisfies $\roundingerror_{\model,\htheta}(\round{\htheta}) = o_p(n^{-1/2})$, so the triangle inequality gives $\imbalance_{h,\model}(\round{\htheta}) \le \imbalance_{h,\model}(\htheta) + o_p(n^{-1/2}) = o_p(n^{-1/2})$.
The matched estimator $\hat\psi^{match}$ is an IPW estimator with weights $\round{\htheta}$ in place of $\htheta$. We verify the conditions of Proposition~\ref{proposition:clt-general}$'$ for these rounded weights.

The normalization condition $N_1^{-1}\sum_i \overset{\neg}{W}_i \round{\htheta}_i = 1$ is preserved by the balance-preserving rounding scheme \eqref{eq:round}. For infinitesimality, observe that rounding changes each weight by at most one, so $\round{\htheta}_i^2 \le 2\htheta_i^2 + 2$. In particular, $\max_i \overset{\neg}{W}_i \round{\htheta}_i^2 \le 2\max_i \overset{\neg}{W}_i \htheta_i^2 + 2$. For the denominator, because $\round{\htheta}_i$ are non-negative integers, $\round{\htheta}_i^2 \ge \round{\htheta}_i$ for all $i$, and so $n\Gamma_{\mathrm{round}}^2 = n N_1^{-1}\sum_i \overset{\neg}{W}_i \round{\htheta}_i^2 \ge n N_1^{-1}\sum_i \overset{\neg}{W}_i \round{\htheta}_i = n$. Combining: $\max_i \overset{\neg}{W}_i \round{\htheta}_i^2 \big/ n\Gamma_{\mathrm{round}}^2 \le (2\max_i \overset{\neg}{W}_i \htheta_i^2 + 2)/n \to_p 0$, using infinitesimality of $\htheta$ and $n\Gamma_{\mathrm{round}}^2 \ge n$.

For weighted $L_2$ consistency, use $\round{\htheta}_i^2 \le 2\htheta_i^2 + 2$ to write
\[
\sum_i \overset{\neg}{W}_i \round{\htheta}_i^2 (\hat\mu_i - m_i)^2 \le 2\sum_i \overset{\neg}{W}_i \htheta_i^2 (\hat\mu_i - m_i)^2 + 2\sum_i \overset{\neg}{W}_i (\hat\mu_i - m_i)^2.
\]
The first term is $o_p(N_1 \Gamma^2)$ by the weighted $L_2$ consistency of $\hat\mu$ with $\htheta$; the second is $o_p(N_1\Gamma^2)$ by $L_2$ consistency. The asymptotic linearity \eqref{eq:asymptotic-linearity-rounded} and central limit theorem \eqref{eq:asymptotic-normality-rounded} then follow from Proposition~\ref{proposition:clt-general}$'$.
\end{proof}

\begin{proof}[Proof of Corollary~\ref{cor:integer-sbw}]
The non-integer version of \eqref{eq:integer-sbw} achieves $o_p(n^{-1/2})$ imbalance for RKHS models with bounded continuous kernel \citep[Theorem 2]{hirshberg2019minimax}.
By Proposition~\ref{prop:rounding}, rounding these weights to integers via \eqref{eq:round} incurs an additional imbalance of $o_p(n^{-1/2})$, so the rounded weights are feasible for \eqref{eq:integer-sbw} with budget $B_n$ equal to the sum of these two rates (see below for explicit rate calculations). Feasibility with probability tending to one follows.

Let $\round{\htheta}^*$ denote the solution to \eqref{eq:integer-sbw}. Because the integer program minimizes $N_1^{-1}\sum_i \overset{\neg}{W}_i \theta_i^2$ over feasible integer solutions, and the rounded continuous solution is feasible, the optimal value satisfies $N_1^{-1}\sum_i \overset{\neg}{W}_i (\round{\htheta}^*_i)^2 \le N_1^{-1}\sum_i \overset{\neg}{W}_i \round{\htheta}_i^2 = O_p(1)$. Moreover, for a bounded continuous kernel the continuous IPW weights are bounded, so $\max_i \overset{\neg}{W}_i \round{\htheta}_i = O_p(1)$; any integer solution with a component exceeding this bound would have larger objective than the rounded continuous solution and hence be suboptimal, so $\max_i \overset{\neg}{W}_i \round{\htheta}^*_i = O_p(1)$. Since the weights $\round{\htheta}^*_i$ are non-negative integers summing to $N_1$, we have $n\Gamma^{*2} = n N_1^{-1}\sum_i \overset{\neg}{W}_i (\round{\htheta}^*_i)^2 \ge n$, and therefore $\max_i \overset{\neg}{W}_i (\round{\htheta}^*_i)^2 \big/ n\Gamma^{*2} = O_p(1/n) \to 0$: the infinitesimality condition of Corollary~\ref{cor:rounding} holds.

The imbalance condition $\imbalance_{h,\model}(\round{\htheta}^*) \le B_n = o_p(n^{-1/2})$ holds by the feasibility constraint. All remaining conditions of Corollary~\ref{cor:rounding} are inherited from the hypotheses, so the asymptotic expansions \eqref{eq:asymptotic-linearity-rounded} and \eqref{eq:asymptotic-normality-rounded} follow.
\end{proof}

\paragraph{Choosing the budget $B_n$: explicit rates for Sobolev models.}
We illustrate with Sobolev models with smoothness $s > d/2$ in dimension $d$. The budget $B_n$ is the sum of two rates: the imbalance rate from \citet[Theorem 2]{hirshberg2019minimax} and the rounding error rate from Corollary~\ref{cor:hdisc-negligible}.

For the imbalance rate, the key quantity is the effective dimension $\sum_{j=1}^\infty \min(\mu_j, n^{-1})$, where $\mu_j$ are the eigenvalues of the kernel. For Sobolev spaces, these decay as $\mu_j \asymp j^{-2s/d}$ \citep[see, e.g.,][]{kuhn2014approximation}. The effective dimension is therefore
\[
\sum_{j=1}^\infty \min(\mu_j, n^{-1}) \asymp \sum_{j \le n^{d/(2s)}} n^{-1} + \sum_{j > n^{d/(2s)}} j^{-2s/d} \asymp n^{d/(2s) - 1}.
\]
Substituting into Theorem 2 of \citet{hirshberg2019minimax} with $\kappa_m = 1$ and $\sigma^2 \asymp 1$ gives an imbalance rate of $n^{-1/2} \cdot n^{(d/(2s) - 1)/2} = n^{d/(4s) - 1}$.

For the rounding error rate, we use the covering number $\log \mathcal{N}(\model, \norm{\cdot}_{L_2}, \epsilon) \asymp \epsilon^{-d/s}$ for the Sobolev unit ball \citep[see, e.g.,][Section 2.1]{kuhn2015approximation}. The entropy integral in Lemma~\ref{lemm:bdd_ent_class} then gives a hereditary discrepancy bound of order $n^{d/(4s)}$, and the rounding error is $n^{-1} \cdot n^{d/(4s)} = n^{d/(4s) - 1}$.

Both rates match, giving $B_n \asymp n^{d/(4s) - 1}$. This is $o(n^{-1/2})$ precisely when $s > d/2$, which is the standard condition for bounded uniform entropy integral.

\paragraph{Binary search in practice.}
In practice, one can find a valid budget by binary search without computing this rate explicitly. Start at $B_n = n^{-1/2}$ and halve repeatedly until the integer program \eqref{eq:integer-sbw} becomes infeasible, then back off one step. Because feasibility is monotone in $B_n$, this terminates in $O(\log n)$ steps and yields a budget within a factor of two of the tightest feasible value. When each optimization is costly---as it can be in large samples, integer programming being hard---the explicit rate $n^{d/(4s)-1}$ provides a better starting point than $n^{-1/2}$.

If you find yourself needing to \emph{double} from $n^{-1/2}$ to achieve feasibility, rather than halve, that is a sign of trouble: the imbalance you can achieve is worse than $n^{-1/2}$, likely indicating poor overlap.

The probabilistic validity of this procedure follows from a union bound over the $O(\log n)$ candidate budgets. The imbalance bound from \citet{hirshberg2019minimax} has exponential concentration, so each candidate budget fails with probability at most $\exp(-cn)$ for some $c > 0$. The hereditary discrepancy bound requires only a norm equivalence that holds uniformly over all budgets, contributing a single event to the union. The total failure probability remains exponentially small.

\paragraph{Matching vs.\ IPW on a subsample.}
Here we prove the variance comparison claimed in Remark~\ref{rem:matching-variance}. Let $\mu_k = \E[\theta(X)^k \mid W=0]$ for $k=1,2$.

\begin{proposition}
\label{prop:matching-efficiency}
Suppose $e(X) < 1/2$ almost surely, $\nu(w,x) = v_w$ is constant in $x$ for each $w\in\{0,1\}$, and $N_1/N_0 \to 0$. Let $\hat\psi^{\mathrm{match}}$ be the integer-SBW matching estimator from Corollary~\ref{cor:integer-sbw}, and let $\hat\psi^{\mathrm{IPW}}$ be an IPW estimator using $N_1$ randomly selected controls with weights $\theta(X_i)/\bar\theta_n$, where $\bar\theta_n = N_1^{-1}\sum_{i \in S_{\mathrm{IPW}}}\theta(X_i)$. Then
\[
N_1 V^{\mathrm{match}} \to_p v_1 + v_0, \qquad N_1 V^{\mathrm{IPW}} \to_p v_1 + v_0\,\frac{\mu_2}{\mu_1^2},
\]
and
\[
\frac{V^{\mathrm{match}}}{V^{\mathrm{IPW}}} \to_p \frac{\rho + 1}{\rho + \mu_2/\mu_1^2}, \qquad \rho := \frac{v_1}{v_0}.
\]
In particular, $V^{\mathrm{match}} \le V^{\mathrm{IPW}}$, with equality if and only if $\theta(X)$ is $P(\cdot\mid W=0)$-a.s.\ constant.
\end{proposition}

\begin{proof}
\textbf{Matching estimator.}
By Corollary~\ref{cor:integer-sbw}, $\hat\psi^{\mathrm{match}}$ is asymptotically linear with influence function $N_1^{-1}(W_i + (1-W_i)\lfloor\hat\theta\rceil^*_i)\{Y_i - m(W_i,X_i)\}$. Since $e(X)<1/2$, the rounded weights lie in $\{0,1\}$ with $|S_0| = N_1$ matched controls (by normalization), and the asymptotic variance is
\[
N_1^2 V^{\mathrm{match}} = \sum_{W_i=1} v_1 + \sum_{i\in S_0} v_0 = N_1(v_1 + v_0).
\]

\textbf{IPW subsample estimator.}
The IPW estimator on a random subsample of $N_1$ controls with normalized propensity weights is asymptotically linear with influence function $N_1^{-1}(W_i - (1-W_i)(\theta(X_i)/\bar\theta_n)\mathbf{1}\{i\in S_\mathrm{IPW}\})\{Y_i - m(W_i,X_i)\}$. The asymptotic variance is
\[
N_1^2 V^{\mathrm{IPW}} = \sum_{W_i=1} v_1 + \sum_{i\in S_\mathrm{IPW}}\frac{\theta(X_i)^2}{\bar\theta_n^2}\,v_0.
\]
Since the selected controls are i.i.d.\ from the control population, $N_1^{-1}\sum_{i\in S_\mathrm{IPW}}\theta(X_i)^2 \to_p \mu_2$ and $\bar\theta_n \to_p \mu_1$ by the law of large numbers. Hence $N_1 V^{\mathrm{IPW}} \to_p v_1 + v_0\mu_2/\mu_1^2$.

\textbf{Comparison.}
Dividing and applying Slutsky's theorem gives the stated ratio. Since $x\mapsto x^2$ is convex, Jensen's inequality gives $\mu_2 = \E[\theta^2\mid W=0] \ge \E[\theta\mid W=0]^2 = \mu_1^2$, so $\mu_2/\mu_1^2\ge 1 \ge 1$ and $\rho+\mu_2/\mu_1^2 \ge \rho+1$, hence $V^{\mathrm{match}}\le V^{\mathrm{IPW}}$. Equality holds iff $\mu_2=\mu_1^2$, i.e., $\mathrm{Var}(\theta\mid W=0)=0$.
\end{proof}

For $e(X)\ll 1/2$, $\theta(X)\approx e(X)$ and $\mu_k\approx \E[e(X)^k\mid W=0]$, so $\mu_2/\mu_1^2 \approx 1 + \mathrm{Var}(e\mid W=0)/\E[e\mid W=0]^2$. With $\rho=1$ (equal noise), the ratio simplifies to $1/\{1+\mathrm{Var}(e)/(2\E[e]^2)\}$, the expression in Remark~\ref{rem:matching-variance}. The general formula shows the gain diminishes as $\rho\to\infty$ (treatment noise dominates) and strengthens as $\rho\to 0$.

\paragraph{Efficiency of integer-weight estimators.}

The variance formulas~\eqref{eq:rounding-variance-deterministic} and~\eqref{eq:rounding-variance-excess} characterize the cost of rounding in terms of the fractional parts $\hat r_i = \hat\theta_i - \floor{\hat\theta_i}$. We now ask: is this cost intrinsic to integer weights, or an artifact of rounding? The answer is that for \emph{regular} estimators---those balancing all $L^2$ functions---the variance formula is exact, not a bound. Any regular integer-weight estimator achieves the same asymptotic variance, regardless of how the weights are constructed.

Write $r(x) = \theta(x) - \floor{\theta(x)}$ for the fractional part of the population propensity odds $\theta(x) = e(x)/\{1-e(x)\}$.

\begin{definition}
\label{def:regular-matching}
A matching estimator with integer weights $\hat\theta_i \in \mathbb{Z}$ is \emph{regular} if it achieves $o_p(1)$ imbalance in every $f \in L^2(P)$:
\[
\frac{1}{N_1}\sum_{i=1}^n \overset{\neg}{W}_i \hat\theta_i\,f(X_i) - \frac{1}{N_1}\sum_{i=1}^n W_i f(X_i) = o_p(1) \quad \text{for all } f \in L^2(P).
\]
\end{definition}

\begin{propositionprime}{\ref{prop:matching-cost}}
Suppose $v_0(x) = \Var[Y \mid W=0, X=x]$ is bounded and $\theta v_0, (1{-}2r)v_0 \in L^2(P)$. Any regular matching estimator satisfies
\begin{equation}
\label{eq:integer-variance}
\Var(\hat\psi^{\textup{match}}) \ge V^{\textup{eff}} + \frac{1}{N_1^2}\sum_{i=1}^n \overset{\neg}{W}_i\, r(X_i)\{1 - r(X_i)\}\,v_0(X_i) + o_p(N_1^{-1}).
\end{equation}
The bound is sharp when each $\hat\theta_i \in \{\floor{\theta(X_i)}, \ceil{\theta(X_i)}\}$.
\end{propositionprime}

\begin{proof}[Proof of Proposition~\ref{prop:matching-cost}$'$]
The proof uses the $a^2 - b^2$ decomposition twice: first to separate the variance excess into a quadratic term and a cross term, then to decompose the quadratic term into the Bernoulli cost and another cross term. Both cross terms vanish by regularity.

\paragraph{Step 1: Variance excess decomposition.}
The treated-side contributions to $V^{\hat\theta}$ and $V^{\textup{eff}}$ are identical, so
\[
N_1^2(V^{\hat\theta} - V^{\textup{eff}}) = \sum_{i=1}^n \overset{\neg}{W}_i\bigl(\hat\theta_i^2 - \theta(X_i)^2\bigr)\,v_0(X_i).
\]
Apply $a^2 - b^2 = (a{-}b)^2 + 2(a{-}b)b$ with $a = \hat\theta_i$ and $b = \theta(X_i)$:
\begin{equation}
\label{eq:variance-decomp-1}
= \sum_{i=1}^n \overset{\neg}{W}_i(\hat\theta_i - \theta_i)^2\,v_0(X_i) + 2\sum_{i=1}^n \overset{\neg}{W}_i(\hat\theta_i - \theta_i)\,\theta_i\,v_0(X_i).
\end{equation}

\paragraph{Step 2: First cross term vanishes.}
We show the second sum in~\eqref{eq:variance-decomp-1} is $o_p(N_1)$. Set $f = \theta v_0$, which is in $L^2(P)$ by assumption. The sum equals $2N_1$ times
\begin{equation}
\label{eq:cross-term-1}
\frac{1}{N_1}\sum_{i=1}^n \overset{\neg}{W}_i(\hat\theta_i - \theta_i)\,f(X_i) = \underbrace{\frac{1}{N_1}\sum_i \overset{\neg}{W}_i\hat\theta_i\,f_i - \frac{1}{N_1}\sum_i W_i f_i}_{o_p(1) \text{ by regularity}} - \underbrace{\frac{1}{N_1}\sum_i \overset{\neg}{W}_i\theta_i\,f_i - \frac{1}{N_1}\sum_i W_i f_i}_{\to_p\, 0 \text{ by LLN}}.
\end{equation}
For the second term, we compute both averages. The control average is
\[
\frac{1}{N_1}\sum_i \overset{\neg}{W}_i\theta_i f_i = \frac{N_0}{N_1} \cdot \frac{1}{N_0}\sum_{W_i=0}\theta(X_i)^2 v_0(X_i) \to_p \frac{1{-}p}{p} \E[\theta^2 v_0 \mid W{=}0].
\]
Using $\E[g \mid W{=}0] = \E[g(1{-}e)]/(1{-}p)$ and $\theta^2(1{-}e) = e\theta$:
\[
\frac{1{-}p}{p}\cdot\frac{\E[\theta^2 v_0(1{-}e)]}{1{-}p} = \frac{\E[e\theta v_0]}{p} = \E[\theta v_0 \mid W{=}1].
\]
The treated average $N_1^{-1}\sum_i W_i f_i = N_1^{-1}\sum_{W_i=1} \theta(X_i) v_0(X_i) \to_p \E[\theta v_0 \mid W{=}1]$ by the law of large numbers. Both limits match, so the second term in~\eqref{eq:cross-term-1} vanishes.

Hence from~\eqref{eq:variance-decomp-1}:
\begin{equation}
\label{eq:variance-quadratic}
N_1(V^{\hat\theta} - V^{\textup{eff}}) = \frac{1}{N_1}\sum_{i=1}^n \overset{\neg}{W}_i(\hat\theta_i - \theta_i)^2\,v_0(X_i) + o_p(1).
\end{equation}

\paragraph{Step 3: Algebraic identity for integer weights.}
For integer $\hat\theta_i$ and $\theta_i = \floor{\theta_i} + r_i$, the squared error decomposes as
\begin{equation}
\label{eq:integer-identity}
(\hat\theta_i - \theta_i)^2 = r_i(1{-}r_i) + (\hat\theta_i - \theta_i)(1{-}2r_i) + \delta_i(\delta_i - 1),
\end{equation}
where $\delta_i = \hat\theta_i - \floor{\theta_i}$ is the integer offset. The remainder $\delta_i(\delta_i - 1) \ge 0$, with equality iff $\delta_i \in \{0,1\}$, i.e., $\hat\theta_i \in \{\floor{\theta_i}, \ceil{\theta_i}\}$.

To verify: $\hat\theta_i - \theta_i = \delta_i - r_i$, so expanding the right-hand side:
\[
r_i(1{-}r_i) + (\delta_i - r_i)(1{-}2r_i) + \delta_i(\delta_i - 1) = r_i - r_i^2 + \delta_i - 2\delta_i r_i - r_i + 2r_i^2 + \delta_i^2 - \delta_i = \delta_i^2 - 2\delta_i r_i + r_i^2 = (\delta_i - r_i)^2.
\]

\paragraph{Step 4: Second cross term vanishes.}
Summing~\eqref{eq:integer-identity} weighted by $v_0$:
\[
\sum_i \overset{\neg}{W}_i(\hat\theta_i - \theta_i)^2 v_0 = \sum_i \overset{\neg}{W}_i r_i(1{-}r_i) v_0 + \sum_i \overset{\neg}{W}_i(\hat\theta_i - \theta_i)(1{-}2r_i) v_0 + \sum_i \overset{\neg}{W}_i \delta_i(\delta_i - 1) v_0.
\]
The second sum is $N_1$ times the imbalance of $\hat\theta$ in $(1{-}2r)v_0 \in L^2(P)$, which is $o_p(1)$ by regularity. The third sum is non-negative.

Combining with~\eqref{eq:variance-quadratic}:
\[
N_1(V^{\hat\theta} - V^{\textup{eff}}) \ge \frac{1}{N_1}\sum_{i=1}^n \overset{\neg}{W}_i\,r_i(1{-}r_i)\,v_0(X_i) + o_p(1),
\]
with equality when all $\delta_i \in \{0,1\}$. \qedhere
\end{proof}

The following proposition shows that balance-preserving rounding achieves equality in~\eqref{eq:integer-variance}, and provides a consistent variance estimator.

\begin{propositionprime}{\ref{prop:matching-achievable}}
Under the conditions of Proposition~\ref{prop:matching-cost}$'$, let $\model$ be a model with bounded uniform entropy integral and suppose:
\begin{enumerate}
\item[(i)] Real-valued weights $\hat\theta$ achieve $o_p(n^{-1/2})$ imbalance uniformly over $\model$.
\item[(ii)] Balance-preserving rounding~\eqref{eq:round} preserves this rate: $\roundingerror_{\model,\hat\theta}(\round{\hat\theta}) = o_p(n^{-1/2})$.
\item[(iii)] $\|\hat\theta v_0\|_\model = O_p(n^{1/2})$ and $\|(1{-}2\hat r)v_0\|_\model = O_p(n^{1/2})$.
\item[(iv)] $\hat\theta \to_p \theta$ in $L^2(P)$.
\end{enumerate}
Then the matching estimator $\hat\psi^{\textup{match}} = N_1^{-1}\sum_i W_i Y_i - \overset{\neg}{W}_i \round{\hat\theta}_i Y_i$ is regular and asymptotically linear:
\[
\hat\psi^{\textup{match}} - \tilde\psi = \frac{1}{N_1}\sum_{i=1}^n\bigl(W_i + \overset{\neg}{W}_i\round{\hat\theta}_i\bigr)\epsilon_i + o_p(N_1^{-1/2}),
\]
with variance achieving equality in~\eqref{eq:integer-variance}:
\[
\Var(\hat\psi^{\textup{match}}) - V^{\textup{eff}} = \frac{1}{N_1^2}\sum_{i=1}^n \overset{\neg}{W}_i\, r(X_i)\{1-r(X_i)\}\,v_0(X_i) + o_p(N_1^{-1}).
\]
A consistent variance estimator is
\begin{equation}
\label{eq:integer-variance-estimator}
\hat V = \frac{1}{N_1^2}\sum_{i=1}^n \bigl(W_i + \overset{\neg}{W}_i\round{\hat\theta}_i^2\bigr)\{Y_i - \hat m(W_i,X_i)\}^2,
\end{equation}
with $\hat V / \Var(\hat\psi^{\textup{match}}) \to_p 1$ and $(\hat\psi^{\textup{match}} - \tilde\psi)/\sqrt{\hat V} \rightsquigarrow N(0,1)$.
\end{propositionprime}

\begin{proof}[Proof of Proposition~\ref{prop:matching-achievable}$'$]
We verify each conclusion in turn.

\paragraph{Regularity.}
Condition~(ii) gives $o_p(n^{-1/2})$ imbalance uniformly over $\model$. Because $\model$ has bounded uniform entropy integral, $\model$ is Donsker, so $o_p(n^{-1/2})$ imbalance uniformly over $\model$ implies $o_p(1)$ imbalance in every $f \in L^2(P)$---the definition of regularity (Definition~\ref{def:regular-matching}).

\paragraph{Asymptotic linearity.}
By condition~(ii) and the triangle inequality~\eqref{eq:att-imbalance}, $\imbalance_{h,\model}(\round{\hat\theta}) \le \imbalance_{h,\model}(\hat\theta) + \roundingerror_{\model,\hat\theta}(\round{\hat\theta}) = o_p(n^{-1/2})$. The asymptotic linearity~\eqref{eq:asymptotic-linearity-rounded} then follows from Corollary~\ref{cor:rounding}.

\paragraph{Variance formula.}
We show $V^{\hat\theta} = V^{\textup{eff}} + o_p(N_1^{-1})$ for the unrounded weights, then add the rounding cost from~\eqref{eq:rounding-variance-excess}.

For the unrounded weights, the variance decomposition~\eqref{eq:variance-decomp-1} gives
\[
N_1^2(V^{\hat\theta} - V^{\textup{eff}}) = \sum_i \overset{\neg}{W}_i(\hat\theta_i - \theta_i)^2 v_0(X_i) + 2\sum_i \overset{\neg}{W}_i(\hat\theta_i - \theta_i)\theta_i v_0(X_i).
\]
The cross term vanishes by the same argument as~\eqref{eq:cross-term-1}: regularity of $\hat\theta$ (which holds by condition~(i) and Donskerity of $\model$) gives $o_p(1)$ imbalance in $\theta v_0 \in L^2(P)$. The quadratic term is
\[
\frac{1}{N_1}\sum_i \overset{\neg}{W}_i(\hat\theta_i - \theta_i)^2 v_0(X_i) \le \|\hat\theta - \theta\|_{L^2(P_0)}^2 \cdot \|v_0\|_\infty = o_p(1)
\]
by condition~(iv).

Second, from~\eqref{eq:rounding-variance-excess}:
\[
N_1(V^{\round{\hat\theta}} - V^{\hat\theta}) = \frac{1}{N_1}\sum_i \overset{\neg}{W}_i\hat r_i(1{-}\hat r_i)\,v_0(X_i) + R\{(1{-}2\hat r)v_0\} + 2R(\hat\theta v_0).
\]
By Proposition~\ref{prop:rounding-efficiency}, condition~(iii) ensures $R(\hat\theta v_0) = o_p(1)$ and $R\{(1{-}2\hat r)v_0\} = o_p(1)$.

Combining and using $\hat r_i \to_p r_i$ from condition~(iv):
\[
N_1(V^{\textup{match}} - V^{\textup{eff}}) = \frac{1}{N_1}\sum_i \overset{\neg}{W}_i\,r_i(1{-}r_i)\,v_0(X_i) + o_p(1).
\]

\paragraph{Variance estimation.}
The estimator~\eqref{eq:integer-variance-estimator} uses squared residuals $\{Y_i - \hat m(W_i,X_i)\}^2$ in place of $v_{W_i}(X_i)$. By uniform consistency of $\hat m$, $(Y_i - \hat m_i)^2 \to_p v_{W_i}(X_i)$, so
\[
\frac{\hat V}{V^{\textup{match}}} = \frac{\sum_i (W_i + \overset{\neg}{W}_i\round{\hat\theta}_i^2)(Y_i - \hat m_i)^2}{\sum_i (W_i + \overset{\neg}{W}_i\round{\hat\theta}_i^2) v_{W_i}(X_i)} \to_p 1
\]
by the law of large numbers. The CLT follows from asymptotic linearity~\eqref{eq:asymptotic-linearity-rounded} and this variance consistency.
\end{proof}

For RKHS models, the conditions of Proposition~\ref{prop:matching-achievable}$'$ are satisfied by the minimax linear weights of \citet{hirshberg2019minimax}.

\begin{corollaryprime}{\ref{cor:integer-sbw}}
Let $\calH$ be an RKHS with bounded continuous kernel that is dense in $L^2(P)$. Assume $\E[1/e(X)] < \infty$ and $\hat\theta v_0, (1{-}2\hat r)v_0 \in L^2(P)$.

For an appropriately chosen sequence of budgets $B_n \ll n^{-1/2}$, the integer program~\eqref{eq:integer-sbw} has a solution $\round{\hat\theta}^*$ with probability tending to one. The matching estimator $\hat\psi^{\textup{match}} = N_1^{-1}\sum_i W_i Y_i - \overset{\neg}{W}_i \round{\hat\theta}^*_i Y_i$ satisfies
\[
\hat\psi^{\textup{match}} - \tilde\psi = \frac{1}{N_1}\sum_{i=1}^n\bigl(W_i + \overset{\neg}{W}_i\round{\hat\theta}^*_i\bigr)\epsilon_i + o_p(N_1^{-1/2})
\]
with variance achieving equality in~\eqref{eq:integer-variance}. The variance estimator~\eqref{eq:integer-variance-estimator} satisfies $\hat V / \Var(\hat\psi^{\textup{match}}) \to_p 1$ and $(\hat\psi^{\textup{match}} - \tilde\psi)/\sqrt{\hat V} \rightsquigarrow N(0,1)$.
\end{corollaryprime}

\begin{proof}[Proof of Corollary~\ref{cor:integer-sbw}$'$]
We verify the conditions of Proposition~\ref{prop:matching-achievable}$'$.

\paragraph{Condition~(i).}
By \citet[Theorem 2]{hirshberg2019minimax}, the minimax linear weights achieve $o_p(n^{-1/2})$ imbalance uniformly over the unit ball of $\calH$.

\paragraph{Condition~(ii).}
By Proposition~\ref{prop:rounding}, balance-preserving rounding preserves the $o_p(n^{-1/2})$ rate.

\paragraph{Condition~(iii).}
By Corollary~\ref{cor:rounding-efficiency-rkhs}, any $f \in L^2(P)$ can be approximated by $f_n \in \calH$ with $\|f_n\|_\calH = O(n^{1/2})$. Since $\hat\theta v_0, (1{-}2\hat r)v_0 \in L^2(P)$ by assumption, both have $\calH$-norm $O_p(n^{1/2})$.

\paragraph{Condition~(iv).}
$L^2$ consistency of $\hat\theta$ toward $\theta$ follows from \citet[Theorem 1]{hirshberg2019minimax}: for bounded continuous RKHS kernels dense in $L^2(P)$, the minimax linear weights converge to the IPW weights in $L^2(P)$ at a rate depending on the kernel's eigenvalue decay.

\paragraph{Feasibility.}
The rounded minimax linear weights are feasible for~\eqref{eq:integer-sbw} with budget $B_n$ equal to the sum of the imbalance and rounding rates (see the explicit rate calculations earlier in this section). The integer-optimal solution $\round{\hat\theta}^*$ minimizes $\sum_i \overset{\neg}{W}_i \theta_i^2$ over feasible integer weights, so has at least as good imbalance as the rounded continuous solution. In particular, $\round{\hat\theta}^*$ satisfies conditions~(i)--(iv) of Proposition~\ref{prop:matching-achievable}$'$.

Applying Proposition~\ref{prop:matching-achievable}$'$: asymptotic linearity follows from~\eqref{eq:asymptotic-linearity-rounded}, the variance achieves equality in~\eqref{eq:integer-variance}, and the variance estimator~\eqref{eq:integer-variance-estimator} is consistent.
\end{proof}


\section{Regression for Orthogonality}
\label{app:regression-orthogonality}

\citet{wang2020debiased} observe that the balancing logic can be reversed. Rather than choosing weights $\hat\gamma$ so that the error $\delta\gamma = \hat\gamma - \gamma^{\text{ipw}}$ is approximately orthogonal to a model $\model$ for the outcome, we can choose the outcome regression $\hat \mu$ so that its error $\delta \mu = \hat \mu - \mu$ is approximately orthogonal to the weight error. Here we sketch the idea in our notation.

Recall that the bias of the AIPW estimator can be written as
\begin{equation}
n^{-1} \sum_{i=1}^n W_i \delta\gamma(X_i) \delta m(X_i) + n^{-1} \sum_{i=1}^n \{W_i \gamma^{\text{ipw}}(X_i) - 1\} \delta m(X_i).
\end{equation}
When the weights are the true inverse propensity weights, the second term is an average of mean-zero terms and is $O_p(n^{-1/2})$. The first term---the inner product of the errors---is what we need to control.

The balancing approach bounds this by controlling $\sup_{m \in \model} n^{-1} \sum_i W_i \delta\gamma(X_i) m(X_i)$, ensuring approximate orthogonality to a model $\model$ containing $\delta \mu$. But we could instead control the other factor: choose $\hat \mu$ so that
\begin{equation}
\sup_{\gamma \in \Gamma} n^{-1} \sum_{i=1}^n \gamma(X_i) \{\hat \mu_1(X_i) - Y_i\} = O_p(n^{-1/2})
\end{equation}
for a class $\Gamma$ containing $\delta\gamma$. This makes the regression error approximately orthogonal to the types of errors we expect in estimating the propensity score.

For the treatment-specific mean, this can be implemented using sample splitting. Let $(X'_i, Y'_i)$ be an independent copy of $(X_i, Y_i)$ among treated units. Choose $\hat \mu_1(\cdot)$ to satisfy
\begin{equation}
\sup_{\gamma \in \Gamma} \left|n^{-1} \sum_{i=1}^n \gamma(X_i) \hat \mu_1(X_i) - n^{-1} \sum_{i=1}^n \gamma(X'_i) Y'_i\right| = O_p(n^{-1/2}).
\end{equation}
Because $\hat \mu_1(\cdot)$ is independent of $W_1 \ldots W_n$, the equicontinuity term in the error decomposition is $O_p(n^{-1/2})$ automatically. This yields $\sqrt{n}$-consistent estimation when the Riesz representer $\gamma^{\text{ipw}}$ is regular in the sense that $\|\hat\gamma - \gamma^{\text{ipw}}\|_\Gamma = O_p(1)$ for a Donsker class $\Gamma$---with essentially no assumptions on the outcome model.

This is the reverse of the balancing approach: there we assume the outcome model is regular and make no assumptions on the propensity score; here we assume the propensity score model is regular and make no assumptions on the outcome.

\ifarxiv
\printbibliography[title=Appendix References]
\end{refsection}
\else
\ifsupp
\bibliography{review}
\fi
\fi
\fi
\end{document}
